% concatenated from 231014828.tex
%%%%% Setup %%%%%%%%%%%%%%%%%%%%%%%%%%%%%%%%%%%%%%

%\documentclass{amsart}
\documentclass[11pt,draft,reqno]{amsart}
  \usepackage{amsmath}
\usepackage{amssymb}
\usepackage{amsfonts}
\usepackage{eucal}
\usepackage{xcolor}
% \usepackage[monochrome]{xcolor}
    \makeatletter
     \def\section{\@startsection{section}{1}%
     \z@{.7\linespacing\@plus\linespacing}{.5\linespacing}%
     {\bfseries%\normalfont\scshape
     \centering
     }}
     \def\@secnumfont{\bfseries}
     \makeatother
 \setlength{\textheight}{20.5 cm}
\setlength{\textwidth}{13.5 cm}
\usepackage[full]{textcomp}
 

%%%%%%%%%%%%%%%%%%%%%%%%%%%%%%%%%%%%%%%%%%%

\newtheorem{theorem}{Theorem}[section]
\newtheorem{lemma}[theorem]{Lemma}
\newtheorem{proposition}[theorem]{Proposition}
\newtheorem{corollary}[theorem]{Corollary}
\theoremstyle{definition}
\newtheorem{definition}[theorem]{Definition}
\newtheorem{example}[theorem]{Example}
\newtheorem{question}[theorem]{Question}
 \theoremstyle{remarks}
\newtheorem{remarks}[theorem]{{\bf  Remarks}}
\newtheorem{remark}[theorem]{{\bf  Remark}}
\numberwithin{equation}{section}
\setcounter{page}{1}
 \usepackage{a4wide}
\font\sevenrm =cmr10 at  7  pt
\def\ddate {\sevenrm \ifcase\month\or January\or
February\or March\or April\or May\or June\or July\or
August\or September\or October\or November\or December\fi\! {\the\day}, \!{\sevenrm\the\year}}


%%%%% Macros List  %%%%%%%%%%%%%%%%%%%%%%

\def \a{{\alpha}}
\def \b{{\beta}}
\def \D{{\Delta}}
\def \d{{\delta}}
\def \e{{\varepsilon}}
\def \g{{\gamma}}
\def \G{{\Gamma}}
\def \k{{\kappa}}
\def \l{{\lambda}}
\def \mh{{\hat {\mu} }}
\def \n{{\eta}}
\def \o{{\omega}}
\def \O{{\Omega}}
\def \p{{\varphi}}
\def \t{{\vartheta}}
\def \ttau{{\theta}}
\def \m{{\mu}}
\def \s{{\sigma}}
\def \Y{{\Theta}}
\def \noi{{\noindent}}
\def \qq{{\qquad}}
\def \dd{{\rm d}}
\def\F{{\mathcal F}}
\def\A{{\mathcal A}}
\def\P{{\mathcal P}}
\def\={\,=\,}

%%%%%%%%%%%%%% Bbb characters %%%%%%%%%%%%%%

\def\E{{\mathbb E \,}}

\def \T{{\mathbb T}}

\def\P{{\mathbb P}}

\def\R{{\mathbb R}}
\def \V{{\mathbb V}}
\def \B{{\mathbb B}}

\def\Z{{\mathbb Z}}

\def\Q{{\mathbb Q}}

\def\N{{\mathbb N}}

\def\NS{{\mathbb N}^*}

\def\C{{\mathbb C}}
 \def\beq{\begin{equation}}
\def\eeq{\end{equation}}
  
\def\ben{\begin{eqnarray}}
\def\een{\end{eqnarray}}\font\sevenrm= cmr10 at 9 pt
\def\ddate {\sevenrm \ifcase\month\or January\or
February\or March\or April\or May\or June\or July\or
August\or September\or October\or November\or December\fi\! {\the\day}, \!{\sevenrm\the\year}}

%%%%%%%%%%%%%%%%%%%%%%%%%%%%
  \hfuzz =0pt
  \scrollmode%%%%%%%%%%%%%%%%%%%%%%%%%%%%
 
\font\phh=cmcsc10
at  8 pt
%\scrollmode
\font\gse= cmb10 at 12 pt
 \font\rr= cmr10 at 9,7pt
\hfuzz =10pt 
   \font\phh=cmr10 at  8,2 pt
\font\ph=cmcsc10 at 10 pt  \scrollmode
\font\gsec= cmb10 at 9,9  pt
\font\gssec= cmb10 at 8,2  pt
 \font\tf=cmbxsl10 at 9,5pt
\font\mog=cmbx10 at 24 pt
\font\trog=cmssqi8 at 10pt
\font\tog=cmbxsl10 at 9pt
\font\tog=cmbxsl10 at 124pt
\font\mmog=cmbx10 at 14 pt
\font\rog=cmssqi8 at 9 pt
 \font\eightrm= cmr10 at 8 pt
 \font\fontr=cmssqi8 at 10 pt
 \font\tf=cmbxsl10 at 9pt
 \font\sevenrm= cmr10 at 7,3 pt
 \font\eightrm= cmr10 at 8 pt
 \font\sevenit=cmti7



%\textwidth16cm

%\usepackage[normalem]{ulem}
  

\begin{document}

%\thispagestyle{empty}

\vskip 1.5cm

 \scrollmode




  \newcommand{\sm}{\vskip} \newcommand{\ens}{\enspace} \


 \title[Divisibility and primality in   random walks]
  {Divisibility and primality in random walks}
 
   \author{ Michel J.\,G.\, WEBER}
  \address{IRMA, UMR 7501, Universit\'e
Louis-Pasteur et C.N.R.S.,   7  rue Ren\'e Descartes, 67084
Strasbourg Cedex, France.
 \vskip 1 pt   E-mail:    {\tt  michel.weber@math.unistra.fr}}
\keywords{Divisibility,  primality, random walk, binomial random walk, Bernoulli random walk, Theta    function, value distribution, quasi-prime numbers, prime numbers, Rademacher random walk,   i.i.d. sums, Cram\'er model,    Riemann Hypothesis,   Farey series,  prime number theorem,  correlation  properties.\\ %\indent\sevenrm{DBRW.tex}
2010 \emph{Mathematics Subject Classification}: {Primary 11K65,    11N37, 60G50 ; Secondary 11A25, 60E05, 60F15.}}
 \begin{abstract}    In this paper we study the divisibility and primality properties of the Bernoulli random walk. We improve or extend some of our divisibility results  to wide classes of iid or independent non iid random walks. We also obtain new primality results       for  the Rademacher random walk. We study the value distribution of divisors of the random walk in the Cram\'er model, and obtain a general estimate of a similar kind to that of the Bernouilli model. 
 %to study similar divisibility or primality  properties in  the Rademacher model and the  . 
   Earlier   results on  divisors and quasi-prime numbers in the Bernoulli model are recorded, as well as  some other recent    in the Cram\'er random model,  based on  an estimate due to Selberg.
\end{abstract}
\maketitle

\tableofcontents
\section{Introduction.}\label{s1}

%In this paper we improve or extend some of our results concerning the Bernoulli model, to iid or independent non iid random walks. We also obtain new primality results        in  the Rademacher model. We study the value distribution of divisors in the Cram\'er model, and obtain a general estimate of a similar kind to that of the Bernouilli model. 
In this paper we study      divisibility and primality properties of the most  fundamental     integer lattice random walks: the      Bernoulli random walk and Rademacher random walk. This is continuing previous works in  \cite{W5}, \cite{W5a}, \cite{W1}, \cite{W4},     \cite{W7},  
  % \cite{W5b},
   \cite{W6}, where these questions
     have been thoroughly investigated, as well as other questions related to their  asymptotic structural independence. %   value distribution of the divisors of $B_n$ and their asymptotic independence  properties,  
 We next extend the study to  the Cram\'er random walk on which Cram\'er's     model built. This  is an important example of independent, non identically distributed random walk.  A     thorough   probabilistic study 
 %study of  the \lq prime   numbers\rq \,   in 
 of this   model was recently made in \cite{W8}.
 %, where we prove  a     sharp density 1 result,  using an   estimate of Selberg. 
     The case of a general square integrable lattice random walk was very recently  investigated. We prove in  \cite{W10},  \cite{W9} sharp divisibility results    valid for quite large classes of independent,   identically distributed  (iid) lattice   random walks, and obtain also divisibility results in independent   non identically distributed   case. The used tools are new in this context, and built on  a remarkable coupling method and a recent improvement of an important necessary condition of arithmetical type for the applicability of the local limit theorem. We   refer to Part IV of our 2009's book   \cite{W2}, the  Chapters 11    and  13, which  are  devoted to divisors, and their applications in Analysis, especially Fourier analysis. 

\vskip 2   pt  Throughout  let     $\rho\in ]0,1[$ and $\mathcal B(\rho)=\{B_n(\rho), n\ge 1\}$ denotes  the binomial random walk,   $B_n(\rho)=\sum_{k=1}^n \b_k(\rho)$, for each $n$, where  $\b(\rho)=
\{\b_i(\rho) , i\ge 1\}$ is a sequence of iid  binomial random variables ($\P\{\b_i(\rho)=0\} =1-\P\{\b_i(\rho)= 1\}=\rho $).  When $\rho =1/2$ we use the simplified notation $\mathcal B =\{B_n , n\ge 1\}$, $B_n =\sum_{k=1}^n \b_k $, for each $n$,    $\b =
\{\b_i  , i\ge 1\}$ 
  being iid Bernoulli random variables.
  %   ($\P\{ \b_i =  0\}=\P\{ \b_i =  1\} =1/2$).  
    Let also $\{R_n, n\ge 1\}$, $R_n=\sum_{k=1}^n \e_k$,  where  $\e=
\{\e_i , i\ge 1\}$ is  a sequence of iid 
Rademacher random variables ($\P\{ \e_i =\pm 1\} =1/2$). Further let   $\mathcal S= \{S_n, n\ge 3\}$, $S_n=\sum_{i=3}^n \xi_i$, be the random walk associated to the Cram\'er model,   $\xi_i$,  $i\ge 3$, being independent
  random variables defined by $\P\{\xi_i= 1\}=  1-\P\{\xi_i= 0\}= \frac1{\log i}$. The basic probability space on which these random variables are defined is noted $(\O,\A,\P)$, and the  corresponding expectation symbol $\E$.
  \smallskip\par
  %Part IV of our 2009's book   \cite{W2}, the  Chapters 11    and  13,   are  devoted to divisors, and their applications in Analysis, especially Fourier analysis.    We  record  earlier     results   on  divisors and quasi-prime numbers in the Bernoulli model,  and  recent   sharp density one results on  the \lq primes\rq\,  in the Cram\'er model,  using an   estimate of Selberg for the lower bound.
  
\smallskip\par 
%The class of   binomial  random walks  in which $\rho$ is varying with $i$, is  of first importance in random models in Number theory. See Freiman-Pitman \cite{FP} for another example arising the study of the number of partitions of a number into distinct parts. Its study is not considered here.

\smallskip\par
 
 %to study similar divisibility or primality  properties in  the Rademacher model and the  . 
  

\medskip\par
 {\it Notation and convention.}    
-- Throughout the paper,  let  
$p$   denotes an arbitrary prime number,   the letter $C$   a universal constant whose value may change at each occurence, and   $C_{\a, \b, \ldots}$     a constant depending only on the parameters ${\a, \b, \ldots}$. 
\vskip 2 pt 
\noi -- We agree that $\sum_{\emptyset}=0$, $\sup_{\emptyset}=0$. 










\bigskip\par
\section{Preamble on summation methods.}\label{s2}




In background  of  the study of the divisibility/primality questions   related to these models are often are summation methods, which may be  of   particular efficiency, as they are linked between them through   Tauberian results. We first begin with some important facts. The uniform  model  (the uniform distribution on $\{0,1,2\ldots, n\}$, for each $n$), also called the   Kubilius model,  is related to Ces\`aro's summation method, whereas the binomial model is  to Euler
  summation method, and there exist Tauberian theorems 
  comparing   summation  methods, in particular these two methods.  Therefore   fine arithmetical results can be   directly used, for one summation  method (one random model), depending on the rate of convergence in this one, to answer a question concerning another random model.  
\begin{definition} A  set  
$A\subset 
\N$   has  Ces\`aro density $\l$ of order 1 ($C_1$--density $\l$),  if 
$$\lim_{n\to \infty}\frac{\#\{ j\le n: j\in A\}}{n}= \l.
$$
\end{definition}Often, the most readily available information about a set  of integers is that it has $C_1$-Ces\`aro density with a given rate of convergence. \begin{definition}  A subset $A$ of $  \N$ is said to have {\it Euler density  $\l$ 
with parameter
$\varrho $} (in short ${ E}_\varrho$ density $\l$) if 
\begin{equation*} 
\lim_{n\to \infty}\sum_{j\in A} {{n}\choose{j}} \varrho^j (1-\varrho)^{n-j}= \l.
\end{equation*} 
\end{definition}
Diaconis and Stein proved in \cite{DS}
(Theorem 1) the following fundamental result:
\begin{theorem}\label{ds}For any $A\subset \N$, and $\varrho\in ]0,1[$ the following assertions are equivalent:
 \begin{eqnarray*} \label{dschar} ({\rm i})  & &\qq \hbox{\it $A$ has  ${E}_\varrho$ density $\l$},
  \cr   ({\rm ii})  & &\qq\lim_{t\to \infty} e^{-t}\sum_{j\in A} \frac{t^j}{j!}=\l,
 \cr ({\rm iii})  & &\qq \hbox{\it  for all $\e>0$}, \quad    \lim_{n\to \infty} \frac{\#\{j\in A : n\le j< n+\e\sqrt n\} }{\e\sqrt n} 
 =\l . 
 \end{eqnarray*}
\end{theorem} 
 
%This  answers a question raised independently by Erd\"os and Gleason.
\begin{remark} \rm The implication (ii)$\Rightarrow$(iii) was in fact
 first established by Jurkat \cite{J}, and a
proof   using an extension of Wiener Tauberian theorem was given by
Moh in \cite{Mo}, not quoted in the paper. 
\end{remark}\vskip 2pt 
An important consequence presented by the authors as the principal  result of the paper,  is   that:
\begin{equation} \label{E.rho} \hbox{\it  The existence and value of ${  E}_\varrho$ density does not depend on
the value of $\varrho$}.  
\end{equation}  
%This is presented by the authors as the principal  result of the paper. 
A study of limit results in the present setting, for binomial random walks,  therefore reduces to the case of the Bernoulli random walk. 
As a consequence of  Theorem 3  in Diaconis and Stein  \cite{DS}, we have
\begin{theorem} If $A\subset \N$ has $E_\varrho$ density $\l$, then $A$ has $C_1$ density $\l$. Conversely the following implication  in which   $g$ is assumed to be regularly varying at infinity with exponent $\a$, $-1<\a\le -1/2$, holds true 
\begin{equation} \label{ 2}   
   \Big|\frac{\#\{ j\le n: j\in A\}}{n}- \l\Big|\le g(n)\qquad\!\!\!\!\!\Longrightarrow  \qquad\!\!\!\!\!  \Big|\P\{B_n(\varrho) \in A\} - \l\Big|\le K g(n)n^{1/2}. 
 \end{equation}
 \end{theorem}
See also Hardy \cite{H1},    Theorem 149.  
 Let
$\{\a_k, k\ge 0\}$ be a  sequence of reals, $\l$ real and $\e(n)$ be positive reals. More generally, it follows from  Lemma \ref{Eulerbound} that 
 \begin{eqnarray*} \label{4.61c}\Big|\frac{1}{n}\sum_{j=0}^{n-1}\a_j- \l \Big|\le \frac{\e(n)}{n^{ 1/2} }
%\qq \qq \Big| \sum_{j=0}^{n-1}\a_j- n\l \Big|\le  \e(n) n^{ 1/2}  
\quad \Longrightarrow \quad  \Big| \sum_{h=0}^n   2^{-n}  
{{n}\choose{h}}\a_h-\l \Big|  \le   \frac{C}{ \sqrt n}\, \max_{1\le \ell\le n} (\e(\ell) \sqrt \ell)     .\end{eqnarray*}
   
 
Let us illustrate this by giving an application to {\it $k$-free} numbers in $\mathcal B(\rho)$, which is implicitly included in \cite{DS}, (Corollary 4).
%\subsection{A sharp estimate of $\P\big\{ B_n(\varrho)\ k\hbox{-{\it free}}\big\}$.}\label{s.1.3}
   
 \begin{definition}An integer $s$ is said {\it $k$-free} if and only if it is not divisible be the $k$-th power of any prime. 
 \end{definition}Leahey and Nymann proved in \cite{LN}  among other papers, by using a formula analogous to 
    \begin{equation*}
    %\label{Bnprime}
    1-\P\{ B_n(\varrho)\ \hbox{is prime}\}= - \sum_{2\le d\le n }  \m(d) \P\{ d|B_n(\varrho)\},  
\end{equation*}  $\m$ being M\"obius function, that
\begin{equation}\label{Bnprime}   
\lim_{n\to \infty}\P\big\{ B_n(\varrho)\ k\hbox{-{\rm free}}\big\}= \frac{1} {\zeta(k)},\end{equation}    where $\zeta$ is the Riemann zeta function.
       A much more precise result holds:
\begin{corollary}  {\it For $k\ge 2$ integer,}
 \begin{equation}\label{k.free.DS} \Big|\P\big\{ B_n(\varrho)\ k\hbox{-{\rm free}}\big\}  - \frac{1}{\zeta(k)} \Big|\le E(n) , 
\end{equation}
where
$$ E(n)\le C\left( n^{\frac1k-1}\cdot  e^{ -Ak^{-8/5}\frac{(\log 
n^{{3/5}}}{(\log\log  n) ^{1/5}}}\right).$$
%Note that $ \frac{1}{\zeta(2)}=\frac{6}{\pi^2}$.
\end{corollary} 
\begin{proof} Let ${\mathcal Q}_k$,  ${\mathcal Q}_k(x)$ respectively denote the set of $k$-free integers, and the number of $k$-free integers $\le x$ ($k\ge 2$, integer).  By a result of Walfisz \cite{Wal},
  \begin{equation}\label{k.free} 
{\mathcal Q}_k (x)= 
\frac{x}{\zeta(k)} + {\mathcal  O} 
\big( x^{\frac1k}. e^{ -Ak^{-8/5}.\frac{(\log 
x^{{3/5}}}{(\log\log  x) ^{1/5}}}\big)
 . 
\end{equation}
Take $A= {\mathcal Q}_k$  and note that 
$$\Big|\frac{\#\{ j\le n: j\in A\}}{n}- \frac{1}{\zeta(k)} \Big|\le g(n)$$
with 
$$g(n)= n^{\frac1k-1}\cdot e^{ -Ak^{-8/5}.\frac{(\log 
n^{{3/5}}}{(\log\log  n) ^{1/5}}},$$
Clearly $g(n)$   is regularly varying at infinity with exponent $\a$, $-1<\a\le -1/2$. Thus \eqref{k.free.DS} follows.
\end{proof}
 \medskip\par
%\subsection{$0$-density sequences} 
  For $0$-density sequences, we have the following complementary result, which does not appear  in the literature.
% It follows from    Hardy \cite[ineq.\,3]{H}  that   the   result holds. 
\begin{proposition}
%\vskip 5 pt {\it
Let $A$ be a set of integers of (natural) density $0$. Then for every $p>1$,
 %}
\begin{equation}\label{hardy1} \lim_{N\to \infty}\frac{1}{N}\sum_{n=1}^N \P\{ B_n(\varrho)\in A\}^p =0.\end{equation} 
 \end{proposition} 
 This follows from Hardy \cite[ineq.\,3]{H}: {\it Let $p>1$. Then for $0<\varrho <1$, $a_n\ge 0$,}
 \begin{equation}\label{hardy} \sum_{n=1}^\infty \left(\sum_{m=0}^n   {{n}\choose{m}}\varrho^m(1-\varrho)^{n-m}a_m\right)^p\le \varrho^{-1}  \sum_{n=1}^\infty a_n^p.
 \end{equation} 

% \begin{proof} We use  Hardy \cite[ineq.\,3]{H}: {\it Let $p>1$. Then for $0<\varrho <1$, $a_n\ge 0$,}
%\begin{equation}\label{hardy} \sum_{n=1}^\infty \Big(\sum_{m=0}^n C_n^m \varrho^m(1-\varrho)^{n-m}a_m\Big)^p\le \varrho^{-1}  \sum_{n=1}^\infty a_n^p.
%\end{equation} 
%By taking $a_m= \chi_{A\cap [1,N]}(m)$ and next letting $N\to \infty$, the claimed result follows.
%\end{proof}. 

  \bigskip\par
 We refer to Apostol \cite{A}, McCarthy \cite{MC}, Montgomery and Vaughan \cite{MV}, Tenenbaum \cite{Te1}   as  bibliographical sources on arithmetical functions, and for results on arithmetical properties of $r$-uples of integers, co-primality,  pairwise co-primality,   least common multiple, and associated arithmetical functions.   We cite  T\'oth's elaborated survey \cite{To} for   sharp limit results of this kind with  estimates of the remainder,  Fernandez and Fernandez \cite{FF}, % and references therein
 for  limiting moments and distribution,   for instance of  {gcd}$ (X^n_1, \ldots, X^n_n)$, {lcm}$(X^n_1, \ldots, X^n_n)$ of  arrays of independent uniformly distributed random variables. These specialized arithmetical functions %, although of interest, 
% are all -{\it in} probability type, and 
are not considered here.  In each 
case, however, what is investigated is often the behavior for $n$ large of $\E f(T_n)^r$, $r\ge 1$, where $f(n)$ is one of these  arithmetical functions, $f$ may be varying with $n$ (for arrays of random variables), $T_n= Y_1+\ldots+Y_{\nu(n)}$  (the expectation symbol $\E$ corresponding to the underlying probability space on which these random variables $Y_j$ are defined). 
%\vskip 2 pt   
Recall that any arithmetical function $f(n)$ can be represented in the form
$$ f(n) =\sum_{d|n} \Phi(d).$$
The function $\Phi(d)$ is defined by the M\"obius inversion formula
$$ \Phi(n)= \sum_{d|n}\mu\Big(\frac{n}{d}\Big) f(d)  .$$
For instance for the Bernoulli model,
$$\E f(n) =\sum_{d|n} \P\{d|B_n\}\,\Phi(d).$$
See also sub-section \ref{sub.bin.way.back}. Hence the interest of having at disposal sharp estimates of  $\P\{d|B_n\}$ when both $d$ and $n$ vary.
 \vskip 2 pt
When $\Phi(d)$ is such that the series $\sum \Phi(d)/d$ is absolutely convergent, we recall that by Wintner's Theorem
$$ \lim_{N\to\infty}\frac{1}{N}\sum_{n=1}^Nf(n)= \sum_{d=1}^\infty \frac{\Phi(d)}{d}.$$
See  Postnikov \cite{Po} p.\,138. 
For many arithmetical functions $\Phi(d)$ is simple. Here are some examples of mostly studied functions. Let  $\mathcal P$ be the set of prime numbers and let the letter $p$ always denote a prime number.  
       %Take$A$ to be the set of multiples of an integer$d$. Then (\cite{DS},\,Example 2),
% \begin{equation}\label{multiple} \Big|\P\big\{d|B_n(\varrho) \}- \frac{1}{d}\Big|\le 
 %e^{-        {
% 8n \varrho(1-\varrho)  /  d^2}  }.
%\end{equation} \bigskip\par
%can be either of the examples below 
  \begin{eqnarray*}
    \s_s(n) = \sum_{d|n} d^s &&\qq\hbox{\rm the sum of $s$-powers of divisors of $n$, $s\in \R$;}
  \cr   J_s(n)=\sum_{d|n} d^s \m\Big(\frac{n}{d} \Big) &&\qq\hbox{\rm the generalized Euler totient function, $s>0$;}
\cr    w(n)=\sum_{p|n}\frac{\log p}{p} &&\qq\hbox{\rm Davenport's function;}
\cr \Phi(n)=\sum_{d|n} \frac{\log d}{d} &&\qq\hbox{\rm Erd\"os and  Zaremba function;}
\cr 
 \omega(n) = \sum_{p|n,p\in \P} 1 &&\qq\hbox{\rm the prime divisor function;}
   \cr  \Omega(n) = \sum_{p|n,p\in \P} p &&\qq\hbox{\rm   the sum of  prime divisor function;}
 \cr\theta(n)= 2^{\o(n)} &&\qq\hbox{\rm    counting the number of prime divisors of  $n$.}\end{eqnarray*}   
   The classical Euler totient function $\p(n)$ counting the number of integers $k $ less than $n$ and such that $(k,n) =1$ is $J_1(n)$, and $\s_0(n)$ is the divisor function $d(n)$ counting the number of divisors of $n$. 
% For instance, we have 
%by Gronwall's estimates,
% \begin{equation} \label{gron}   \limsup_{n\to \infty} \frac{\log
%\big(\frac{\s_s(n)}{n^s}\big)}{\frac{(\log n)^{1-s}}{\log\log n}} =\frac{1}{1-s}, \qq (0<s<1)\qq  \limsup_{n%\to \infty}\frac{\s_{-1} (n)}{n\log\log n}= e^\g,
%\end{equation}
 %    where $\g$ is Euler's constant.
 

%\vskip 3 pt 
%The case of the divisor function,  corresponding to the first example with $s=0$, does not seem to be
% reduced to this approach.
\vskip 2 pt To the best of our knowledge, {\it all}     divisibility results for a random walk existing in the literature are {\it in}-probability results, except for \cite{W6}. It is a pending question to know whether it is possible to derive almost sure   counterparts, from existing {\it in}-probability results using almost sure convergence criteria, which is a far more complicated task.     These criteria are   not  familiar to number theorists. 
%Like the Bernoulli random walk, 
A random walk  often has   natural  structural independence properties. Therefore  a study of almost sure divisibility/primality  in a random walk is possible, at the price of many additional technicalities. That's make the research in this domain  attractive. In the case of the Bernoulli random walk, it is not complicated to show   that   the system 
 $$ \big\{ \chi\{d|B_n\},    \    2\le d\le n,\ n\ge 1\big\}$$
is a mixing system, namely        the   correlation
function   satisfies for each     $d$ and $\d$ greater than 2, 
 \begin{equation}
 \label{delta.def.lim}
  \lim_{m-n\to \infty }  \P\big\{ d|B_n  ,    \d|B_m\big\}-\P \{ d|B_n \}\P \{    
\d|B_m \}= 0.
\end{equation}
The estimation of the correlation 
 \begin{equation}
 \label{delta.def}
  {\Delta} \big( (d,n), (\d, m)\big)= \P\big\{ d|B_n  ,   \d|B_m\big\}-\P \{ d|B_n \}\P \{    
\d|B_m \}.
\end{equation}
  was made in \cite{W4},   and is substancially improved in    
  \cite{W14}. We have the following formulation
  \begin{eqnarray}\label{6.4}& & {\bf \Delta} \big( (d,n),  (\d, m)\big)  =
\cr & &\qq \quad{1\over d\d}  \sum_{j=1}^{d-1}\sum_{h=1}^{\d-1}  e^{ i\pi  (  {j\over d}n +   {h\over \d}m) }  \cos^{m-n}{ \pi   {h\over
\d} }\, \Big(\cos^n{ \pi  (  {j\over d} +  
{h\over \d})   }-\cos^{n} {
\pi j\over d}     \cos^{n} {
\pi h\over \d}\Big),
 \end{eqnarray}  
which   can be more simplified    using symmetries.



\medskip  \par

The paper is organized as follows. In Sections \ref{s2}-\ref{s6} we  study the divisors in the Bernoulli random walk. We also study examples. Hitherto, the   divisibility or primality properties of general   i.i.d.   or independent non i.i.d. lattice   random walks, were not much explored and the results we present are mostly the first. In Section \ref{s8}, we describe the value distribution of divisors in the Rademacher random walk and present some applications; we essentially study primality properties. In Section \ref{s9}, we study the value distribution of divisors in the Cram\'er random walk, and record some primality   results obtained by the author. Complements to previous sections are gather in Appendix.






\section{Divisibility in the Bernoulli random walk.}
 \label{s3}
We study  the probability
  $ \P \{d|B_n \}$ when both $d$ and $n$ vary.    We first prove a basic but  handy estimate. Next we state a sharper and nearly optimal estimate using Theta function, earlier established by the author in \cite{W1}. 
  Applications are given to: 
   $\E \s_{-1}(B_n)$, $\E d(B_n)$,     $\P\big\{  B_n\, , B_m \ \hbox{co-prime} \big\}$, $\P\{ B_n\, \hbox{$z$-quasi-prime} \}$, $\P\{ B_n\,\hbox{  prime} \}$. We will  also estimate the probability $\P\{D|B_nB_m\}$ for $n<m$. The study of the divisors of $B_nB_m$  is one piece of the estimation of the correlation appearing in \eqref{delta.def} when $n$ and $m$ are relatively near,  details are given in Section \ref{s6}.

   %Derivation of  estimates of $\E f(B_n)$ where $f$ is an arithmetical function also follow. 
    %We estimate for instance $\E \Psi(B_n)$ where $\Psi(n)= \sum_{d|n} \frac{(\log d )(\log\log d)}{d}$ arises from  a recent improvment of a 
%result of   Erd\"os and  Zaremba in \cite{W11} by the author.

 

% $\lim_{n\to \infty} \P\big\{d|S_n\big\}$ exists or not, (b) given   $d$ and $n$, to estimate $\P\big\{d|S_n\big\}$. It is clear that the study of (a) may require a direct estimation of $\P\big\{d|S_n\big\}$, but for $d$ fixed and $n$ large, which is another question than the one contained in (b) where we want to find, if possible, optimal estimates of this probability, $d$ and $n$ being fixed or/even uniformly over all $d$, or for all $d$ in  a test set of divisors.
 %\medskip\par
 %Of course this  observation is general and applies to any other example of study, such as 
 %$$\P\big\{ S_n\,{\rm squarefree}\big\},\qq \P\big\{ d|S_n\,,\, \d|S_m\big\},\qq\P\big\{ S_n\, ,\,z{\rm-quasiprime}\big\}, \qq \E f(S_n),$$ 
%where $f$ is an arithmetical function, etc ... 


   % (b) given   $d$ and $n$, to estimate $\P\big\{d|S_n\big\}$. %It is clear that the study of (a) may require a direct estimation of $\P\big\{d|S_n\big\}$, but for $d$ fixed and $n$ large, which is another question than the one contained in (b) where we want to find, if possible, optimal estimates of this probability, $d$ and $n$ being fixed or/even uniformly over all $d$, or for all $d$ in  a test set of divisors.
% \medskip\par
 %Of course this  observation is general and applies to any other example of study, such as 
% $$\P\big\{ S_n\,{\rm squarefree}\big\},\qq \P\big\{ d|S_n\,,\, \d|S_m\big\},\qq\P\big\{ S_n\, ,\,z{\rm-quasiprime}\big\}, \qq \E f(S_n),$$ 
%where $f$ is an arithmetical function, etc ... 
 
 
\subsection{Basic estimates.}
% of ${ {\P} \{ d | B_n \}}$ and  ${ {\E}\, d(B_n)}$, %These are  derived from a slightly more precise formulation of   (\ref{multiple}).
 %and  provide some applications.
\begin{proposition}   \label{Lemma44}
  For any integers   
 $d\le n$,
  \begin{equation}\label{multiple1}\big|\P \{ d | B_n \} -\frac{1}{d}\big|
   \le      {2\over d}\sum_{ 1\le j<d/2}
e^{-        {
 2n j^2/  d^2}  } . 
\end{equation}
\end{proposition}
\begin{proof} From the formula 
\begin{equation}\label{basic}d\,\chi(d| B_n)=\sum_{j=0}^{d-1} e^{2i\pi  {j \over d}B_{n }},
\end{equation}
  we get by integration,
%properties of characteristic functions of sums of independent random variables,    we get  after integration
\begin{equation}\label{basic.i} \P\big\{ d | B_n\big\}={1\over d}\sum_{j=0}^{d-1} e^{i\pi n{j\over d}}\big(\cos { \pi j\over d}\big)^{n}.
  \end{equation}
 As $e^{i  n(\pi-x)} \cos^n    (\pi-x)  =(-1)^ne^{-i 
nx}  (-1)^n\cos^n   x= e^{-i 
nx}  \cos^n   x  $, we have in fact  by distinguishing the case $d$ even from the case $d$ odd,   the simplified form
\begin{equation}\label{Lemma41} \P\big\{ d | B_n\big\} ={1\over d}+ {2\over d}\sum_{ 1\le j<d/2}
(\cos{ \pi n{j\over d}})\big(\cos { \pi j\over d}\big)^{n}.
\end{equation}
%\leqno(2.3)
    Thus   
    $$\big|\P\big\{ d | B_n\big\} -\frac{1}{d}\big|
   \le      {2\over d}\sum_{ 1\le j<d/2}
e^{n\log (  1-2 \sin^2{ \pi j\over 2d})} .$$
     As $\sin x\ge (2/\pi )x$, $0\le x\le \pi/2 $, it follows that $\log (  1-2 \sin^2{ \pi j\over 2d})\le -2 \sin^2{ \pi j\over 2d}\le  -2  ({
  j\over  d}) ^2   $. Thus 
\begin{eqnarray*} \big|\P\big\{ d | B_n\big\} -\frac{1}{d}\big|
  &\le &    {2\over d}\sum_{ 1\le j<d/2}
e^{-        {
 2n j^2/  d^2}  }.
\end{eqnarray*}   
\end{proof} 
This   implies estimate in   Example 2 in Diaconis and Stein \cite{DS}    without using Taylor's expansion or Ramus' formula.  See also Remark \ref{rem.Ex.2}.
 %   Note    that Ramus' formula  is just what gives   \eqref{basic.i}.
Since $\sum_{j\ge 1} e^{- { 2n j^2/  d^2}}\le
Cd/\sqrt n $,   we  deduce  from Proposition \ref{Lemma44} the useful estimate below.
\begin{corollary}\label{c.useful.est}
\begin{equation*}   \sup_{1\le d\le n}\,\big|\P\big\{d|B_n\big\}- {1\over d}  \big| \, \le \,   {C \over\sqrt n} .
\end{equation*}
%\begin{eqnarray*}  \big|\P\big\{ d | B_n\big\} -\frac{1}{d}\big| \le   \frac{C}{\sqrt n},\end{eqnarray*} 
where $C$ is an absolute constant. 
\end{corollary}

  \bigskip \par Proposition \ref{Lemma44}  is   already sharp enough  to
imply  the   useful bound below.
 \begin{proposition} \label{unifd}  $$\sup_{n\ge 1} \sum_{d<\sqrt n} \big|\P\big\{  d|   B_{ n}
\big\}-{1\over d} 
\big|<\infty.$$
\end{proposition}\begin{proof}
  Let $A $ and   $m$ be positive {\it reals}. Then for $d\ge 1$,
$$ \int_{m\over d+1}^{m\over d}e^{-At^2}{  d t\over t}\ge e^{-A({m\over d})^2}\int_{m\over d+1}^{m\over d} {  d t\over t}=e^{-A({m\over
d})^2}\log (1+{ 1\over d})\ge  { 1\over 2d}e^{-A({m\over
d})^2} , $$
since $\log 1+x\ge \frac{x}{2}$, if $0\le x\le 1$. Writing  $n=\nu^2$ and letting $m=\nu j$, $A=2$  gives
$${1\over d}\sum_{j=1}^d e^{-2({\nu j\over
d})^2}\le 2 \sum_{j=1}^d  \int_{\nu j\over d+1}^{\nu j\over d}e^{-2t^2}{  d t\over t}.$$
Let 
 $ H= \sum_{1\le d<\nu} \big\{{1\over d}\sum_{1\le j\le d} e^{-2({\nu j\over
d})^2}\big\}$.  
By permuting the order of summation, we get
\begin{eqnarray*}H&\le& 2\sum_{1\le d<\nu}\Big\{\sum_{1\le j\le d}   \int_{\nu j\over d+1}^{\nu j\over d}e^{-2t^2}{  d t\over t}\Big\}
=2\sum_{1\le j<\nu}\sum_{j\le d<\nu} \int_{\nu j\over d+1}^{\nu j\over d}e^{-2t^2}{  d t\over t}
\cr &\le &2
\sum_{1\le j<\nu} \int_{  {\nu j\over \nu+1}  }^{\nu  }e^{-2t^2}{  d t\over t}=  2 \int_{  {\nu +1 \over \nu }}^{\nu  }\Big(\sum_{j=1}^{\min(\nu ,
{\nu +1 \over \nu }t)}1\Big)  e^{-2t^2}{  d t\over t}  
\cr &\le & 4\int_{{ 1 \over 2 } }^\infty  e^{-2t^2}  
d t  = C. \end{eqnarray*}
Consequently, by using (\ref{multiple1}),
 \begin{equation*}  
  \sum_{d<\sqrt n} \big|\P\big\{  d|   B_{ n}
\big\}-{1\over d} 
\big|\le 2\sum_{d<\sqrt n} {1\over d}\sum_{ 1\le j<d/2}
e^{-        {
 2n j^2/  d^2}  }\le 2H\le C, \end{equation*}
as claimed.\end{proof}


{Some applications}:
\vskip 3 pt
\noi(1) {\it An estimate  of $\E d(   B_{ n})$.} Proposition \ref{unifd}  implies the following estimate.
\begin{proposition}[Weber \cite{W11},\,Prop.\,6] \label{dbn} For  some universal constant $C$, we have 
$$ \sup_{n\ge 2}\, \big| \E d(   B_{ n})- \log n \big| \le C   .$$  \end{proposition}
We refer to the cited article  for the proof.
\smallskip \par
\noi This estimate is certainly improvable. The same question arises with $\E d^2(   B_{ n})$.  We remark that \begin{eqnarray}\E d^2(B_n)&=&
%\E d (B_n)+ 
\sum_{   d \le n\atop    \d  \le n}\P\{ [d,\d]|B_n\}\, =\, 
%\E d (B_n)+
 \sum_{ [d,\d]\le n}\P\{ [d,\d]|B_n\}
%\cr &=& \sum_{ u\le n}R(u)\P\{ u|B_n\}
%\cr &=& \E d (B_n)+2\sum_{  \d \le n}\ \sum_{ \d< d \le n}{1\over [d,\d]}+\  2\sum_{  \d \le n}\ \sum_{ \d< d \le n}\big(\P\{ [d,\d]|B_n\}- {1\over [d,\d]}\big)
.
\end{eqnarray}
%By Lemma \ref{unifd}  $$\sup_{n\ge 1} \sum_{[d,\d]<\sqrt n} \big|\P\big\{  [d,\d]|   B_{ n}\big\}-{1\over [d,\d]} \big|<\infty.$$
 A relevant result is    \begin{theorem}[Shi-Weber \cite{Sh.We},\,Th.\,2.1] \label{shi.weber}
  {\rm (i)}   If $0<\s< 1$, then
 \begin{eqnarray*}
 \sum_{ [d,\d]\le N}  \frac{1}{[d,\d]^{\s }}
 &=&\frac 3{\pi^2} \cdot \frac 1{1-\s}N^{1-\s} (\log N)^2+\mathcal O\left( N^ {1-\s}(\log N)\right).
\end{eqnarray*}
 
{\rm (ii)} Further, if $\s =1$,
 \begin{eqnarray*}
 \sum_{ [d,\d]\le N}  \frac{1}{[d,\d]}
 &=&\frac 1{\pi^2} (\log N)^3+\mathcal O\left((\log N)^2\right).
\end{eqnarray*}
\end{theorem}
\begin{remark}[Large zeta values imply large values of a very small selection of Dirichlet polynomials]\rm Combined with the arithmetic perturbation technic of Dirichlet polynomials     introduced in \cite{Sh.We},  Proposition 1.1, 
  large zeta values can be  compared with large values of a very small selection of Dirichlet polynomials. For any $N\ge1$,  and any  real numbers $\s,\e, t$,
  \begin{equation}\label{e2}\sum_{n=1}^N \frac{1}{n^{\s +it}}  =\sum_{[d,\d]\le N}   \frac{J_\e (d)J_\e (\d)}{[d,\d]^{{\s +2\e+it}}}
\sum_{m=1}^{ \frac{N}{[d,\d]}}            \frac{1} {m^{{\s+2\e +it}}}.
\end{equation}
%Proposition \ref{p1}-(ii) with Theorem \ref{l1}-(ii), we get as special case: 
    Let $\s_0>0$. %We have as special case (Proposition 1.1 and Theorem 2.1): 
    Thus by Theorem \ref{shi.weber}-(ii), for any   $\s\ge \s_0$,  $|t|\ge 1$, $\e> 0$,\begin{eqnarray*}\big|\zeta(1 + it) \big|&\le& C(\e)\, (\log |t|)^3  \max_{M\in\mathcal E({|t|})} \Big|\sum_{m=1}^{M} \frac{1}{m^{{1+2\e +it}}}\Big|+ C  |t|^{-\s},
\end{eqnarray*}
where $\mathcal E(|t|))\subset [1, |t|]$ and 
%  $\mathcal E(|t|)$ being the set of  integers of the form $\lfloor \frac{|t|}{d}\rfloor$, $d$ integer, $1\le d\le |t|$. Note that   
$\#(\mathcal E(|t|))=\mathcal O\big( \log |t|\big)$. Hence a much smaller selection than the one derived from Abel's summation. These pointwise implications extend to Dirichlet polynomials indexed on intervals, and have by integration standard in-measure counterparts. This is displayed elsewhere. These questions are relevant in the recent work of     Matom\"aki and  Ter\"av\"ainen \cite{MT}.
\end{remark}
\vskip 3 pt
\noi(2) {\it Infinite M\"obius inversion.}  The following application of Proposition \ref{dbn} to infinite M\"obius inversion is proved in \cite{W11}.
\begin{theorem}\label{t2} Let $g:\N\to \R^+$. Assume that  for some $0<\e<1$,
  % Let$ \{a_n , n\ge 1\} $ be a sequence of positive reals.
 %  Assume that 
\begin{equation}\label{locvar[g]}    \g_g(\e)= \sup_{n\ge 1} \sup_{\frac{|2m-n |}{\sqrt{n\log\log n}}\,\le 1+\e}\frac{ g(m) }{ g(n) }<\infty.
\end{equation}
Then the infinite M\"obius inversion 
\begin{equation}\label{HW} 
g(x) =\sum_{m=1}^\infty f(mx)\qq \Leftrightarrow \qq f(x) =\sum_{n=1}^\infty \m(n)g(nx),
\end{equation} holds for any $x$, as soon as
\begin{equation}\label{imi[g]} \sum_{n\ge 1} g(n)  \log n<\infty.\end{equation} 
% for any real $c>0$,
%\begin{equation}\label{slicereg1}  \g_1(c)=\sup_{   n\ge 1   }\sup_{ m\ge cn  }{a_{m}\over
%a_n}  <\infty.
%\end{equation}
%Then  \eqref{conv.}  implies \eqref{conv.concl}.
\end{theorem}
See also Weber \cite{W13}, Remark 2,  for a rectification  of  Theorem 270 in Hardy et   Wright \cite{HW}, more precisely, the assertion that {\it under} assumption $ {\sum_{m,n=1}^\infty} |f(mnx)|<\infty$,     the implication    $(\Longleftarrow)$  in \eqref{HW} holds,     is not clear.\medskip \par \noi
%\underline{Some applications}:
(3) {\it Estimates of $\E \s_{-1}(B_n)$ and $\P\big\{  B_n\,{\it and}\, B_m \ {\it coprime} \big\}$.}
    Recall that $\s_s(k)$ (section \ref{s2})  
 %     Let $\s_s(k) := \sum_{d|k} d^s$,  
 %$s\in \R$,\,  $s\not=0$ be the function sum of $s$-th powers of divisors. 
%This 
is a multiplicative function and  that $\s_{-s}(n)=n^{-s} \s_s(n)  $. \begin{corollary}
$$\Big|\E \s_{-1}(B_n)-   \zeta(2) \Big| \le\  C\,{  \log n\over\sqrt n} .$$
\end{corollary}



 

\begin{proof}Writing that $\s_{-1}(B_n)=\sum_{d\le n }\frac{1}{d}\chi\{d|B_n\}$ since $B_n\le n$, we have 
%$$  \s_{-1}(B_n)=\sum_{d| B_n}\frac{1}{d} =\sum_{d\le n\,,\,d|B_n}\frac{1}{d} =\sum_{d\le n }\frac{1}{d}\chi\{d|B_n\}  $$
$$  \E \s_{-1}(B_n)= \sum_{d\le n }\frac{1}{d}\,\P\{d|B_n\}.  $$
Now $$\E \s_{-1}(B_n)- \sum_{d\le n}{1\over d^2}=
\sum_{d\le n}{1\over d}\big( \P\big\{d|B_n\big\}  - {1\over d }\big). $$
Whence,
%By Corollary \ref{c.useful.est}, 
$$\big|\E \s_{-1}(B_n)- \sum_{d\le n}{1\over d^2}\big|\le
{C \over\sqrt n}\ \sum_{d\le n}{1\over d}\le\  C\,{  \log n\over\sqrt n} .$$
Consequently,
%$$  = \zeta(2) = \frac{\pi^2}{6},$$
  $$\Big|\E \s_{-1}(B_n)-   \zeta(2) \Big| \le\  C\,{  \log n\over\sqrt n} .$$
\end{proof}
\begin{corollary} We have for $n>m\ge 2$,
\begin{eqnarray*} \P\big\{  B_n\,{\it and}\, B_m \ {\it coprime} \big\}  &=&\sum_{d\le m\wedge (n-m)} {\m(d)\over d^2}\cr & &+\mathcal O\Big(   {\log m \over  \sqrt{m\wedge (n-m)}} +\Big({m\wedge (n-m)\over m\vee (n-m)}\Big)^{1/2}\Big).
\end{eqnarray*}
Moreover,
$$\lim_{m\to\infty\atop  {n }/{m} \to\infty}\P\big\{ (B_n,B_m)=1\big\}
 =\frac{1}{\zeta(2)}
 .$$
 \end{corollary}
% \frac{1}{\zeta(2)}=\frac{6}{\pi^2}
 \begin{proof} By using M\"obius inversion formula,
%  and again  Corollary \ref{useful.est},
$$\P\big\{ (B_n,B_m)=1\big\} =\sum_{d\le 
m\wedge (n-m)}\m(d)\P\{d|S_m\}\,\P\{d|S_{n-m}\}$$
$$=\sum_{d\le m\wedge (n-m)}\m(d)\Big(\P\big\{d|S_m\big\}-{1\over d}+{1\over d}\Big)\,\Big(\P\big\{d|S_{n-m}\big\}-{1\over d}+{1\over d}\Big) $$
$$=\sum_{d\le m\wedge (n-m)}\m(d)\Big(\P\big\{d|S_m\big\}-{1\over d}+{1\over d}\Big)\,\Big(\P\big\{d|S_{n-m}\big\}-{1\over d}+{1\over d}\Big) $$
$$=\sum_{d\le m\wedge (n-m)}\m(d)\Big( {1\over d}+\mathcal O\big({ 1\over\sqrt m}\big)\Big)\, \Big({1\over d}+\mathcal O\big({ 1\over\sqrt {n-m}}\big)\Big) $$
$$=\sum_{d\le m\wedge (n-m)} {\m(d)\over d^2}+\mathcal O\Big(\sum_{d\le m}
{1\over d }\big({1\over  \sqrt{ m}}+{1\over  \sqrt{ n-m}}\big)+\Big({m\wedge (n-m)\over \sqrt{m( n-m)}}\Big)\Big)$$
$$=\sum_{d\le m\wedge (n-m)} {\m(d)\over d^2}+\mathcal O\Big(   {\log m(\sqrt m+\sqrt{ n-m} )\over  \sqrt{ m(n-m)}} +\Big({m\wedge (n-m)\over m\vee (n-m)}\Big)^{1/2}\Big).$$
For instance: 
\vskip 2pt --- If  $(\log m)^2<n-m\le m$, 
\begin{eqnarray*} \Big|\P\big\{  B_n\,{\it and}\, B_m \ {\it coprime} \big\}  -\sum_{d\le  n-m } {\m(d)\over d^2}\Big|  & = & \mathcal O\Big(   {\log m \over  \sqrt{ n-m }} +\Big({ n-m\over m }\Big)^{1/2}\Big).
\end{eqnarray*}

--- If  $n\ge 2m$,  
\begin{eqnarray*} \Big|\P\big\{  B_n\,{\it and}\, B_m \ {\it coprime} \big\}  -\sum_{d\le m  } {\m(d)\over d^2}\Big|  & = & \mathcal O\Big(   {\log m \over  \sqrt{  m }} +\Big({  m\over n-m }\Big)^{1/2}\Big).
\end{eqnarray*}
This implies the second claim.\end{proof}
 
This suggests to study the following problem: Let $n_1<n_2<\ldots< n_k$ be positive integers such that 
$$\min\{ n_j/n_i, 1\le i<j \le k\}$$
is large. To estimate the probability that $B_{n_1}, B_{n_2},\ldots, B_{n_k}$ are mutually coprime, namely
$$\P\big\{ (B_{n_i}, B_{n_j})=1, \forall 1\le i<j \le k\big\}.$$ 
%This seems doable.
 That problem is related to a nice notion: for a subset of an integer lattice, to be {\it visible from the origin}, in this case the   set of random points   $(x,y)$ in $\N^2$, $x\neq y$, with coordinates in $\{B_{n_1}, B_{n_2},\ldots, B_{n_k}\}$. Referring to section 3.8 in Apostol \cite{A}, we say  that two lattice points $P$ and $Q$ are mutually visible  if the segment which join  them contains no lattice points other than the endpoints $P$ and $Q$. We also recall that two lattice points $(a,b)$ and $(m,n)$ are mutually visible if  and only if $(a-m, b-n)= 1$, or  equivalently  if and only if $(a-m,b-n)$ is visible from the origin. The standard case presented in \cite{A} has clear link with Farey series, and Dirichlet's easy   estimate on the number of corresponding visible  points is a particular case of a much more general result having a direct link with Riemann Hypothesis  (RH). See   \cite{W13} for more.
%and references therein.
\smallskip \par
 %As  an application, we prove  that    for any    $\eta>1$, any finite sequence of reals $\{c_k, k\in K\}$,
%$\sum_{k,\ell\in K} c_kc_\ell \, \frac{\gcd(k,\ell)^{2}}{k\ell} \le  C(\eta) \sum_{\nu\in K} c_\nu^2(\log\log\log \nu)^\eta \Psi(\nu)
%$,where $C(\eta)$ depends on $\eta$ only.  

 \begin{remark}\rm Let $X$ and $Y$ be independent random variables having geometric distribution with parameter $1-e^{-\b}$ with $\b>0$: $\P\{X=n\}=(1-e^{-\b})e^{-\b(n-1)}$ for all $n\in\N$. Bureaux and Enriquez \cite{BE} have shown that 
 that error term
 $$ \P\{(X,Y)=1\}= \frac{1}{\zeta(2)}+\frac{\b}{\zeta(2)}+\mathcal O(\b^{3/2-\e}), \quad \hbox{for each}\ \e>0\ \hbox{as}\ \b\downarrow 0, $$
 is equivalent to RH. 
 
 A natural question arises from this nice result. It seems possible (although we have not investigated it) to   estimate the probability $\P\{X\,\hbox{prime}\}$ with accuracy, up to some extend. Would the accuracy of the error term be also sensitive? 
 \end{remark}  

%\begin{remark}[Ramus' formula] Ramus proved in 1834 the formula (see also Problem 2(d) p.39 in the book of Riordan \cite{Ri}, 1958) $$\sum_{k=0\atop k\equiv m  (d)}^n  C^k_n = {1\over d}\sum_{k=1}^d \o^{-km}(1+\o^k)^n = {1\over d}\sum_{k=1}^d\big(2\cos {k\pi \over d}\big)^n\cos\big({(n-2m)k\pi\over d}\big),$$where  $\o$ is the $n$-th primitive root of the unit. Indeed, $(1+x)^\nu $ has derivative $\nu(1+x)^{\nu-1} =  \sum_{k=1}^{ \nu} kx^{k-1}C^k_{ \nu}   $. Letting $x=1$ yields$ \sum_{k=1}^{ \nu} k C^k_{ \nu}=\nu2^{ \nu}$. \end{remark}  
 
   
%      The following  estimate    is implicit   in the proof (see   p.9). Let $\a>1$ and put $\p_n= \sqrt{ \a n\log  n}$, $\D_n= \frac{n}{2}-\p_n$. Then
%\begin{eqnarray}\label{twosided}   \#\{d<\sqrt {\D_n}: d| B_{ n} 
% \} - \sqrt { n}\,\chi\{|B_{ n}-\frac{n}{2}|> \p_n\}
%&\le & \#\{d<\sqrt B_{ n}: d| B_{ n}   \} 
% \cr &\le& \#\{d<\sqrt{ n}: d| B_{ n}   \}.
%\end{eqnarray}  
%On the other hand for   any positive integer $m$,
%\begin{equation}\label{sqdiv}d(m)=2\#\{d<\sqrt m: d|m \} + \chi\{m\  \hbox{is square}\}, 
 %\end{equation}
 %so that 
% \begin{equation*} \Big|d(B_{ n})-2\#\{d<\sqrt{B_{ n}}: d|B_{ n} \} \Big|\le  \chi\{B_{ n}\  \hbox{square}\}, 
% \end{equation*}

% $$  \#\{d<\sqrt B_{ n}: d| B_{ n}   \}\le \#\{d<\sqrt{ n}: d| B_{ n}   \}$$
%\begin{eqnarray*}  \#\{d<\sqrt B_{ n}: d| B_{ n}  
% \} 
%&\ge&   \#\{d<\sqrt { \D_n}: d| B_{ n}  
% \}\cdot  \chi\{|B_{ n}-\frac{n}{2}|\le \p_n\}
%\cr &=&   \#\{d<\sqrt {\D_n}: d| B_{ n} 
% \} - \#\{d<\sqrt {\D_n}: d| B_{ n}\}\cdot  \chi\{|B_{ n}-\frac{n}{2}|> \p_n\}
%\cr
% &\ge& \#\{d<\sqrt {\D_n}: d| B_{ n} 
% \} - \sqrt { n}\,\chi\{|B_{ n}-\frac{n}{2}|> \p_n\}
%\end{eqnarray*}
% \begin{eqnarray}
% \label{dbn.demi}\frac{1}{2}\log n-C_1\le \E \#\{d<\sqrt B_{ n}: d| B_{ n}   \}  \le
% \frac{1}{2}\log n+C_2.
 %\end{eqnarray}  
%Combined with  \begin{equation}\label{sqdiv}d(m)=2\#\{d<\sqrt m: d|m \} + \chi\{m\  \hbox{is square}\}, 
% \end{equation}
%valid for   any integer $m$, and a separate estimate of $ \P\{ B_n\  \hbox{square}\}$, it allows to one conclude.

\subsection{The local limit theorem for binomial sums}
The first local limit theorem    was  established  nearly  three centuries ago, in the binomial case. We refer to     Szewczak    and  Weber \cite{SW}, Th. 1.1 for the statement   under the  form below. 
 \begin{theorem}[{\rm De Moivre--Laplace} (1733)]
\label{moivre}  Let
$0<\rho<1$. Let $X$ be such that
${\mathbb P}\{X=1\}=p=1-{\mathbb P}\{X=0\}
$. Let
$X_1, X_2,\ldots$ be independent copies of
$X$ and let $S_n=X_1+\ldots +X_n$, $n\ge 2$. Let also  $0<\g<1$.
Then for all $k$ such that $|k-np|  \le \g n\rho(1-\rho)$ and $    n \ge \max(\rho/(1-\rho),(1-\rho)/\rho)$, letting
$ x= \frac{k-n\rho}{\sqrt{  n\rho(1-\rho)}}$,
we have
 \begin{equation*} {\mathbb P}\{B_n(\rho)=k\}  \ =\
 \frac{e^{-  \frac{x^2}{  2  } }}{\sqrt{2\pi n\rho(1-\rho)}} \   e^E ,
 \end{equation*}
with
$|E|\le   \frac{3|x|+ 2 |x|^3}{(1-\g)\sqrt{ n\rho(1-\rho)}} +  \frac{1}{ 4n\min(\rho,1-\rho )(1 -  \g      )}$.

\vskip   3 pt  If $x= o(n^{1/6})$,  then
\begin{equation*}\Big| {\mathbb P}\{B_n(\rho)=k\}- \frac{1 }{\sqrt{2\pi n\rho(1-\rho)}}\,  e^{-  \frac{(k-n\rho)^2}{  2 n\rho(1-\rho) }}\Big|
\ \le  \   \frac{1 }{\sqrt{2\pi n\rho(1-\rho)}}\ e^{-  \frac{(k-n\rho)^2}{  2 n\rho(1-\rho) }}\Big( \frac{c_1 |x|^3}{\sqrt{ n}}+ \frac{c_2|x|}{\sqrt n} +  \frac{c_3}{ n}\Big).
\end{equation*}
The constants $c_i$, $i=1,2,3$ are explicit and depend on $p$ only. \end{theorem}

A useful  consequence is: {\it For any $b>0$, $n\ge 2$,}
\beq\label{binomial.llt.1}
  \sup_{k:|k-n\rho|\le \sqrt{b n\log n}}\Big|{\mathbb P}\{B_n(\rho)=k\}- \frac{1 }{\sqrt{2\pi n\rho(1-\rho)}}\,  e^{-  \frac{(k-n\rho)^2}{  2 n\rho(1-\rho) }}\Big|
\, \le  \,  C_b \Big( \frac{ \log^{3/2} n}{\sqrt n}\Big).  
\eeq 

\subsection{A concentration inequality   for binomial sums.}\label{sub2.2}
 The following  is Th.\,2.3   in MacDiarmid \cite{di}, which addresses a more general case.
\begin{lemma} \label{di.1}
Let $X_1, \dots, X_n$     be independent random variables, with $0 \le X_k \le 1$ for each $k$.
Let $S_n = \sum_{k=1}^n X_k$ and $\mu = \E S_n$. Then for any $\epsilon >0$,
 \begin{eqnarray*}
{\rm (a)}\quad  &&
 \P\big\{S_n \ge  (1+\epsilon)\mu\big\}
    \le  e^{- \frac{\epsilon^2\mu}{2(1+ \epsilon/3) } } ,
    \cr {\rm (b)} \quad & &\P\big\{S_n \le  (1-\epsilon)\mu\big\}\le    e^{- \frac{\epsilon^2\mu}{2}}.
 \end{eqnarray*}
\end{lemma}













\subsection{Sharper uniform estimate and  Theta     function.}













  One might believe in view of Proposition \ref{unifd} that the central comparison term  is $\frac1d$. But this is {\it not} so. The right comparison term -{\it with a sharper remainder term\,}- turns up to be more complicated, and is equal to ${\Theta(d,n)\over d}$ where $\Theta(d,n)$ is the Theta     function, 
  \begin{equation} \label{Theta}\displaystyle{ \Theta (d,n)  =  \sum_{\ell\in \Z} e^{in\pi{\ell\over   d }-{n\pi^2\ell^2\over 2 d^2}},     }
   \end{equation}
  $d\ge 1,\ n\ge 1$ being integers.  
\medskip \par The following Theorem proved in Weber \cite{W1}, 
provides a uniform estimate on the whole range of  divisors.
%, with a sharper rate of approximation.   
%\begin{equation}\label{unifdiv} \sup_{2\le d\le n}\Big|\P\big\{ d | B_n\big\}-  {1\over d}\sum_{ 0\le |j|<d/2 }
%  e^{i\pi n{j\over d}  -n  
%{\pi^2j^2\over 2d^2}}\Big|= {\mathcal O}\big((\log n)^{5/2}n^{-3/2}\big)
%.\end{equation}
  \begin{theorem}[\cite{W1},\,Theorem II]\label{estPdlBn} 
 % ([We6]  Theorem II)   
{  \rm (1)}    We have the following uniform estimate: 
\begin{equation}\label{unif.est.theta}\sup_{2\le d\le n}\Big|\P\big\{d|B_n\big\}- {\Theta(d,n)\over d}  \Big|= {\mathcal O}\Big(\frac{(\log n)^{5/2}}{n^{3/2}}\Big).  
\end{equation}
 
  
%{  \rm (2)}  Further, 
%\vskip 1pt \begin{equation}\label{(1.2)}
  %\big|\P\big\{d|B_n\big\}- {1\over d}  \big|\le
 %   \begin{cases}  C\Big(\frac{(\log n)^{5/2}}{n^{3/2}}+ { 1\over d}
 %e^{ - {n \pi^2 \over 2d^2}} \Big)   & \  {\it if}\  d\le \sqrt n,\cr 
  %    {C \over\sqrt n}   & \  {\it if} \  \sqrt n\le d\le n.
%\end{cases}
%\end{equation}
%Thus $\lim_{n\to \infty}\P\big\{d|B_n\big\}= {1/ d}  $, and we have the bound
%\begin{equation}\label{(1.5)}  \big|\P\big\{d|B_n\big\}- {1\over d}  \big| \le C  \, {  d\over n},
%\qq \qquad {\rm if}\ 2\le d\le \sqrt n.
%\end{equation}  Whence also,
%\begin{equation}\label{useful.est}   \sup_{1\le d\le n}\,\big|\P\big\{d|B_n\big\}- {1\over d}  \big| \, \le \,   {C \over\sqrt n} .
% \end{equation}

 
{  \rm (2)} Further,
  for any $\a>0$,  
\begin{equation}\label{unif.est.theta.alpha} \sup_{d<  \pi   \sqrt{   n \over 2\a\log n}}\big|\P\big\{  d|   B_{ n} \big\}-{1\over d} 
\big|= {\mathcal O}_\e\big(n^{-\a+\e }\big),\qq \quad (\forall \e>0)  
\end{equation}   
 and for any $0<\rho<1 $, 
\begin{equation}\label{unif.est.theta.rho}  \sup_{d<  (\pi/\sqrt 2) n^{(1-\rho)/2} }\big|\P\big\{  d|   B_{ n} \big\}-{1\over d} 
\big|= {\mathcal O}_\e\big(e^{-(1-\e) n^\rho}\big),\qq \quad (\forall 0<\e<1). 
\end{equation} 
 \end{theorem} 
\vskip 2 pt The uniform estimate \eqref{unif.est.theta} covers the   critical region of divisors $d\sim \sqrt n$, and   we don't know  any  other   estimate of this kind covering this case.   This estimate  also admits a version  for binomial sums, which follows by simply copying the proof made in the Bernoulli case.
 
 \begin{remark}\label{unif.est.theta.alpha.rem}\rm  The proof we give is intrinsic, in particular we don't see how to  obtain this estimate from an approach using the uniform model.
 \end{remark}
\begin{remarks} \label{rem.Ex.2}\rm  --- Note  the
considerable variation of the error term between \eqref{unif.est.theta.alpha} and \eqref{unif.est.theta.rho}. 
These latter estimates are easy to prove and are true for binomial sums.  
In fact (\cite{DS},\,Example 2),
 \begin{equation}\label{multiple} \Big|\P\big\{d|B_n(\varrho) \}- \frac{1}{d}\Big|\le 
 e^{-        {
 8n \varrho(1-\varrho)  /  d^2}  }.
\end{equation} 
\vskip 3 pt
 \noi ---  Lemma 2.1 in      Cilleruelo,   Fern\'andez and Fern\'andez \cite{CFF} is only a weak  form of it. 
\vskip 3 pt
\noi  --- Letting next $h(x)$ be increasing to infinity with $x$ and such that $h(x)\le  \sqrt x $, it follows that
 \begin{equation}\label{multiple.h} \sup_{d< \sqrt{  n/ h(n)}}\Big|\P\big\{d|B_n(\varrho) \}- \frac{1}{d}\Big|\le 
 e^{-    8  
   \varrho(1-\varrho) h(n)  }.
\end{equation} 
   For instance,  
for any $\a>0$,   
 \begin{equation}\label{unifdivalpha}  \sup_{d< \sqrt{  n/ (\a \log n)}}\Big|\P\big\{d|B_n(\varrho) \}- \frac{1}{d}\Big|\le 
 n^{-    8  \a
   \varrho(1-\varrho)  },  
\end{equation} and  for any $0<\b<1 $,   
\begin{equation}\label{unifdivro} \sup_{d<   n^{(1-\b)/2} }\big|\P\big\{  d|   B_{ n} \big\}-{1\over d} 
\big|\le e^{- 8   
   \varrho(1-\varrho) n^\b}. 
\end{equation}
\end{remarks} 

   \begin{remark} \rm    
Poisson summation formula: for   $x\in \R,\  0\le
\d\le 1
$, 
\begin{equation}\label{poisson}\sum_{\ell\in \Z} e^{-(\ell+\d)^2\pi x^{-1}}=x^{1/2} \sum_{\ell\in \Z}  e^{2i\pi \ell\d -\ell^2\pi x},
 \end{equation}  
 implies that   \begin{equation}\label{poisson.comp.theta.llt}
 \frac{\Theta(d,n)}{ d}  =\sqrt{  { 2\over \pi
n}}  \sum_{ 
z\equiv 0\, (d)} e^{-{ (2z-n)^2\over 2 n}}.
 \end{equation} 
 \end{remark}
 \begin{remark} \rm   Formulas \eqref{unif.est.theta}, \eqref{poisson.comp.theta.llt}
show that the general problem of estimating $\P\big\{d|B_n\big\}$  reduces to the one of   estimating the $d$ series
\begin{equation}\label{vd.div.series}
\sum_{j }e^{-\frac{(jd-\rho)^2}{2n}},
\end{equation}
the summation being either over   even rational numbers $j$, or over   odd rational numbers $j$,   $\rho$ denoting   the residue class modulo $d$ of $n$.

An indication of the behavior of these series when $d$ and $n$ {\it simultaneously} vary, is given by  the already sharp estimate, 
\begin{equation}\label{Theta.dmu} 
 \big|{\Theta (d,n)\over d}-{1\over d}\big| \ \le \ \begin{cases}   { C\over d}
 e^{ - {n \pi^2 \over 2d^2}} & \qq \text{   if}\quad d\le \sqrt n,\cr        {C \over\sqrt n}   & \qq \text{  if}\quad \sqrt n\le d\le n.
 \end{cases}
 \end{equation} 
    



We also quote the classical bound in the study of   inclusion theorems related to  summation methods. 
\begin{lemma} [\cite{H},\,Th.\,140] \label{exp.square}
\begin{equation}\sum_{k\in \Z} e^{-ck^2/n} =\sqrt{\Big(\frac{n\pi}{c}\Big)} +\mathcal O(1)
\end{equation}
when $n\to \infty$, uniformly in any interval of positive $c$.
%$$\pi x^{-1}=c /n\qq x=\frac{\pi n }{c}$$
\end{lemma}
Indeed by applying Poisson summation formula    with $x=\frac{\pi n }{c}$ and $\d=0$, \begin{equation} \sum_{k\in \Z} e^{-k^2\pi x^{-1}}=x^{1/2}   \sum_{k\in \Z}e^{-k^2\pi x}=\sqrt{\Big(\frac{n\pi}{c}\Big)} +\sum_{|k|\ge 1 } e^{-k^2n /c}  
    \end{equation} 
 We next use the  estimate below. Let $a$ be any positive real. Then,  
\begin{equation}\label{(2.17)} \sum_{H=1}^\infty e^{-aH^2}\le \ \ \begin{cases}   3e^{-a} & \qq {\rm  if}\quad \  a\ge 1,\cr
       {3 /\sqrt a}   & \qq {\rm  if}\quad  \ 0<a\le 1 .
       \end{cases}
       \end{equation}
Here $a=n/c\ge 1 $ for $n$ large, in which case
\begin{equation}\sum_{k\in \Z} e^{-ck^2/n}=\sqrt{\Big(\frac{n\pi}{c}\Big)} +\mathcal O(  e^{-n/c}), 
    \end{equation} 
 thereby implying the claimed estimate with a much better remainder.
 This   is however too weak to imply estimate \eqref{Theta.dmu} in which  $d$   and $n$ may vary simultaneously.
As to \eqref{(2.17)}, this is elementary since \begin{eqnarray*}
\sum_{H=1}^k e^{-aH^2}&\le &e^{-a } +\sum_{H=2}^k e^{-a H^2 }  \le  e^{-a } + \int_{1}^\infty
  e^{-au^2}du
 \cr 
   (u=v/\sqrt{2a})\qq &  =  & e^{-a } + {1\over \sqrt{ 2a}}\int_{\sqrt{ 2a}}^\infty
  e^{-v^2/2}dv \cr &\le  &e^{-a }\Big(1 + \sqrt{{\pi\over   2a }} \Big),  \end{eqnarray*}
 where   we used the Mill's ratio bound $ e^{ x^2/2} \int_x^\infty e^{-t^2/2}\, dt  \le \sqrt{{\pi\over 2}}     $,
 valid for all $x\ge 0$. Thus $\sum_{H=1}^\infty e^{-aH^2}\le e^{-a } (1 + \sqrt{{\pi/  2  }}  )\le    3e^{-a}$   if  $a\ge 1$, and
$\sum_{H=1}^\infty e^{-aH^2}\le   (1/ \sqrt a )\big(1 + \sqrt{{\pi/   2  }} \big)\le      3/ \sqrt a  $,  if  $a\le 1$.  
\end{remark} 

\vskip 4 pt
\begin{remark}\rm
A study of the special case   of square-free divisors  is made in Shi-Weber \cite{Sh.We1}.
 \end{remark} 

%\begin{corollary} \begin{equation}\label{Bn.prime.intro}\P\{ B_n\,\hbox{ prime} \} \,=\,  \sum_{d\le n\,,\,  P^+(d )\le \sqrt n}{  \m(d)\over d}\Theta(d,n)+ {\mathcal O}\Big({ \log^{5/2} n  \over n}\Big),\end{equation} 
%where $P^+(m )$  denotes the largest prime factor of $m$.
%\end{corollary}
\medskip \par
  In the same 2007's paper \cite{W1}, the distribution of the smallest prime divisor of $B_n$ was further studied.  Some notation: let $(a,b)$ denote the greatest common divisor of the integers $a$ and $b$. Recall that    $P^-(m )$ denotes    the smallest
prime divisor  of the integer $m>1$ (by convention $P^-(1)=+\infty$). We refer to   Tenenbaum \cite{Te}, III.6, for a study of integers free of   small divisors, and to the deep and instructive survey of Hildebrand and Tenenbaum \cite{HT}, and III.5 in the above, concerning    integers free of large divisors. 

More precisely, we    derived   from Theorem \ref{estPdlBn}  and elementary Erathost\`ene's sieve,   sharp estimates of  the probability $\P \{ P^-( B_n ) > \zeta  \}$,  
providing   the right central comparison term  together with an already   satisfactory and exploitable error term. %As a matter of interest, this also serves as an illustration of the application of this method. 
  \begin{theorem}[\cite{W1}, Th.\,I]\label{P-Bn}
  There exist a positive real $c>0$ and   constants $C_0$, $\zeta_0$ such that for $n$ large
enough,   we have the
 following estimate
$$
\Big|\hskip1pt\P\big\{ P^-( B_n ) > \zeta \big\}-     { e^{-\gamma} \over  \log
\zeta }\,\Big|  \le  { C_0 \over   \log^2 \zeta } \qq \qquad (\,\zeta_0\le  \zeta\le n^{c/\log\log n}\,) $$
 where    $\g$ is  Euler's constant.
\end{theorem}
\begin{remark}\rm Sharpenings of Theorem \ref{P-Bn} are the object of a separate work.
\end{remark} 
We illustrate this on
% by giving 
some examples.
% of applications.
 
 
%[[[
%$$\frac{\Theta(d,n)}{ d}  =\sqrt{  { 2\over \pi
%n}}  \sum_{ 
%z\equiv 0\, (d)} e^{-{ (2z-n)^2\over 2 n}}$$
 %$$\P\{ B_n\,\hbox{  prime} \} \,=\, \sqrt{  { 2\over \pi
%n}}   \sum_{d\le n\,,\,  P^+(d )\le \sqrt n}   \m(d)\sum_{ 
%z\equiv 0\, (d)} e^{-{ (2z-n)^2\over 2 n}}+ {\mathcal O}\Big({ \log^{5/2} n  \over n}\Big)$$
% ]]]
 
 \vskip 3 pt  \noi (1) {\it Estimate of $\P\{ B_n\,\hbox{  prime} \}$.}
We begin  with clarifying   the link between this probability and the Prime Number Theorem (PNT).
% Let  $2=p_1<p_2<\ldots  $ be the sequence of consecutive prime numbers. It is easy to show that if $\P\{ B_n\,\hbox{  prime} \}\ge (\log n)^{-a}$, $a>0$, then $g(k)=p_{k+1}-p_k$  verifies $g(k)= {\mathcal O}\big( \sqrt{k\log k \log\log k} \big)  $.According to a well-known result of   Cram\'er \cite{C},  the validity of RH would  imply that  $g(k)= {\mathcal O}\big(  k^{1/2}    (\log k)^{3/2}\big) $. Besides, the Cram\'er conjecture states that  $$\limsup_{k\to \infty}{g(k)\over (\log p_k)^2}=1 .$$  But in the actual state of knowledge on the PNT, this turns up of no peculiar interest.
%,  since  that probability cannot be estimated. This   is   easy to explain. 
We first remark that
$$\P\Big\{ \big|B_n-\frac{n}{2}\big|>\frac12 \sqrt{3n\log n}\Big\}\ll \frac{1}{n}.$$
 Next by the local limit theorem
 there exists a numerical constant $C_0$ such that for all positive $n$
 \begin{eqnarray*}  \sup_{k}\, \Big|  {\mathbb P}\big\{\mathcal B_n=k\}
 -\sqrt{\frac{2}{\pi n}} e^{-{ (2k-n)^2\over
2 n}}\Big|\ \le \
  \  \frac{C_0}{n^{3/2}}  .
  \end{eqnarray*}
 Accordingly, 
 $$\P\{ B_n\,\hbox{  prime} \}=\sqrt{\frac{2}{\pi n}} \sum_{p\in\mathcal P\cap [\frac{n}{2}-\frac12 \sqrt{3n\log n},\frac{n}{2}+\frac12 \sqrt{3n\log n}]}   e^{-{ (2p-n)^2\over
2 n}}+\mathcal O\Big(\frac{1}{n}\Big) .$$
  The summation index is of type $\mathcal P\cap [x, x+y]$, $y\asymp \sqrt{x\log x}$; thus its cardinality relies upon extensions of the PNT  and cannot in turn be estimated.
  The best   unconditional estimate (slightly improved by Heath-Brown) is due to Huxley who showed in \cite{Hu}    that the PNT extends to intervals of the type $[x,x+x^\t]$,   $x^{\frac{7}{12} }  \le \t\le \frac{x}{(\log x)^4}$, namely
\begin{equation*}
%\label{huxley} 
\#\{[x,x+x^\t]\cap \mathcal P\} \sim  {x^\t}/{ \log x }.
\end{equation*} 
 %,  even assuming the RH   where it is required that  $y\ge 2\sqrt{x }\log x$, see Heath-Brown \cite{HB}. 
  Assuming   RH, the best result known to us states as follows,
\begin{equation}\label{RH.pi}\pi(x)-\pi(x-y)\, =\, \int_{x-y}^x \frac{\dd t}{\log t}+ 
\mathcal O\Big( x^{\frac{1}{2}}\log \Big\{ \frac{y}{x^{\frac{1}{2}}\log x}\Big\}\Big)
\end{equation}
for $y$ in the range 
$ 2x^{\frac{1}{2}}\log x \le y\le x$. Hence for $M\ge 2$ fixed,
\begin{equation}\pi(x)-\pi(x-Mx^{\frac{1}{2}}\log x)
\sim\,  x^{\frac{1}{2}} \big\{M+ 
\mathcal O\big(  \log M\big)\big\}.\end{equation}
%Thus for $M\ge 2$ fixed,
%\begin{equation}\pi(x)-\pi(x-Mx^{\frac{1}{2}}\log x)
%\sim\,  x^{\frac{1}{2}} \big\{M+ 
%\mathcal O\big(  \log M\big)\big\}.\end{equation}

See Heath-Brown \cite{HB}.
%,  see also the recent paper \cite{HB1} and the references therein.
The interested reader will fruitfully refer to Montgomery and Soundararajan \cite{MS},\,p.\,595, also to Montgomery and Vaughan \cite{MV}, a model of   workbook. Accordingly any  result based on an assumption on the size's order of $\P\{ B_n\,\hbox{  prime} \}$ like for instance $\ge ( \log n)^{-\a}$, $\a>0$  -which has an easy interpretation on the gap between consecutive primes- is of no peculiar interest.
\vskip 3pt On the other hand, it follows from Remark \ref{rem.Snprime.M}   that  for some constant $C>0$,
\begin{equation}\label{bd1}(\log n) \, \P\{ B_n\,\hbox{prime} \}\ge C,
\end{equation}
for all   integers $n$ in a  set  of density 1.  \bigskip \par
  
  
  
 By using Theorem \ref{P-Bn}, we get  the following remarkable compact formula with sharp rate of convergence. 
\begin{theorem}\label{eta(n;y)Bprime}
 \begin{equation*}\P\{ B_n\,\hbox{\rm prime} \} \,=\,  \sum_{d\le n\,,\,  P^+(d )\le \sqrt n}{  \m(d)\over d}\Theta(d,n)+ {\mathcal O}\Big({ \log^{5/2} n  \over n}\Big).\end{equation*} 
 Further
\begin{equation*}\P\{ B_n\,\hbox{\rm prime} \} \,=\, \sqrt{  { 2\over \pi
n}}  \sum_{ k\in\Z}  \Big( \sum_{d\le n\atop  P^+(d )\le \sqrt n}   \m(d) \,e^{-{ (2kd-n)^2\over 2 n}}\Big)+ {\mathcal O}\Big({ \log^{5/2} n  \over n}\Big).\end{equation*} 
\end{theorem}  
%\begin{equation}\label{poisson.comp.theta.llt}
 %\frac{\Theta(d,n)}{ d}  =\sqrt{  { 2\over \pi
%n}}  \sum_{ 
%z\equiv 0\, (d)} e^{-{ (2z-n)^2\over 2 n}}.
% \end{equation}
 %It is implicit in the right-term that $d\le n$. 
We first prove
\begin{proposition} \label{p.eta(n;y)B}  For any positive integer $n$ and any positive real $y$,
 we have 
  \begin{equation}\label{eta(n;y)B.est}\P\{P^-(B_n )> y\} \,=\,  \sum_{ d\le n\,,\, P^+(d )\le y}{  \m(d)\over d}\Theta(d,n)+ {\mathcal O}\Big( { y\,\log^{5/2} n  \over n^{ 3/2}}\, 
e^{-\frac{\log n}{\log y}/2 }
\Big).
\end{equation} 
%Also for $n\ge y \ge 2$,
%\begin{eqnarray*} \P\{P^-(B_n )> y\}   &=& \frac{e^{-\g}}{\log y}\Big\{1+\mathcal O\Big(\frac{1}{\log y}\Big)\Big\} + \mathcal O\Big(  \frac{y}{\sqrt n}\, e^{-\frac12 \frac{\log n}{\log y}}\Big)   .
%\end{eqnarray*} 
%If $y= e^{\log^\b n}$,$\frac{\log n}{\log y}=  \log^{1-\b} n$.
  \end{proposition}
 
We need a Lemma.
 \begin{lemma} \label{l.eta(n;y)B}For any positive integer $n$ and any positive real $y$,
 we have
 \begin{equation*} \P\{P^-(B_n )> y\} =\sum_{  P^+(d )\le y} \m(d)\P\{ d|B_n \}
.
\end{equation*}
  %\label{eta(n;y)B}
 \end{lemma}
 
 
\begin{proof} Recall that the characteristic function of the set of integers $n$ such that $P^-(n )> y$,
\begin{equation*}\eta(n;y) =\begin{cases} 1 \quad & \ \hbox{ if $P^-(n )> y$} \cr
0 \quad & \ \hbox{ if $P^-(n )\le  y$}.
\end{cases}
\end{equation*}
is multiplicative. Thus by the   M\"obius inversion formula
\begin{equation}\label{eta(n;y)}\eta(n;y) =\sum_{d|n\atop P^+(d )\le y} \m(d) \qq \qq (n\ge 1).
\end{equation}
 See  Tenenbaum \cite[p.\,395-396]{Te1}. 
%-396
 %Montgomery and Vaughan, \cite{MV},\, th.\, 2.1, Apostol \cite{A},\,p.\,28.
  Therefore by integration the Lemma follows.   \end{proof}


\begin{proof}[Proof of Proposition \ref{p.eta(n;y)B}] A direct consequence of  Lemma \ref{l.eta(n;y)B} and Estimate  \eqref{unif.est.theta} of  Theorem \ref{estPdlBn} is  that for $n\ge y\ge 2$,
 \begin{equation*} \P\{P^-(B_n )> y\} \,=\,  \sum_{ d\le n\,,\, P^+(d )\le y}{  \m(d)\over d}\Theta(d,n)+ {\mathcal O}\Big(\sum_{ d\le n\,,\, P^+(d )\le y}{ | \m(d)|\over d}\,{ \log^{5/2} n  \over n^{ 3/2}}\, 
  \Big).
\end{equation*} 
\vskip 3 pt Now let $\Psi(x,y)= \#\{m\le x:P^+(m)\le y\}$. By Theorem 1 p.\,359 in \cite{Te1}, we have 
\begin{eqnarray*}
\sum_{ d\le n\,,\, P^+(d )\le y} |\m(d)|\,\le \, \Psi (n,y) \,\ll \, y\,e^{-(\frac{\log n}{\log y})/2}\qq \quad (n\ge y\ge 2).\end{eqnarray*}
Therefore 
%(zzz verifier car avant $- \frac{\log n}{\log y}$)
\begin{equation*} \P\{P^-(B_n )> y\} \,=\,  \sum_{ d\le n\,,\, P^+(d )\le y}{  \m(d)\over d}\Theta(d,n)+ {\mathcal O}\Big(y\,{ \log^{5/2} n  \over n^{ 3/2}}\, 
e^{-\frac{\log n}{\log y}/2 }
\Big).
\end{equation*} 
%Now recall  that (Montgomery and Vaughan, \cite{MV} and\, th.\,2.1 in  Apostol \cite{A},\,p.\,28),    
 % $$  \sum_{  d|n  } \frac{\m(d)}{d}=  \frac{\p(n)}{n}=\prod_{ p|n   }\big(1-\frac1p\big),$$
 %  $\p(n)$ being Euler totient function,  we get \begin{eqnarray*} \P\{P^-(B_n )> y\} &=& \sum_{  d|\prod_{p\le y} p  } \frac{\m(d)}{d}manque \Theta + \mathcal O\Big( \sum_{d\le n\atop  P^+(d )\le y}\big|\P\{ d|B_n \}-\frac{1}{d}\big|\Big)   
%\cr &=& \prod_{ p|\prod_{p\le y} p   }\big(1-\frac1p\big)  + \mathcal O\Big(  \frac{1}{\sqrt n} \sum_{ d\le n\atop P^+(d )\le y}1\Big)   
%\cr &=& \prod_{ p\le y   }\big(1-\frac1p\big)  + \mathcal O\Big(  \frac{y}{\sqrt n}\,e^{-\frac12\,\frac{\log n}{\log y}}\Big)   
%\cr &=&\frac{e^{-\g}}{\log y}\Big\{1+\mathcal O\Big(\frac{1}{\log y}\Big)\Big\} + \mathcal O\Big(  \frac{y}{\sqrt n}\,e^{-\frac12\,\frac{\log n}{\log y}}\Big) ,
%\end{eqnarray*}
%if $n\ge y \ge 2$. 

 



%Dirichlet mult [M9e], (76)
%\begin{eqnarray} \sum_{y\le  n\le
%x \atop P^+(n)\le y  }  {1   \over n }&\le & C \Big[     e^{-{1\over 2}{\log x\over \log y}}+   \log y  \Big]\le C\log y .
 %\end{eqnarray}

% \begin{equation}\label{p.c.t.}
% \frac{\Theta(d,n)}{ d}  =\sqrt{  { 2\over \pi
%n}}  \sum_{ 
%\ell \in \Z} 
%%e^{-{ (2d\ell-n)^2\over 2 n}}
%\sum_{h=0}^\infty \frac{(-1)^h}{h!}  \Big({ (2d\ell-n)^2\over 2 n}}\Big)^h.
% \end{equation}
 
 % \begin{equation*}
%%\label{eta(n;y)B1}
%   \sum_{ d\le n\,,\, P^+(d )\le y}\m(d)\,{ \Theta(d,n)\over d}= \sqrt{  { 2\over \pin}}  \sum_{ \ell \in \Z} 
%%e^{-{ (2d\ell-n)^2\over 2 n}}
%\sum_{h=0}^\infty \frac{(-1)^h}{( 2 n)^hh!}  \sum_{ d\le n\,,\, P^+(d )\le y}\m(d)    (2d\ell-n)^{2h} \,.
%\end{equation*}  
%By Proposition \ref{unifd}, 
%$$\sup_{n\ge 1} \sum_{d<\sqrt n} \big|\P\big\{  d|   B_{ n}
%\big\}-{1\over d} 
%\big|<\infty.$$

%Tenenbaum \cite{Te},\, th.\,3.5.1.
\end{proof}
%\begin{equation}\label{Theta.d.} 
 %\big|{\Theta (d,n)\over d}-{1\over d}\big| \ \le \ \begin{cases}   { C\over d}
 %e^{ - {n \pi^2 \over 2d^2}} & \qq {\rm  if}\quad d\le \sqrt n,\cr&\cr
  %     {C \over\sqrt n}   & \qq {\rm  if}\quad \sqrt n\le d\le n.
 %\end{cases}
% \end{equation}

% $$ rajout$$
%\begin{eqnarray*}
%& &\sum_{ d\le n\,,\, P^+(d )\le y}\m(d)\,{ \Theta(d,n)\over d}\,=\,\sum_{ d\le \sqrt n\,,\, P^+(d )\le y} \,{ \m(d)\over d} 
%\cr & & \quad +\mathcal O\Big(   \sum_{ d\le \sqrt n\,,\,  P^+(d )\le y} {\,e^{ - {n \pi^2 \over 2d^2}}\over d}
 %  \Big)+ \mathcal O\Big(  {1\over \sqrt n} \sum_{ d> \sqrt n\,,\, P^+(d )\le y}  1
 %  \Big)+.\end{eqnarray*}
 %$$ rajout$$

%By using Proposition \ref{uniftheta}, we get 
%\begin{equation}\label{eta(n;y)B}\P\{P^-(B_n )> y\} \,=\,  \sum_{ d\le n\,,\, P^+(d )\le y}{  \m(d)\over d}\Theta(d,n)+ {\mathcal O}\Big({ \log^{5/2} n  \over n^{ 3/2}}\sum_{ d\le n\,,\, P^+(d )\le y} |\m(d)|\Big),
%\end{equation}
\begin{proof}[Proof of Theorem \ref{eta(n;y)Bprime}]It follows from Proposition \ref{p.eta(n;y)B} by taking $y=\sqrt n$, and  since $B_n$ is prime  if and only if $P^-(B_n )> \sqrt n$, that 
\begin{equation*}
%\label{eta(n;y)B.est}
\P\{ B_n\,\hbox{  prime} \} \,=\,  \sum_{ d\le n\,,\, P^+(d )\le \sqrt n}{  \m(d)\over d}\Theta(d,n)+ {\mathcal O}\Big( {  \log^{5/2} n  \over n }  
 \Big).
\end{equation*} 
   The proof is completed by  applying \eqref{poisson.comp.theta.llt}.
\end{proof}
 
%Let $\Pi_y$ denote the product of the prime numbers less of equal to $y$. We note that
%\begin{equation*}\#\{ d\le n\,;\, \m(d)=1 :P^+(d )\le y\}
%=\#\{ d\le n  :d|\Pi_y\}\le d(\Pi_y)=2^{\pi(y)},\end{equation*} 
%which is interesting for small values of $y$.

\vskip 9 pt

%Analogously to Proposition \ref{p.eta(n;y)B} we have

%\begin{proposition} \label{l.eta(n;y)B.pair} {\rm (a)} For any positive integers $n,m$ and any positive reals $y,z$,
% \begin{eqnarray}\label{eta(n;y)B}\P\{P^-(B_n )> y\,,\,P^-(B_m )> z\} &=&
 %\P\{P^-(B_n )> y\}\, \P\{P^-(B_m )> z\}
% %\Big(\sum_{  P^+(d )\le y} \m(d)\P\{ d|B_n \} \Big) \Big(\sum_{  P^+(\d )\le z} \m(\d)\P\{ \d|B_m \} \Big)
%\cr & & +  \sum_{  P^+(d )\le y\atop P^+(\d )\le z } \m(d) \m(\d) \D\big((d,n)\,,\, (\d,m)\big).
%\end{eqnarray}
%where $\D\big( (d,n), (\d, m)\big)$ is defined in \eqref{Delta.def}.
%\vskip 2 pt {\rm (b)} In particular, 
%\begin{eqnarray}\label{bn.bm.prime}\,\P\{B_n,B_m\,\hbox{ are prime}\} &=&
% \P\{B_n\,\hbox{  prime}\} \, \P\{B_m\,\hbox{  prime}\}
% %\Big(\sum_{  P^+(d )\le y} \m(d)\P\{ d|B_n \} \Big) \Big(\sum_{  P^+(\d )\le z} \m(\d)\P\{ \d|B_m \} \Big)
%\cr & & +  \sum_{  P^+(d )\le \sqrt n\atop P^+(\d )\le \sqrt m } \m(d) \m(\d) \D\big((d,n)\,,\, (\d,m)\big).
%\end{eqnarray}
%\end{proposition}
%\begin{proof}By \eqref{eta(n;y)}
%\begin{equation}\label{eta(n;y)aa}\eta(B_n;y)\eta(B_m;z) =\sum_{d|B_n\atop P^+(d )\le y} \m(d)\sum_{\d|B_m\atop P^+(\d )\le z} \m(\d).
%\end{equation}
%On integrating we get
%\begin{eqnarray}\label{eta(n;y)aaa}& &\P\{P^-(n )> y,P^-(m )> z\} =\sum_{  P^+(d )\le y\atop  P^+(\d )\le z} \m(d)  \m(\d)\P\{d|B_n,\d|B_m \}
%\cr &=&\sum_{  P^+(d )\le y\atop  P^+(\d )\le z} \m(d)  \m(\d)\P\{d|B_n\}\P\{\d|B_m \}
%\cr & & + \sum_{  P^+(d )\le y\atop  P^+(\d )\le z} \m(d)  \m(\d)\Big( \P\{d|B_n\,,\,\d|B_m \}- \P\{d|B_n\}\P\{\d|B_m \}\Big)
%\cr &=&\P\{P^-(n )> y\}\P\{P^-(m )> z\}
 % + \sum_{  P^+(d )\le y\atop  P^+(\d )\le z} \m(d)  \m(\d) \D\big((d,n)\,,\, (\d,m)\big).\end{eqnarray}
% This proves assertion (a).   By choosing $y=\sqrt n$, $z=\sqrt m$, assertion (b) follows immediately. 
% \end{proof}
%In particular  for $n$ large (as $\sqrt {n+2}-\sqrt {n }=o(1)$),
%\begin{eqnarray}\label{twin.prime}\,\P\{B_n\,\hbox{and}\,B_{n+2}\,\hbox{prime}\} &=&
% \P\{B_n\,\hbox{prime}\} \, \P\{B_{n+2}\,\hbox{prime}\}
%\cr & & +  \sum_{  P^+(d )\le \sqrt n\atop P^+(\d )\le \sqrt {n} } \m(d) \m(\d) \D\big((d,n)\,,\, (\d,n+2)\big).\end{eqnarray}
%$$[[\P\{B_n\,\hbox{and}\,B_{n+2}\,\hbox{prime}\}=\frac12\P\{B_n=2\}+\frac14\P\{B_n\,\hbox{prime}\}+\frac14 \P\{B_n\,\hbox{and}\,B_{n }+2\,\hbox{prime}\}$$ 
%$$ \sum_{  P^+(d )\le \sqrt n\atop P^+(\d )\le \sqrt {n } } \m(d) \m(\d)\Big( \P\{d|B_n\,,\,\d|B_{n+2} \}- \P\{d|B_n\}\P\{\d|B_{n+2}\}\Big)$$
%$$= \sum_{  P^+(d )\le \sqrt n\atop P^+(\d )\le \sqrt {n } } \m(d) \m(\d)\Big( \P\{d\d|B_n B_{n+2} \}- \P\{d|B_n\}\P\{\d|B_{n+2}\}\Big)]]$$

% See also Lemma \ref{lem61},  Remark \ref{ind} and Theorem \ref{t5}, notably  for estimates of $\D\big( (d,n), (\d, m)\big)$.
 
 
%Taking   $y= \sqrt n$, $z=\sqrt m$, we get 
%\begin{eqnarray}\label{bn.bm.prime}\,\P\{B_n,B_m\,\hbox{ are prime}\} &=&
% \P\{B_n\,\hbox{ is prime}\} \, \P\{B_m\,\hbox{ is prime}\}
% %\Big(\sum_{  P^+(d )\le y} \m(d)\P\{ d|B_n \} \Big) \Big(\sum_{  P^+(\d )\le z} \m(\d)\P\{ \d|B_m \} \Big)
%\cr & & +  \sum_{  P^+(d )\le \sqrt n\atop P^+(\d )\le \sqrt m } \m(d) \m(\d) \D\big((d,n)\,,\, (\d,m)\big).
%\end{eqnarray}
 

  Let $a$ be real and  $\s_{a}(n)= \sum_{d|k} d^a  $.  In a similar way we get for $a=-1$,%\begin{equation*}
%\label{poisson.comp.theta.llt}
%\frac{\Theta(d,n)}{ d}  =\sqrt{  { 2\over \pin}}  \sum_{ z\equiv 0\, (d)} e^{-{ (2z-n)^2\over 2 n}}. \end{equation*} }
  \begin{theorem}\label{sigma-1.Bn}
\begin{equation*}\E \s_{-1}( B_n) \,=\,   \sqrt{  { 2\over \pi
n}} \sum_{d\le n}  \frac1d \sum_{ 
z\equiv 0\, (d)} e^{-{ (2z-n)^2\over 2 n}} 
% \sum_{d\le n}{ \Theta(d,n)\over d^2} 
+ {\mathcal O}\Big({ \log^{3} n   \over n^{3/2}}\Big) .\end{equation*} 
%[[[$$\frac{\Theta(d,n)}{ d}  =\sqrt{  { 2\over \pi
%n}}  \sum_{ 
%z\equiv 0\, (d)} e^{-{ (2z-n)^2\over 2 n}}$$$${\color{blue}
%\sum_{d\le n}{ \Theta(d,n)\over d^2} =  \sqrt{  { 2\over \pi
%n}} \sum_{d\le n}  \frac1d \sum_{ 
%z\equiv 0\, (d)} e^{-{ (2z-n)^2\over 2 n}} \le e^\gamma \log\log n}$$]]]
\end{theorem}
%We don't know another way to prove this.
%Robin showed that the 
%RH  is true if and only if $$  \s_{-1}(n)  < e^\gamma\log\log n,$$
%for $n>5040$, where   $\gamma$ is Euler's constant. Here we have an estimate of $\E \s_{-1}( B_n)$. Can we find a path from it to $\s_{-1}(n)$? Yes, at least theoretically. The answer is given in the next sub-section.

   
 \vskip 4 pt \noi (2) {\it Estimates related to Erd\"os and  Zaremba function.}
Let $\Phi(n)=\sum_{d|n} \frac{\log d}{d}$.  This function appears in the study of good lattice points in numerical integration. Erd\"os and  Zaremba  \cite{EZ} showed that
$$ \limsup_{n\to \infty} \frac{\Phi(n)}{(\log\log n)^2}=e^\g,$$ $\g$ being Euler's constant. The proof
   is based on the identity
 \begin{eqnarray}\label{formule}\Phi(n)
=\sum_{i=1}^r \sum_{\nu_i=1}^{\a_i}\frac{\log p_i^{\nu_i}}{p_i^{\nu_i}}\sum_{\d|n p_i^{-\a_i}}\frac{1}{\d} , \qq\qq ( n= p_1^{\a_1}\ldots p_r^{\a_r}),
\end {eqnarray}

 Consider  the function 
 $$\Psi(n)= \sum_{d|n} \frac{(\log d )(\log\log d)}{d} .$$  In this case  a  formula similar to \eqref{formule}   no longer holds, the "log-linearity" being  lost due to the extra factor $\log\log d$. 
We showed in the recent work  \cite{W11}
that
  \begin{eqnarray*}  \limsup_{n\to \infty}\, \frac{\Psi(n)}{(\log \log n)^2(\log\log\log n)}\,=\, e^\g.\end{eqnarray*}
%The study of $\Psi(n)$ much involves Davenport's function  $w(n)=\sum_{p|n}\frac{\log p}{p}$.
  
%devising
This required  a new approach, and the proof is considerably more complicated. 
  \begin{corollary}
 $$  \E \Psi(B_n)= \sum_{d\le n} \frac{(\log d )(\log\log d)\Theta(d,n)}{d^2} + {\mathcal O}\Big(\frac{(\log n)^{7/2}(\log\log n)}{n^{ 3/2}}
   \Big)$$
  \end{corollary}
  \begin{proof}Writing $\sum_{d| B_n} =\sum_{d\le n\,,\,d|B_n}$ we get,
$$\E \Psi(B_n)=\sum_{d\le n}  \frac{(\log d )(\log\log d)}{d}\big(\P\big\{d|B_n\big\}-{\Theta(d,n)\over d} +{\Theta(d,n)\over d}\big).$$
 Whence by Theorem \ref{estPdlBn},$$  \E \Psi(B_n)= \sum_{d\le n} \frac{(\log d )(\log\log d)\Theta(d,n)}{d^2} + {\mathcal O}\Big(\frac{(\log n)^{7/2}(\log\log n)}{n^{ 3/2}}
 %\sum_{d\le n} \frac{(\log d )(\log\log d)}{  d  } 
 \Big)
  . $$
 
\end{proof}

%Recall a well-known device.
%\begin{lemma}[Kronecker's lemma]\label{kl}
%Let $a_n\uparrow \infty$ and the series $\sum_{n=1}^\infty x_n/a_n$ converges. Then $(1/a_n)\sum_{k=1}^n  x_k\to 0$.
%\end{lemma}
%Apply it with the choice  $x_n=\chi_{\{P^-(B_n)>\zeta_n\}}-\P\{P^-(B_n)>\zeta_n\}$, $a_n=\sum_{k\le n}\P\{P^-(B_k)>\zeta_k\}$, $n$ running in a partial index $\mathcal N$. If $\sum_{k\ge 1}\P\{P^-(B_k)>\zeta_k\}=\infty$, then
 %$$\sum_{n=1}^\infty \frac{\chi_{\{P^-(B_n)>\zeta_n\}}-\P\{P^-(B_n)>\zeta_n\}}{\sum_{k\le n}\P\{P^-(B_k)>\zeta_k\}}  $$
 
%$$\P\{P^-(B_k )> \zeta_k\} =\sum_{  P^+(d )\le \zeta_k} \m(d)\P\{ d|B_k \}
%$$
\subsection{A  way back along Pascal's matrices.}\label{sub.bin.way.back}
 Let $f(j)$ be an arithmetical function. Consider the column vector ${\boldsymbol \nu}=\big(f(1),f(2),\ldots, f(n)\big)$. Let $P_n$ be the $n$-size   lower triangular Pascal matrix:
\begin{eqnarray*}P_n={ \left[\begin{matrix}
1        \cr
1   &1       \cr
1        &2            &1    \cr
1       &3       &3 &1       \cr
 \vdots         &   \vdots                  &\vdots   &\vdots   &\ddots    
 \cr {{n}\choose{1}}& {{n}\choose{2}}&{{n}\choose{3}}&{{n}\choose{4}}  &\cdots  &{{n}\choose{n}}\end{matrix} \right]}
 %\qq j=1,\ldots, n.
\end{eqnarray*}
Then $P_n{\boldsymbol \nu}=  \,{\boldsymbol w}$ where ${\boldsymbol w}=\big(\E f(B_1),2\E f(B_2),\ldots, 2^n\E f(B_n)\big)$. Call and Velleman (\cite{CV},\,Th.\,1)  
 identified $P^{-1}_n$, which is $DP_nD^{-1}$ with $D$ as indicated:
\begin{eqnarray*}P^{-1}_n={ \left[\begin{matrix}
1        \cr
-1  & 1       \cr
1        &-2            &1    \cr
-1       &3       &-3 &1       \cr
 \vdots         &   \vdots                  &\vdots   &\vdots   &\ddots    
 \cr {{n}\choose{1}}& {{n}\choose{2}}&{{n}\choose{3}}&{{n}\choose{4}}  &\cdots  &{{n}\choose{n}}\end{matrix} \right]}
 %\qq j=1,\ldots, n.
\end{eqnarray*}
\begin{eqnarray*}D={ \left[\begin{matrix}
1        \cr
  &-1       \cr
       &{}      &1    \cr
  &{}          &{}         &-1       \cr
 &{}          &{}   &{}          &{}  &\ddots    
 \cr &{}          &{}   &{}          &{}    &   &(-1)^{n+1}\end{matrix} \right]}
 %\qq j=1,\ldots, n.
\end{eqnarray*}
 So that knowing ${\boldsymbol w}$  we have -at least theoretically- a way back to know ${\boldsymbol \nu}$. In  condensed form     
$P_n^{-1}=Q_n$,  where  
$Q_n$ is the $n\times n$ matrix defined by   \begin{equation}Q_n(i,j)=\begin{cases} (-1)^{i-j}  {{i-1} \choose{j-1}} \quad&\hbox{if $i\ge j$,}\cr
 0 \quad&\hbox{otherwise.}
 \end{cases}
 \end{equation}  
   \begin{proposition}\label{way.back}Let $f(j)$ be an arithmetical function. Then
\beq\label{} f(i)=\sum_{j=1}^i  (-1)^{i-j}2^j{{i-1}\choose{j-1}}\E f(B_j),\qq\quad i\ge 2.
\eeq \end{proposition}
  
%\beq\label{} f(i)=\sum_{j=1}^i  (-1)^{i-j}2^j{{i-1}\choose{j-1}}\E f(B_j).
%\eeq
%\beq\label{} \s_{-1}(i)=\sum_{j=1}^i  (-1)^{i-j}2^j{{i-1}\choose{j-1}}\E \s_{-1}(B_j).
%\eeq
  
  
  
  
  
   %%%%%%%%%%%%%%%%%%%%%%%%%%%%%%%%%%%%%%%%%%%%%%%%%%%%%%%%%%%%%%%%%%%%%%%%%%%%%%%%%%%%%%%%%%%%%%%%%%%%%%%%%%%%%%%%

 
 

 %%%%%%%%%%%%%%%%%%%%%%%
%%%%%%%%%%%%%%%%%%%%%%%
 \subsection{Extension to all residue classes mod\,($d$).}  
%
     Theorem \ref{estPdlBn} was recently  extended to all residue classes mod\,($d$) in   Weber \cite{W9} with application to divisors in iid square integrable random walks.
    % This amounts  to study the  value distribution of divisors of $B_{ n}+u$,   $u $ being any non-negative integer, is necessary. We recently established  the validity of such extension.
   The  estimate   obtained  is   uniform over  all residue classes. 
     %We begin with a preliminary  observation. The restriction $d\le n$ in the estimate of Theorem \ref{estPdlBn} is superfluous. That  estimate is in turn also uniform over all $d\ge 2$.  
% If $d>n$, as the first term is 0,   \eqref{unif.est.theta} provides an   estimate  of   the central term ${\Theta(d,n)/ d} $, for $d>n$. 
 
%This point   has some  degree of importance.   The proof of the Theorem below  is   similar to the one of Theorem \ref{estPdlBn}, which corresponds to the case $u=0$. We refer to \cite{W9}, for the proof. %This term  in fact  equals (see proof of  Theorem II in \cite{W1}) to   $\sum_{ 1\le j<d/2}(\cos{ \pi n{j\over d}})e^{ -n   {\pi^2j^2/2d^2}} $, up to the error term   $Ce^{- {\pi^2n / 72}}$,  $C$ being absolute, uniformly for all $d\ge 2$.   
 \begin{theorem}[\cite{W9}, Th.\,3.1]\label{estPdlBn.u} 
%For any integer $u$, 
There exist  two absolute constants $C$ and $n_0$ such that for all $n\ge n_0$,
\begin{equation*}
 \sup_{u\ge 0}\,\sup_{d\ge 2}\Big|\P\big\{  d|   B_{ n}+u \big\}-  {1\over d}\sum_{ 0\le |j|< d }
e^{i\pi (2un){j\over d}}\  e^{  -n  
{\pi^2j^2\over 2d^2}}\Big|\le C\, (\log n)^{5/2}n^{-3/2}.
\end{equation*}

\end{theorem}

  Put 
 \begin{equation}\label{theta.u.}
\Theta_u(d,n)  =  \sum_{\ell\in \Z}  e^{i\pi (2u+n){j\over d}}\  e^{  -n  
{\pi^2j^2\over 2d^2}}.
 \end{equation}
 Note that $\Theta_0(d,n)=\Theta(d,n)$.
 %Noticing that $$   \sum_{ \ell\in\Z}
%e^{i\pi (2u+n){\ell\over d}}\  e^{  -n  
%{\pi^2\ell^2\over 2d^2}} =e^{   u  
%{\pi^2\ell^2\over  d^2}}\Theta (d,2u+n),    $$
As a corollary we get,
\begin{corollary}\label{cor.estPdlBn.u} 
For some absolute constant $C$, we have  
\begin{equation*} \sup_{u\ge 0}\, \sup_{d\ge 2}\Big|\P\big\{  d|   B_{ n}+u \big\}-{\Theta_u(d,n)\over d} \Big| \le C \,(\log
n)^{5/2}n^{-3/2} ,
\end{equation*}
for all $n\ge n_0$.
%\begin{equation*}
%\sup_{u\ge 0} \sup_{2\le d\le n}\Big|\P\big\{  d|   B_{ n}+u \big\}- {1\over d}-{2 \over d}\,\Theta_+ (d,(2u+n))\Big|= {\mathcal O}\big((\log n)^{5/2}n^{-3/2}\big).
%\end{equation*}
\end{corollary}
   % The role of these symmetry properties is important.
    The proof is based on the lemma below. 
    %the proof being elementary, however it is necessary to display calculations, and in order to spare it to the reader  we   included a proof. Otherwise the proof of Corollary \ref{cor.estPdlBn.u}  
% the proof   is similar to that of Theorem II in \cite{W1}.
    % but  We first operate a reduction due to symmetries.% in order to make it rigorous,  there is no other way  than detailing the calculations.

 \begin{lemma}\label{reduction.sym} For any integers $d\ge 2$, $n\ge 2$ and $u\ge 0$, 
 \begin{eqnarray*}
  \P\big\{ d| B_n+u\big\}&=& {1\over d} + {2\over d}\sum_{1\le j< d/2} \cos\big(\pi (2u+n){j\over d} \big)\big(\cos { \pi j\over d}\big)^{n}
.
\end{eqnarray*}
 \end{lemma}
 
%\begin{remark}
% In the recent work   \cite{SW}, a similar but more complicated reduction step is made    in a study   of  squarefree divisors of Bernoulli sums.
%\end{remark}

% \begin{proof} Let $d=2\d$  with $\d \ge 1$. Then 
% \begin{equation}\label{basica}
% \P\big\{ 2\d | B_n+u\big\}  =  {1\over 2\d}\sum_{j=0}^{\d-1} e^{i\pi (2u+n){j\over 2\d}}\big(\cos { \pi j\over 2\d}\big)^{n}
%  +{1\over 2\d}\sum_{j=\d}^{2\d-1} e^{i\pi (2u+n){j\over 2\d}}\big(\cos { \pi j\over 2\d}\big)^{n}
%\end{equation}
 
%Letting first $j=2\d-v$, $v=1,\ldots, \d$ in the last sum, we have  
% \begin{eqnarray*}\sum_{j=\d}^{2\d-1} e^{i\pi (2u+n){j\over 2\d}}\big(\cos { \pi j\over 2\d}\big)^{n}
%&=&\sum_{v=1}^\d  e^{i\pi (2u+n){2\d-v\over 2\d}}\big(\cos { \pi(2\d-v)\over 2\d}\big)^{n}
%\cr &=&\sum_{v=1}^\d  e^{i  (2u+n)(\pi -{\pi v\over 2\d})}\big(\cos (\pi -{\pi v\over 2\d})\big)^{n}. 
%\end{eqnarray*}
%But 
%\begin{equation}\label{formula.sym}e^{i  m(\pi-x)} \cos^n    (\pi-x)  =(-1)^me^{-i 
%mx}  (-1)^n\cos^n   x= (-1)^{m+n}e^{-i 
%nx}  \cos^n   x  .
%\end{equation} Thus
%\begin{eqnarray*}\sum_{j=\d}^{2\d-1} e^{i\pi (2u+n){j\over 2\d}}\big(\cos { \pi j\over 2\d}\big)^{n}
 % &=&\sum_{v=1}^\d  e^{i  (2u+n)(\pi -{\pi v\over 2\d})}\big(\cos (\pi -{\pi v\over 2\d})\big)^{n}
%  \cr &=&
 % \sum_{v=1}^\d  e^{-i  (2u+n) {\pi v\over 2\d}}\big(\cos {\pi v\over 2\d}\big)^{n}
 % \cr &=&
 % \sum_{v=1}^{\d -1} e^{-i  (2u+n) {\pi v\over 2\d}}\big(\cos {\pi v\over 2\d}\big)^{n}, 
%\end{eqnarray*}
%since for $v=\d$, we have $e^{-i  (2u+n) {\pi v\over 2\d}}\big(\cos {\pi v\over 2\d}\big)^{n}=e^{-i  (2u+n) {\pi  \over 2 }}\big(\cos {\pi  \over 2 }\big)^{n}=0$.
%\vskip 3 pt 
%Carrying   this back to \eqref{basica} gives,
% \begin{eqnarray}\label{basicb}
%& & \P\big\{ 2\d | B_n+u\big\} \,=\, {1\over 2\d}\sum_{j=0}^{\d-1} e^{i\pi (2u+n){j\over 2\d}}\big(\cos { \pi j\over 2\d}\big)^{n}
 %\cr & &\quad+{1\over 2\d} \sum_{j=1}^{\d-1}  e^{-i  (2u+n) {\pi j\over 2\d}}\big(\cos {\pi j\over 2\d}\big)^{n}
% \cr &=&{1\over 2\d}+ {1\over 2\d} \sum_{j=1}^{\d-1} \Big\{e^{i\pi (2u+n){j\over 2\d}}+e^{-i  (2u+n) {\pi j\over 2\d}}\Big\}\big(\cos {\pi j\over 2\d}\big)^{n}
% \cr &=&{1\over 2\d}+{2\over 2\d}\sum_{j=1}^{\d-1} \cos \big(\pi (2u+n){j\over 2\d}\big)\big(\cos { \pi j\over 2\d}\big)^{n}.
%\end{eqnarray}

%Now let  $d=2\d+1$ with $\d \ge 1$ and write that
% \begin{eqnarray}\label{basicc.}
% \P\big\{ 2\d+1 | B_n+u\big\}&=& {1\over 2\d+1}\sum_{j=0}^{\d} e^{i\pi (2u+n){j\over 2\d +1}}\big(\cos { \pi j\over 2\d +1}\big)^{n}
%\cr & &\quad +{1\over 2\d+1}\sum_{j=\d+1}^{2\d} e^{i\pi (2u+n){j\over 2\d +1}}\big(\cos { \pi j\over 2\d +1}\big)^{n}.
%\end{eqnarray} 

%With the variable change $j=2\d +1-v$,  $v=1, \ldots \d$,    the second sum writes using \eqref{formula.sym},
%\begin{eqnarray}\label{basicd.}
 % \sum_{j=\d}^{2\d} e^{i\pi (2u+n){j\over 2\d +1}}\Big(\cos { \pi j\over 2\d}\Big)^{n}
%&=&  \sum_{v=1}^{\d} e^{i\pi (2u+n){(2\d +1-v)\over 2\d +1}}\Big(\cos { \pi (2\d +1-v)\over 2\d+1}\Big)^{n}
%\cr &=&    \sum_{v=1}^{\d} e^{-i\pi (2u+n){v\over 2\d +1}}\big(\cos { \pi v\over 2\d+1}\big)^{n}.
%\end{eqnarray} 
%Thus 
%\begin{eqnarray}\label{basicc}
%& & \P\big\{ 2\d +1| B_n+u\big\}\,=\, {1\over 2\d+1}\sum_{j=0}^{\d} e^{i\pi (2u+n){j\over 2\d +1}}\big(\cos { \pi j\over 2\d +1}\big)^{n}
%\cr & &\quad +{1\over 2\d+1}\sum_{j=1}^{\d} e^{-i\pi (2u+n){j\over 2\d +1}}\big(\cos { \pi j\over 2\d+1}\big)^{n}
%cr &=& {1\over 2\d+1} + {1\over 2\d+1}\sum_{j=1}^{\d} \Big\{ e^{i\pi (2u+n){j\over 2\d +1}}+e^{-i\pi (2u+n){j\over 2\d +1}}\Big\}\big(\cos { \pi j\over 2\d +1}\big)^{n}
%\cr &=& {1\over 2\d+1} + {2\over 2\d+1}\sum_{j=1}^{\d} \cos\big(\pi (2u+n){j\over 2\d +1} \big)\big(\cos { \pi j\over 2\d +1}\big)^{n}
%.
%\end{eqnarray} 
%From \eqref{basicb} and \eqref{basicc.} follows that
%\begin{eqnarray}\label{basicd}
%  \P\big\{ d| B_n+u\big\}&=& {1\over d} + {2\over d}\sum_{1\le j< d/2} \cos\big(\pi (2u+n){j\over d} \big)\big(\cos { \pi j\over d}\big)^{n}
%.
%\end{eqnarray} 
%\end{proof}




%\begin{remark}
%For any $\alpha\in (0,1)$, there is a constant $C_\a$ such that, for integers $n\ge 1$, $d\le n$, and $r\in \{0,1,\ldots, d-1\}$, there holds
%$$\Big|\sum_{l\equiv r\,{\rm mod}\,d}   {n \choose l} \alpha^l (1-\alpha)^{n-l} -\frac1d \Big|\le \frac{C_\a}{\sqrt n}.$$
% [Gould Volume 6, formula (6.20)]
%$$ \sum_{l\equiv r\,{\rm mod}\,d}   {n \choose l} \alpha^l (1-\alpha)^{n-l}  = \frac1d  \sum_{k=0}^{d-1}(\o_d^k)^{-r}\big[(1-\a)+\a \o_d^k]^n,$$
%where $\o_d=e^{2\pi i/d}$. \end{remark}











  
\section{Divisibility properties in  i.i.d.   random walks.} \label{s4}
    Theorem \ref{estPdlBn} extends to     
i.i.d.  square integrable random walks.  \begin{theorem}[Weber \cite{W9},\,Th.\,2.1]\label{uniform.semi.local.limit.theorem} Let $X  $  be  a square integrable  random variable
   %  with basic probability space $(\O, \A, \P)$,
      taking values in the  lattice $\mathcal L(v_0,1)=\big\{v_k=v_0+ k,k\in \Z\big\}$, and let $f(k)=\P\{X=v_k\}$, for each $k$.   
Assume that 
   \beq \label{bp}\t_X   :=\sum_{k\in \Z}\min(f(k), f(k+1))>0.
   \eeq  
   Let $  X_i$, $i\ge 1 $  be  independent, identically distributed random variables having same law than $X$, and let $S_n=\sum_{j=1}^nX_j$, for each $n$.  
\vskip 2 pt   Let $ \bar \m= \{ \m_k, k\in \Z\}$ be a   sequence of non-negative reals such that  $0<\m_k < f(k) $, if $f(k)>0$,   $\m_k=0$ if $f(k)=0$,    let  $ \m= \sum_{k\in \Z}\m_k$, and assume that $1- \m< \t_X$.   
   Let also non-negative reals $\bar \tau=  \{ \tau_k, k\in \Z\}$ be  solutions of the (solvable) equation:   
 $ {\tau_{k-1}+\tau_k\over 2} =f(k)  - \m_k,
 $ for all $k \in \Z$.  
Let   $\t = \sum_{k\in \Z}  \tau_k = 1- \m.$ Further let $s(t) =\sum_{k\in \Z} \m_k\, e^{ 2i \pi   v_kt}$ and    $\rho$ be such  that  $1-\t<\rho<1$.   
  \vskip 1 pt      Let $\widetilde X$ be the   random variable associated  to $X$  and $\bar \m$, and defined by the relation
 $ \P\{ \widetilde X =v_k\}= {\tau_k}/{\t}$, for all  $k\in \Z   $.
 %Let $\mathcal F=\{F_{d},   d\ge 2\}$ be  class where $F_{d}= d\N$.
 Then the following results hold:
\vskip 2 pt {\rm (i)}
There exists  
$\theta=\theta(\rho,\t)$ with  $ 0< \theta <\t$,  
$C$  and $N$ 
 such that for $  n \ge N$,
\begin{align*}   \sup_{u\ge 0}\,\sup_{d\ge 2} \Big| \P \{ d|S_n+u  \}   -  {1\over d}-  {1\over d}\sum_{ 0< |\ell|<d }& \Big( e^{ (i\pi {\ell\over   d }-{ \pi^2\ell^2\over 2 d^2}) } \t\,
 \E \,e^{2i \pi {\ell\over   d }\widetilde X  }   +s\big(  {\ell\over   d }\big)\Big)^n \Big|
  \cr  &\le    \frac{C }{ \theta ^{3/2}}\   \frac{(\log  n)^{5/2}}{ n^{3/2}}+2\rho^n.
\end{align*}
 \vskip 2 pt   {\rm (ii)} Let   $\mathcal D$ be a test set of divisors $\ge 2$,   $\mathcal D_\p$ be the section of $\mathcal D$ at height $\p$  and $|\mathcal D_\p|$ denote its   cardinality. Then,
    \begin{eqnarray*}
    %\label{est.1.}
      \sum_{n=N}^\infty \   \sup_{u\ge 0} \, \sup_{\p\ge 2}\, {1\over |\mathcal D_\p |} \sum_{d\in \mathcal D_\p } \,\Big| \P \{d|S_n+u  \}   - {1\over d}\Big|
    & \le  &    \frac{C_1}{\t}
 \,  
         +  \frac{C_2 }{ \theta ^{3/2}} +\frac{2\rho^2}{1-\rho},
\end{eqnarray*}
where $C_1=\frac{2e^{     { \pi^2/ 4 }  }}
     {(1-e^{   -    {  \pi^2/ 16} }   ) }$, $C_2= C\,  \sum_{n=N}^\infty  \frac{(\log {n})^{5/2} }{ ({n})^{3/2}}$\,.
     \end{theorem} 
    Condition \eqref{bp}  defines an already very large class of random variables. A random variable $X$ is said to have a {\it Bernoulli part} if  condition \eqref{bp} is satisfied.  That notion is due to  McDonald \cite{M}, who introduced an ingenious coupling method  in the study of a famous limit theorem: the local limit theorem  which has  interface  with Number Theory.
    %the applicability of the local limit theorem  to this class.
    This method is particularly relevant in our proof. Concretely, let $  \{ \tau_k, k\in \Z\}$     of   non-negative reals such that
\begin{equation}\label{basber0}  \tau_{k-1}+\tau_k\le 2f(k), \qq  \qq\sum_{k\in \Z}  \tau_k =\t.
\end{equation}
For instance  
 $\tau_k=  \frac{\t}{\nu_X} \, \min(f(k), f(k+1))  $ is suitable.
   Define   a pair of random variables $(V,\e)$   as follows:
  \begin{eqnarray}\label{ve} \qq\qq\begin{cases} {\mathbb P}\{ (V,\e)=( v_k,1)\}=\tau_k,      \cr
 {\mathbb P}\{ (V,\e)=( v_k,0)\}=f(k) -{\tau_{k-1}+\tau_k\over
2}    .  \end{cases}\qq (\forall k\in \Z)
\end{eqnarray}
 Then we have the following decomposition which iterates well to iid sums $S_n$.    \begin{lemma} \label{bpr} Let $L$
be a Bernoulli random variable    which is independent of  $(V,\e)$, and put
 $Z= V+ \e  L$.
We have $Z\buildrel{\mathcal D}\over{ =}X$.
\end{lemma}

 
\begin{remark}[{\tf Reduction of binomial case   to  standard Bernoulli case.}] \rm Obviously binomial random variables have a Bernoulli part.  For  binomial random variables, the above decomposition is   direct. 
Let  $\b $  be a binomial random variable with  $\P\{\b  =1\}= \a = 1-\P\{\b  =0\}$. Let also $\varsigma$  be standard Bernoulli random variable  ($
\P\{\varsigma  =1\}= 1/2=\P\{\varsigma  =0\}$). 

Assume $0<\a<1/2$.  Let $\e $    be such that    $\P\{\e  =1\}= 2\a = 1-\P\{\e  =0\}$ and $\b, \e, \varsigma$ are independent. Trivially  $\e\varsigma \buildrel{\mathcal L}\over{=}\b$. 

 If now $1/2<a<1$, let $\tau_0 $
be  verifying
$0<  \tau_0< 2(1-\a)=2\min (\a,1-\a)$. 
Define a pair of random variables
$(V,\e)$   as follows.  
\begin{eqnarray*}   \begin{cases} \P\{ (V,\e)=(1,1)\}=0  
      \cr   \P\{ (V,\e)=(1,0)\}=\a -{\tau_0 \over
2}    . \end{cases}  \qq \begin{cases}   \P\{ (V,\e)=(0,1)\}=\tau_{0 } 
      \cr 
 \P\{ (V,\e)=( 0,0)\}=1-\a -{ \tau_0\over
2}    . \end{cases} 
\end{eqnarray*}
  Let $ \varsigma$ be  independent from $(V,\e)$.     Then 
$V+\e  \varsigma\buildrel{\mathcal L}\over{=}\b$.  

\end{remark}
  
   Consequently:     
    {\it      The study of the  limit properties (CLT,\, LLT,\,...) of the sequence
$S_n= k_1\b_1+
\ldots + k_n\b_n$, $n\ge 1$,  where 
$\b_j$ are independent binomial random variables, and   $\{k_j, j\ge 1\}$   a given increasing sequence of positive integers,   reduces to the case when these random variables are   i.i.d      standard Bernoulli.}
     
     
     
     
     
     
     
     
     
     
     
     
     
     
     
     
     
     
     
     
     
          
     
     \section{Divisibility properties of  non i.i.d. square integrable random walks.} \label{s5}
      Let $X=\{X_i , i\ge 1\}$ be  a sequence of  independent  variables taking values in $\Z$, 
  %a common lattice $\mathcal L(v_{0},D )$, namely defined by the
 %sequence $v_{ k}=v_{ 0}+D k$, $k\in \Z$, where $v_{0} $ and $D >0$ are   integers. 
and let $S_n=\sum_{k=1}^n X_k$, for each $n$. 
To our knowledge there is no study of divisors, in particular of their value distribution, in general independent   random walks. However divisibility is present in the study   the local limit theorem.  The sequence $X$ is  \emph{asymptotically uniformly distributed}, in short a.u.d., if   for any $d\ge 2$ and $m = 0,1,\ldots,d-1$, we
have
 \begin{equation} \label{aud1}\lim_{n\to \infty}\ \P\{S_n \equiv m \,{\rm (mod)}\,d\}=\frac1d.
  \end{equation}
 
  Let us assume   that the random variables $X_i$ take values in a common integer lattice $\mathcal L(v_{0},D )$, namely defined by the
 sequence $v_{ k}=v_{ 0}+D k$, $k\in \Z$,   $D >0$, and are  square integrable,  
and let     
\begin{equation} \label{MnBn}M_n= {\mathbb E\,} S_n , \qq B_n^2={\rm Var}(S_n)\to \infty.
\end{equation}
 
  We say  that the local limit theorem (in the usual form) is \emph{applicable} to   $X$    if
 \begin{equation}\label{def.llt.indep}    \sup_{N=v_0n+Dk }\Big|B_n\, {\mathbb P}\{S_n=N\}-{D\over  \sqrt{ 2\pi } }e^{-
{(N-M_n)^2\over  2 B_n^2} }\Big| = o(1), \qq \quad n\to\infty.
\end{equation} %We refer to our joint monograph   with    Szewczak.
%   \section{Divisibility and local limit theorem.} \label{s5}
%The following notion plays an important role in the study of the local limit theorem.   %an important notion in the  study  of the local limit theorem, in particular in the well-known Rozanov's necessary condition.    The
 % The sequence  $X$
%$\{S_n, n\ge 1\}$
 %is said to be {asymptotically uniformly distributed with respect
%to lattices of span $d$}, in short a.u.d.($d$), if for $m = 0,1,\ldots,d-1$, we
%have
 %\begin{equation} \label{aud1}\lim_{n\to \infty}\ \P\{S_n \equiv m \,{\rm (mod)}\,d\}=\frac1d.
  %\end{equation}
%Equivalenty  for $m = 0,1,\ldots,d-1$, we
%have 
% \begin{equation}\label{uad.lim1}\lim_{n\to \infty}\ \P\{d|S_n-m\}=\frac1d.
 %\end{equation} 
 

 
%As it will be clear along the next lines, this notion plays an important role in the theory of the local limit theorem. 
%Dvoretsky and Wolfowitz \cite{DW} proved the following characterization. Assume that $X$ is composed with independent random variables taking   only the values 
%$$ 0, 1,\ldots, h-1.$$
 % In order that the partial sums $\{S_n, n\ge 1\}$  be a.u.d.($h$), it is necessary and sufficient that
%\begin{equation} \label{aud.dw.ns} \prod_{k=1}^\infty\bigg( \sum_{m=0}^{h-1}\P\{X_k=m\}\,e^{\frac{2i\pi }{h}rm}\bigg) \,=\, 0, \qq \quad (r=1,\ldots, h-1).
% \end{equation}
%Equivalently,
%\begin{equation} \label{aud.dw.ns.} \prod_{k=1}^\infty\big(\E  e^{\frac{2i\pi }{h}rX_k}\big) \,=\, \lim_{N\to \infty}  \big(\E  e^{\frac{2i\pi }{h}rS_N}\big) \,=\, 0, \qq \quad (r=1,\ldots, h-1).
%\end{equation}

By   well-known Rozanov's result \cite{Ro}, 
%for a sequence $X=\{X_i , i\ge 1\}$ of  independent  variables taking values in $\Z$,  and    $S_n=\sum_{k=1}^n X_k$, for each $n$, 
the local limit theorem is applicable to   $X$  {\it only} if $X$ satisfies the a.u.d. property.
  
  

%By definition a local limit theorem is applicable to the sequence $X$, if
 %It is well-known that the local limit theorem is applicable to $X$   only if  
 
 %That well-known limit theorem hides a curious fact  and in the same time raises some question. 
%When the random variables $X_i$ are identically distributed,    \eqref{def.llt.indep} reduces to   \begin{equation}\label{llt.iid}    \sup_{N=v_0n+Dk }\Big|  \s \sqrt{n}\, {\mathbb P}\{S_n=N\}-{D\over  \sqrt{ 2\pi } }e^{-
%{(N-n\m)^2\over  2 n\s^2} }\Big| = o(1),
%\end{equation}
%where  $\m={\mathbb E\,} X_1$, $\s^2={\rm Var}(
%X_1)$. By Gnedenko's Theorem \cite{G}, see also \cite{P}, p.\,187,  \cite{SW}, Th.\,1.4, \eqref{llt.iid}   holds
%if and only if the span $D$ is maximal (there are no  other real numbers
%$v'_{0}
%$ and
%$D' >D$ for which
%${\mathbb P}\{X
%\in\mathcal L(v'_0,D')\}=1$).

%Note that  the transformation
%\begin{equation}\label{llt.transf.}
% X'_j= \frac{X_j-v_0}{D},
 % \end{equation}
%allows one to reduce  to the case  $v_0=0$, $D=1$.
   \vskip 10pt     In a recent work     based on   Poisson summation formula, 
the stronger result    is proved,  establishing an explicit link between the local limit theorem and  the a.u.d. property, through a     quantitative   estimate of  the difference ${\mathbb P} \{ S_n\equiv\,  m\! \hbox{\rm{ (mod $h$)}} \}- {1}/{h}$. 


  
\begin{theorem}[Weber \cite{W10},\,Th.\,1.4] \label{l1a}    Let $X=\{X_i , i\ge 1\}$ be  a sequence of  independent  variables taking values in $\Z$,  
and let $S_n=\sum_{k=1}^n X_k$, for each $n$. Assume that, 
\vskip 3 pt {\rm ($\mathcal A$) \it For some function $1\le \phi(t)\uparrow \infty $ as $t\to \infty$, and some constant $C$, we have for all $n$} 
 \begin{equation}\label{phi.cond}
\sup_{m\in \Z}\Big|B_n\P\big\{ S_n=m\big\}- {1\over   \sqrt{ 2\pi } }\ e^{-
{(m-M_n)^2\over  2 B_n^2} }\Big|\,\le \,  {C\over \,\phi(B_n)}.
 \end{equation}
%  The local limit theorem is applicable to   $\{S_n, n\ge 1\}$  only if the a.u.d. property is satisfied. 
\vskip 3 pt   Then there exists a numerical constant $C_1$, such that for all $0<\e \le 1$, all  $n $ such that $B_n\ge 6$, and all $h\ge 2$, 
\begin{align*}
\sup_{\m=0,1,\ldots, h-1} \,&\Big|{\mathbb P}\big\{ S_n\equiv\, \m\ \hbox{\rm{ (mod $h$)}}\big\}- \frac{1}{h}\Big|
\cr   &\le   {1\over \sqrt{2\pi}\,  B_n   }+\frac{2C}{h\,\sqrt{\e}\,\phi(B_n)}
+  {\mathbb P}\Big\{   \frac{|S_n -M_n |}{B_n}> \frac{1}{\sqrt \e}\Big\}+C_1 \,e^{-1/(16\e)}.
\end{align*}
\end{theorem}
    It follows from the proof that $C_1=2e\sqrt{\pi}$ is suitable.
    Choosing $\e= \phi(B_n)^{-2/3}$ and using Tchebycheff's inequality, we get the following
 \begin{corollary}\label{}For all   $n $ such that $B_n\ge 6$, and all $h\ge 2$, 
we have 
\begin{align}\label{eps.phi}
\sup_{\m=0,1,\ldots, h-1} \,  \Big|{\mathbb P}\big\{ S_n\equiv\,  \m\ \hbox{\rm{ (mod $h$)}}\big\}- \frac{1}{h}\Big|
 \le H_n ,
\end{align}
with
\begin{align}\label{eps.phi.Hn}
   H_n=  {1\over \sqrt{2\pi}\,  B_n   }+\frac{1+ 2 {C}/{h}
}{ \phi(B_n)^{2/3} } 
+  C_1 \,e^{-(1/ 16 )\phi(B_n)^{2/3}}.
\end{align}
\end{corollary}
 Theorem \ref{l1a}  implies Rozanov's result,  since by  definition such a function $\phi$ exists   if the local limit theorem is applicable to   $X$.          
 
\bigskip\par
  Instead of assumption ($\mathcal A$), consider now the following assumption  
  \vskip 5 pt ($\mathcal B$)    {\it $\t_{X_j} =\displaystyle{\sum_{k\in \Z}}{\mathbb P}\{X_j= k\}\wedge{\mathbb P}\{X_j= k+1 \} >0$,  for each $j$, and 
 $\nu_n =\displaystyle{\sum_{j=1}^n} \t_j\uparrow \infty$}.
 
   \vskip 3 pt We have the following sharp results.
 \begin{theorem}[\cite{W10},\,Th.\,4.3]\label{saud1} Let $D=1$. Suppose that assumption ($\mathcal B$) is fulfilled.   Let $\a\!>\!\a'\!>\!0$, $0\!<\!\e\!<\!1$. Then for   each $n$ such that 
 %${\sin x\over
%x}\ge (\a^\prime/\a)^{1/2}$
 %$  \tau_n\ge (\a^\prime/\a)^{1/2}$ where
% by \eqref{phi.tau}
%if 
$$|x|\le\frac12 \sqrt{  \frac{ 2\a\log (1-\epsilon)\nu_n}{ (1-\epsilon)\nu_n }}\qq \Rightarrow \qq{\sin x\over
x}\ge (\a^\prime/\a)^{1/2},$$
%then ${\sin x\over
%x}\ge (\a^\prime/\a)^{1/2}$,
%  \begin{equation*} \p_n=   \big( {2\a\log B_n \over B_n}\big)^{1/2},
%\qquad\qquad \tau_n= {\sin\p_n/2\over
%\p_n /2},
%\leqno(2.4)
%\end{equation*}
   we have
 \begin{eqnarray*} \sup_{u\ge 0}\,\sup_{d<  \pi  \sqrt{   (1-\epsilon)\nu_n \over 2\a\log (1-\epsilon)\nu_n}} \ \Big| \P \{d|S_n+u  \}   -  {1\over d} \Big|
    &\le &2 \,e^{- \frac{\epsilon^2 }{2}\nu_n}+
 \,\big( (1-\epsilon)\nu_n\big)^{-\a'}  .
\end{eqnarray*}
\end{theorem}
   
 \begin{theorem}[\cite{W10},\,Th.\,4.6]\label{saud2} Let $D=1$. Suppose that assumption ($\mathcal B$) is fulfilled. Let $0<\rho<1 $ and   $0<\e<1$. Then for   each $n$ such that 
$$|x|\le\frac12 \,\sqrt{  \frac{ 2 }{ ((1-\epsilon)\nu_n)^{1-\rho} }}\qq \Rightarrow \qq{\sin x\over
x}\ge \sqrt{1-\e}$$
%For each $n$ such that $\widetilde\tau_n\ge \sqrt{1-\e}$, where
%$$  \widetilde\tau_n= {\sin\psi_n/2\over
%\psi_n /2}\qq \qq \psi_n= \big({2B_n^\rho \over B_n}\big)^{1/2},$$  
we have\begin{eqnarray*} 
\sup_{u\ge 0}\,\sup_{d<  (\pi/\sqrt 2) ((1-\e)\nu_n)^{(1-\rho)/2} }\ \big| \P \{d|S_n+u  \}   -  {1\over d} \big|  
   &\le  &  2 e^{- \frac{\epsilon^2 }{2}\nu_n}+e^{- ( (1-\epsilon)\nu_n)^\rho}
  .
\end{eqnarray*}
\end{theorem}
 \vskip 2 pt We refer the reader to \cite{W10} for the proofs of these results.
 We also refer to our  recent monograph   with    Szewczak \cite{SW} for more.
 
 

 
\section{Distribution of   divisors of ${ B_nB_m}$.}\label{s6}
This task is more complicated.
%The study of the divisors of $B_nB_m$ is   a  more complicated  enterprise  than the one of the divisors of   $B_n$.  
 A key function is  the  (multiplicative) function
\begin{eqnarray}\label{rho}\varrho_k(D):=\#\big\{1\le y\le D: \  D| y^2+ky \big\} ,\qq   k=0,1,\ldots
\end{eqnarray} 
 which  it is not so surprising in view of the elementary decomposition
$B_nB_m= B_n^2 +B_n(B_m-B_n)$. 
\vskip 2 pt
As  mentioned in section \ref{s3}, that study is used in   
  \cite{W14} in the  control of   the correlation function \eqref{delta.def}, more precisely in the case      $n<m\le n+n^c$, $c>0$, $c$ being small, where the following intermediate bound appears (hence also $\varrho_k(.)$),
\begin{eqnarray*}\P \{ 
d|B_n\, ,\, \d|B_m \}
& =&\P \{  d|B_n\, ,\, \d|B_m\, ,\,B_m-B_n=0 \}
 +\P \{ 
d|B_n\, ,\, \d|B_m\, ,\,B_m-B_n>0 \}
\cr & =&2^{-(m-n)}\,\P \{  d|B_n\, ,\, \d|B_n \}  
 +\P \{ 
d|B_n\, ,\, \d|B_m\, ,\,B_m-B_n>0 \}
\cr &\le& 2^{-(m-n)}\,\P \{ [d, \d]|B_n  \}
+\P \{ 
d\d|B_n B_m\, ,\,B_m-B_n>0 \}.
%\cr &= &2^{-(m-n)}\,\P \{ [d, \d]|B_n  \}
% \cr & & \qquad +{1\over 2^n2^{m-n}}\sum_{k=1}^{m-n} {{m-n}\choose{k}} {1\over
%d\d}\sum_{j=0}^{d\d-1}\sum_{h=0}^n  {{n}\choose{h}}   e^{2i\pi{j\over d\d}  (h^2+kh)},
\end{eqnarray*} 
Note that $\P \{ 
d|B_n\, ,\, \d|B_m \}>0$, only if $m-n\ge (d,\d)$. 
 \vskip 3 pt  
  By using independence,  
  the characteristic function of $B_nB_m$ expresses as follows
$$\E e^{i\upsilon B_nB_m}={1\over 2^n2^{m-n}}\sum_{k=0}^{m-n} {{m-n}\choose {k}}\sum_{h=0}^n   {{n}\choose{h}} e^{i\upsilon (h^2+kh)}.$$
And so, we begin with 
\begin{eqnarray}\label{basicnm}\qq   \P \{ D|B_nB_m \} &=& \E\Big({1\over D} \sum_{j=0}^{D-1} e^{2i\pi  {j \over D}B_nB_m}\Big) 
\cr &=&  {1\over
2^{m-n}}\sum_{k=0}^{m-n}{{m-n}\choose {k}}{1\over D}\sum_{j=0}^{D-1}\Big({1\over
2^n}\sum_{h=0}^n   {{n}\choose{h}} e^{2i\pi{j\over D}  (h^2+kh)}\Big). 
\end{eqnarray}  

 \begin{theorem}\label{p4.5} For any $\e>0$, there exists a constant $C_\e$ depending on $\e$ only, such that  for any
 positive  integers
 $n,m, D$ with 
 $m>n$
\begin{equation}\label{4.11} \Big|\P \{ D|B_nB_m \}-  {1\over D  2^{ m-n}} \sum_{k=0}^{m-n} { {{m-n}\choose {k}}     } \varrho_k(D) 
\Big|
\le C_\e
\left({     D  ^{1 + \e}\over   n }\right)^{1/2 }. 
\end{equation} 
\end{theorem}
  The function    $\varrho_k(D)$ will be  made  explicit    in Theorem \ref{romult1}. A combination of   both results   will lead to  a useful bound of $\P \{ D|B_nB_m
\}$  stated in   Theorem \ref{bnmbound}. 
\vskip 2 pt
 We first prove Theorem \ref{p4.5}.  The main idea of the proof consists with   controlling the difference
$$\D=  {1\over
2^n }  \sum_{h=0}^n   {{n}\choose{h}} e^{2i\pi{j\over D}  (h^2+kh)}- {1\over n}\sum_{h=0}^{n-1}  e^{2i\pi{j\over D} 
(h^2+kh)},$$
 in which the first average   is the  last Euler sum   in (\ref{basicnm}), and  the second one,  its Ces\`aro means counterpart.  Let 
  $0\le
\ell
\le n$ be some integer,   
$\ell=LD+\m$,
$0\le \m< D$. Obviously,
\begin{eqnarray}\label{sumS} S  := {1\over \ell}\sum_{h=0}^{\ell-1}  e^{2i\pi{j\over D} 
(h^2+kh)}  & =& {1\over D}\sum_{h= 0}^{D-1}e^{2i\pi{j\over D} 
(h^2+kh)}  +   {1\over \ell}\sum_{h= 0}^{\m}e^{2i\pi{j\over D}  
(h^2+kh)} 
\cr & &  +L\big({1\over \ell}-{1\over LD}\big)\sum_{h= 0}^{D-1}e^{2i\pi{j\over D} 
(h^2+kh)}.
 \end{eqnarray}
 
 
It follows from (\ref{sumS}),    that 
  for   $j=1,\ldots, D$, 
\begin{eqnarray*} \Big|S-{1\over D}\sum_{h= 0}^{D-1}e^{2i\pi{j\over D} 
(h^2+kh)} \Big|  &\le  &  {1\over  \ell} \Big(\big|\sum_{h= 0}^{D-1}e^{2i\pi{j\over D} 
(h^2+kh)} \big|  +  \big|\sum_{h= 0}^{\m}e^{2i\pi{j\over D} 
(h^2+kh)}\big|\Big).
\end{eqnarray*} 
   We   rewrite this for   later use under the following form.
  \begin{lemma} \label{basicnm2}Let
$\ell\ge 1$, $\ell\equiv \m \, {\rm mod}(D)$ and
$0\le \m< D$.  For   $j=1,\ldots, D$, 
\begin{eqnarray*} \Big| \sum_{h=0}^{\ell-1}  e^{2i\pi{j\over D} 
(h^2+kh)}-{\ell\over D}\sum_{h= 0}^{D-1}e^{2i\pi{j\over D} 
(h^2+kh)} \Big|   
  & \le   &  \big|\sum_{h= 0}^{D-1}e^{2i\pi{j\over D} 
(h^2+kh)} \big|    +  \big|\sum_{h= 0}^{\m}e^{2i\pi{j\over D} 
(h^2+kh)}\big| .
\end{eqnarray*}  \end{lemma}
  Thus the study of $\D$ reduces to the one of the difference
$$ \D'= {1\over
2^n }  \sum_{h=0}^n   {{n}\choose{h}} e^{2i\pi{j\over D}  (h^2+kh)}- {1\over D}\sum_{h=0}^{D-1}e^{2i\pi{j\over D} 
(h^2+kh)} .
$$ 

 By  linearity of Euler summation, it suffices to have at disposal a bound for Euler sums, and we have a  classical tool from the theory of divergent series.  
\begin{lemma}  \label{Eulerbound}Let
$\{a_k, k\ge 0\}$ be a  sequence of reals. Put $A_\ell= \sum_{k=0}^\ell a_k$,  $\ell\ge 0$.
 There exists an   absolute constant $C$ such that for every positive integer $n$
\begin{equation} \label{4.6}\big| \sum_{h=0}^n   2^{-n}{{n}\choose{h}} a_h\big|\le {C\over \sqrt n}\max_{\ell=0}^n\big|A_\ell\big|
.\end{equation}
   \end{lemma}
  This is easily seen by using Abel summation and well-known property of the binomial distribution (Feller \cite{F}, p.\,151). (A  consequence is   Theorem 149 in Hardy \cite{H}.)
\vskip 4 pt
Applying   it  with the choice $$a_h=  e^{2i\pi{j\over D} (h^2+kh)}-
{1\over D}\sum_{y= 0}^{D-1}e^{2i\pi{j\over D}  (y^2+ky)}  \qq h=0,1,\ldots,n,$$
  we get in view of   Lemma \ref{basicnm2},
 \begin{eqnarray}
\D'&\le & { C\over \sqrt n}\max_{\m=0}^{D-1}  \big|\sum_{h= 0}^{\m}e^{2i\pi{j\over D} 
(h^2+kh)}\big|  . 
\end{eqnarray}
 
 Whence the   intermediate estimate.
 \begin{corollary} For some absolute constant $C$, 
 \begin{align}  \label{4.81}   \Big|\sum_{h=0}^n   2^{-n}{{n}\choose{h}} e^{2i\pi{j\over D} (h^2+kh)}-& {1\over D}\sum_{h=
0}^{D-1}e^{2i\pi{j\over D}  (h^2+kh)} \Big|
\cr &\le  {C\over \sqrt n}\bigg(\max_{\m=0}^{D'-1} \big|\sum_{h=0}^{\m} e^{2i\pi{j'\over D'} 
(h^2+kh)} 
\big| \bigg), 
\end{align} where  
$j'=  j/(j,D)
$,
$D'=D/(j,D)$.
 \end{corollary}
 The Lemma below  is just a slight reformulation of Sark\"osy's estimate \cite{Sa}. 
\begin{lemma}  Let  $\a$ be a real number and    $a,q$ be positive integers such that $(a,q)=1$ and
$|\a-a/q|<1/q^2$. Then, for any positive integer  $M$, 
\begin{eqnarray}\label{sarko}
\sup_{k\ge 0}\,\big|\sum_{x=1}^M e^{2i\pi
\a (x^2 +kx)}\big|^2&\le &\sum_{u=1-M}^{M-1}\Big|\sum_{y=\max(1-u,1) }^{\min(M,M-u)}e^{4i\pi \a uy}\Big|
\cr &\le &49
\Big\{{M^2\over  q}+(M\log q) + q\log q \Big\}.
\end{eqnarray} 
\end{lemma}

%It implies for instance by taking $\a=1/N$, $a=1$, $q=N$, $M=m$ that
%$$ \Big|\sum_{k=1}^{m} e^{2i\pi \frac{k^2}{N}  } \Big| \le  7
%\Big\{{m \over  \sqrt N}+\sqrt{m\log N} + \sqrt{N\log N} \Big\},$$
%which for $m\le N$ is comparable to Lehmer's bound, up to a $\sqrt{\log N}$ factor. 
\begin{proof}  The last inequality in \eqref{sarko} is precisely what is established in the proof of Lemma 4 p.\,128 in \cite{Sa}.
Let $T=\sum_{x=1}^M e^{2i\pi
\a (x^2 +kx)}$. On expanding $|T|^2$, we get
\begin{eqnarray*}|T|^2  &=&
\sum_{y=1}^M\sum_{u=1-y}^{M-y} e^{2i\pi \a u(u+2y+  k )} =\sum_{u=1-M}^{M-1}\sum_{y=\max(1-u,1) }^{\min(M,M-u)}  e^{2i\pi \a u(u+2y+  k )}
\cr &\le &
\sum_{u=1-M}^{M-1}\Big|\sum_{y=\max(1-u,1) }^{\min(M,M-u)}  e^{4i\pi \a   uy }\Big|.
\end{eqnarray*}So that,
$$\sup_{k\ge 0}\Big|\sum_{x=1}^M\sum_{y=1}^M e^{2i\pi \a (x^2-y^2 +k(x-y))}\Big|^2\le 49
\Big\{{M^2\over  q}+(M\log q) + q\log q \Big\},$$
or
$$\sup_{k\ge 0}\Big|\sum_{x=1}^M e^{2i\pi
\a (x^2 +kx)}\Big| \le 7
\Big\{{M \over  \sqrt q}+\sqrt{ M\log q} + \sqrt{q\log q} \Big\}.$$
\end{proof}

 

 We can now pass to the proof of Theorem \ref{4.11}.
 \begin{proof}[Proof of Theorem \ref{4.11}]
Returning to \eqref{4.81}, we get  
\begin{eqnarray*}\sup_{k\ge 0}\big|\sum_{y=1}^{\m} e^{2i\pi{j'\over D'} 
(y^2+ky)} 
\big|&\le & 7\Big\{{\m \over  \sqrt{ D'}}+\sqrt{ \m\log D'} + \sqrt{D'\log D'} \Big\}
 \cr (\m \le
D')\quad & \le & 7  \Big\{ \sqrt{ D'}+2\sqrt{ D'\log D'} \Big\}
\cr   & \le& C_\e  (D')^{1/2+ \e}\le C_\e  (D )^{1/2+ \e}, 
\end{eqnarray*}
 where  
$j'=  j/(j,D)
$,
$D'=D/(j,D)$.
Thus
\begin{align}  \label{4.81a}   \Big|\sum_{h=0}^n   2^{-n}{{n}\choose{h}} e^{2i\pi{j\over D} (h^2+kh)}-& {1\over D}\sum_{h=
0}^{D-1}e^{2i\pi{j\over D}  (h^2+kh)} \Big|
\cr &\le  {C\over \sqrt n}\bigg(2+\max_{\m=1}^{D'-1} \big|\sum_{h=1}^{\m} e^{2i\pi{j'\over D'} 
(h^2+kh)} 
\big| \bigg)
  \le  {C_\e  (D )^{1/2+ \e}\over \sqrt n}.
\end{align} 
 
 Whence, 
\begin{equation}\label{4.9} \sup_{ 0\le j< D }\sup_{0\le k\le m-n}\Big|\sum_{h=0}^n   2^{-n}{{n}\choose{h}} e^{2i\pi{j\over D} (h^2+kh)}- {1\over D}\sum_{h=
0}^{D-1}e^{ i\pi{2j\over D}  (h^2+kh)} \Big|\le C_\e \left({     D  ^{1 + \e}\over   n }\right)^{1/2 } \! \!.
\end{equation}  
As
\begin{equation}\label{rho.} {1\over  D^2}\sum_{j=0}^{D-1} \sum_{h= 0}^{D-1}e^{ i\pi{2j\over D} 
(h^2+kh)}={\#\big\{1\le y\le D: D|y^2+ky\big\}\over D }
={\varrho_k(D)\over D } \!, 
\end{equation} 
it follows  
\begin{equation}\label{4.10}  \sup_{0\le k\le m-n} \bigg|{1\over D}\sum_{   j=0}^{D-1}\sum_{h=0}^n   2^{-n}{{n}\choose{h}} e^{2i\pi{j\over D} (h^2+kh)}
   -{\varrho_k(D)\over D } \Bigg|  \le C_\e \Big({     D  ^{1 + \e}\over   n }\Big)^{1/2 }\!\!
.
\end{equation} 
By combining now (\ref{4.10}) with (\ref{basicnm}), we obtain
$$ \Big|\P \{ D|B_nB_m \}-   \sum_{k=0}^{m-n} { {{m-n}\choose {k}} \over 2^{ (m-n)}  } {\varrho_k(D)\over D } 
\Big|
\le C_\e\,
\Big({     D  ^{1 + \e}\over   n }\Big)^{1/2 }. 
$$
\end{proof} 
 
     
   
  
  
     \medskip \par
 
 
 It remains to compute   $\varrho_k(D)$.
  Introduce some   notation. Let   $v_p(d)$ denote the $p$-adic valuation of $d$, $p$ being a prime number, namely $p^{v_p(d)} || d$,   also \begin{eqnarray}\label{D.demi}D_{1/2} \,:=\,\prod_{p|D} p^{\lfloor {v_p(D)\over 2}\rfloor }. 
\end{eqnarray}

   %We shall prove the following exact formula. 
\begin{theorem} \label{romult1}   We have  
  \begin{equation*} \varrho_k(D)=\begin{cases}  \  \,
%\prod_{p|D} p^{\lfloor {v_p(D)\over 2}\rfloor }} 
D_{1/2}\qq & \hbox{ if \ $k=0$,}\cr
%\displaystyle{   \prod_{ v_p(k)<{v_p(D)/2}}(2p^{v_p(k)})\cdot\prod_{   v_p(k)\ge {v_p(D)/2} }  p^{\lfloor{v_p(D)\over2}\rfloor}}\qq &\hbox{ if $k\ge 1$.}
\  \, 2^{\#\{p\,:\, 1\le v_p(k)<v_p(D)/2 \}}\, (k, D_{1/2})\qq &\hbox{ if \ $k\ge 1$.}
\end{cases}
\end{equation*}
     \end{theorem}
  
The proof is given in Appendix. 
%Note  that $ (k, D_{1/2})= \sum_{d|k\,|\,d|D_{1/2}}\p(d)$. 
We deduce
\begin{corollary}\label{rokest} For any positive integers  $D,k$,   
$$\varrho_k(D)\le 2^{\o ((k, D_{1/2}))}\, (k, D_{1/2})
%2^{\o(D) }(k\wedge  \sqrt D) 
 \qq
  {\rm and}\qq
\varrho_0(D)\le  
\sqrt D.
$$   
\end{corollary}  
%[[]]Further, by Corollary \ref{rokest} applied with $D=d\d$,
%\begin{eqnarray*}{1\over  2^{m-n}}\sum_{k=1}^{m-n}C^k_{m-n}{\varrho_k(d\d)\over d\d
%}&\le& {1\over  2^{m-n} }\sum_{k=1}^{m-n} C^k_{m-n} 2^{\o ((k, (d\d)_{1/2}))}\, (k, (d\d)_{1/2}) \cr 
%&= &\E  2^{\o ((B_{m-n}, (d\d)_{1/2}))}\, (B_{m-n}, (d\d)_{1/2}) .
%\end{eqnarray*}

\begin{proof}Immediate. 
\end{proof}
\begin{theorem} \label{bnmbound} We have for any positive integers $n,m,D$,  
\begin{eqnarray*}   \P \{ D|B_nB_m \} 
 &\le &{1  \over D }\sum_{d| D_{1/2} } 2^{\o (d)}d
 \Big(\sum_{u|  (D_{1/2}/d )}\m(u)\, \P\{ud|B_{m-n}\}\Big)
\cr &  &   +{1\over 2^{  m-n }\sqrt D } +C_\e
\left({     D  ^{1 + \e}\over   n }\right)^{1/2 }. 
\end{eqnarray*} 
%$$  \P \{ D|B_nB_m \}    \le   {1\over 2^{  m-n }\sqrt D}+    {2^{\o(D)}( m-n)  \over D } + C_\e {     D  ^{{1/ 2} + \e}\over   \sqrt n }  .    $$
   \end{theorem}  
\begin{proof} By Theorem \ref{p4.5},
 for any $\e>0$, there exists a constant $C_\e$ depending on $\e$ only, such that  for any
 positive  integers
 $n,m, D$ with 
 $m>n$
\begin{equation}\label{4.11a} \Big|\P \{ D|B_nB_m \}-  {1\over D  2^{ m-n}} \sum_{k=0}^{m-n} { {{m-n}\choose {k}}     } \varrho_k(D) 
\Big|
\le C_\e
\left({     D  ^{1 + \e}\over   n }\right)^{1/2 }. 
\end{equation} 
 By Corollary \ref{rokest}, 
 %and since 
 %$ (k, D_{1/2})= \sum_{d|k\,|\,d|D_{1/2}}\p(d)$,
\begin{eqnarray}\label{323}
& & {1\over D  2^{ m-n}}  \sum_{k=0}^{m-n}   {{m-n}\choose {k}} \,  \varrho_k(D)
\ = \  {\varrho_0(D)\over 2^{ 
m-n }D }+{1\over D  2^{ m-n}} 
\sum_{k=1}^{m-n}   {{m-n}\choose {k}}\,
 \varrho_k(D)  
\cr &  \le  &  {1\over 2^{  m-n }\sqrt D }+  {1  \over D 2^{  m-n }}  \sum_{k=1}^{m-n}   {{m-n}\choose {k}} 2^{\o ((k, D_{1/2}))}\, (k, D_{1/2})
\cr & =  &  {1\over 2^{  m-n }\sqrt D }+ 
 {1  \over D 2^{  m-n }}\sum_{d| D_{1/2}} 2^{\o (d)}d
 \sum_{k=1\atop (k, D_{1/2})=d}^{m-n}   {{m-n}\choose {k}}  
%\cr &  =  &  {1\over 2^{  m-n }\sqrt D }+  {1  \over D 2^{  m-n }} \sum_{ d|D_{1/2}} \, \p(d)\sum_{k=1\atop d|k}^{m-n}   C^k_{m-n} 2^{\o ((k, D_{1/2}))}%\cr &  \le  &  {1\over 2^{  m-n }\sqrt D }+  {2^{\o(D)}  \over D 2^{  m-n }}  \sum_{k=1}^{m-n}   C^k_{m-n}k
 %\cr &  \le   &{1\over 2^{  m-n }\sqrt D}+   .... 
 .\end{eqnarray}
Let $\d$ be the arithmetical function defined by
\begin{eqnarray} \label{Delta}   \d(n)=\begin{cases}1\quad & {\rm if}\ n=1,\cr
0\quad & {\rm if}\ n \not=1. \end{cases} \end{eqnarray} Recall that 
\begin{eqnarray} \label{sp} \sum_{d|n}\m(d)=\d(n) . \end{eqnarray} 
 
 Consider first the sub-sum corresponding to $d=1$. We have  
\begin{eqnarray*}
 {1  \over  2^{  m-n }}
 \sum_{k=1\atop (k, D_{1/2})=1}^{m-n}   {{m-n}\choose {k}}
 &=&
 {1  \over   2^{  m-n }}
 \sum_{k=1}^{m-n}   {{m-n}\choose {k}}\d\big((k, D_{1/2})\big)
 \cr&=& 
 {1  \over D 2^{  m-n }}
 \sum_{k=1}^{m-n}   {{m-n}\choose {k}}\Big(\sum_{u|(k, D_{1/2} )}\m(u)\Big)
\cr&=& 
 \sum_{u|  D_{1/2} }\m(u)
 {1  \over   2^{  m-n }} \sum_{k=1\atop u|k}^{m-n}   {{m-n}\choose {k}}\cr&=& 
 \sum_{u|  D_{1/2} }\m(u)\,\big(\P\{u|B_{m-n}\}-{1  \over  2^{  m-n }}
\big)\end{eqnarray*}
Now let $(k, D_{1/2})=d>1$. Let $k=\k d$.
%, $D_{1/2}=\d d$
Then $(k,D_{1/2}/d)=1$, and so
\begin{eqnarray*}
& & {1  \over   2^{  m-n }}
 \sum_{k=1\atop (k, D_{1/2})=d}^{m-n}   {{m-n}\choose {k}}
\ = \ 
 {1  \over   2^{  m-n }}
 \sum_{1\le \k\le (m-n)/d\atop 
 (\k,D_{1/2}/d)=1}     {{m-n}\choose {\k d}}
\cr&=&
 {1  \over   2^{  m-n }}
 \sum_{1\le \k\le (m-n)/d}   {{m-n}\choose {\k d}}\d\big((\k, D_{1/2}/d)\big)
 \cr&=&{1  \over   2^{  m-n }}
 \sum_{1\le \k\le (m-n)/d}   {{m-n}\choose {\k d}}\Big(\sum_{u|(\k, D_{1/2}/d) }\m(u)\Big)
 \cr&=&\sum_{u|  D_{1/2}/d }\m(u)
 {1  \over   2^{  m-n }} \sum_{1\le \k\le (m-n)/d\atop u|\k}   {{m-n}\choose {\k d}}
    \cr &\buildrel{(\k = u\tau)}\over{=}& 
 \sum_{u|  D_{1/2} /d}\m(u)
 {1  \over   2^{  m-n }} \sum_{1\le ud\tau \le  m-n}     {{m-n}\choose {ud\tau}}
 \cr &=&
 \sum_{u|  (D_{1/2}/d ) }\m(u)
 {1  \over   2^{  m-n }} \sum_{1\le x \le  m-n\atop 
 ud|x}    {{m-n}\choose {x}}\cr&=& 
  \sum_{u|  (D_{1/2}/d )}\m(u)\,\big(\P\{ud|B_{m-n}\}-{1  \over  2^{  m-n }}
\big)
  \cr&=&
  \sum_{u|  (D_{1/2}/d )}\m(u) \P\{ud|B_{m-n}\}-{\d(D_{1/2}/d ) \over  2^{  m-n }}.\end{eqnarray*} 
Consequently,
 \begin{eqnarray*}
& &{1  \over  2^{  m-n }}\sum_{d| D_{1/2}} 2^{\o (d)}d
 \sum_{k=1\atop (k, D_{1/2})=d}^{m-n}   {{m-n}\choose {k}}  
\cr &=&
 \sum_{d| D_{1/2}} 2^{\o (d)}d
 \Big(\sum_{u|  (D_{1/2}/d )}\m(u)\,\big(\P\{ud|B_{m-n}\}-{\d(D_{1/2}/d )  \over  2^{  m-n }}\Big)
 \cr &=&
 \sum_{d| D_{1/2} } 2^{\o (d)}d
 \Big(\sum_{u|  (D_{1/2}/d )}\m(u)\, \P\{ud|B_{m-n}\}\Big)-{1  \over  2^{  m-n }} .
 \end{eqnarray*}
We observe   that $\d(D_{1/2}/d ) =0$ as long as $d< D_{1/2}$, and if $d=D_{1/2}$, then the corresponding summand   produces the contribution $2^{\o (D_{1/2})}D_{1/2}   \P\{ D_{1/2}|B_{m-n}\}-{1  \over  2^{  m-n }}$  (recalling that $\m(1)=1$), whence the last line. 
\vskip 3 pt
By reporting this in \eqref{323} we get
\begin{eqnarray}\label{323a} 
 {1\over D2^{ m-n }  } \sum_{k=0}^{m-n}  {{m-n}\choose {k}} \, \varrho_k(D) 
 & =  & {1  \over D 2^{  m-n }}\sum_{d| D_{1/2}} 2^{\o (d)}d
 \sum_{k=1\atop (k, D_{1/2})=d}^{m-n}   {{m-n}\choose {k}}  +{1\over 2^{  m-n }\sqrt D }
\cr & =  &  {1  \over D }\sum_{d| D_{1/2} } 2^{\o (d)}d
 \Big(\sum_{u|  (D_{1/2}/d )}\m(u)\, \P\{ud|B_{m-n}\}\Big)
 \cr & &  +{1\over 2^{  m-n }\sqrt D }-{1  \over  D2^{  m-n }} 
   .\end{eqnarray}
By inserting now this into  \eqref{4.11a}, we finally get
\begin{eqnarray}\label{4.11b}  \P \{ D|B_nB_m \} 
&\le &{1  \over D }\sum_{d| D_{1/2} } 2^{\o (d)}d
 \Big(\sum_{u|  (D_{1/2}/d )}\m(u)\, \P\{ud|B_{m-n}\}\Big)
\cr &  &   +{1\over 2^{  m-n }\sqrt D }-{1  \over D 2^{  m-n }}+C_\e
\left({     D  ^{1 + \e}\over   n }\right)^{1/2 }
\cr &\le &{1  \over D }\sum_{d| D_{1/2} } 2^{\o (d)}d
 \Big(\sum_{u|  (D_{1/2}/d )}\m(u)\, \P\{ud|B_{m-n}\}\Big)
\cr &  &   +{1\over 2^{  m-n }\sqrt D } +C_\e
\left({     D  ^{1 + \e}\over   n }\right)^{1/2 }. 
\end{eqnarray} 
This achieves the proof of Theorem \ref{bnmbound}.
\end{proof}

\vskip 30 pt

 
 
 
 
 
 

\section{Divisibility and primality in the Rademacher random walk}\label{s8}
 %%%%%
%%%%%
The primality properties in that case  are more studied than divisibility. The necessary value distribution results  of divisors     are taken from  \cite{W7a}, which contains an application to the functional equation of the
 Zeta-Riemann function. An extremal divisor case is thoroughly investigated in \cite{W6}. The proof deeply relies on mixing properties of the system $\big\{ \chi\{d|R_n\},    \    2\le d\le n,\ n\ge 1\big\}$.
 
 \subsection{Divisibility.}Put, 
\begin{eqnarray}\label{radem(2.7)}\qq \qq  
\Theta_1(\d,M) & =& 2\ \displaystyle{\sum_{\ell\in \Z} e^{-{2M\pi^2\ell^2\over \d^2}},} \cr   
\Theta_2(\d,M)&=&2\ \displaystyle{\sum_{\ell=0}^\infty \big\{e^{-{ M\pi^2(2\ell )^2/ 2\d^2}}-e^{-{
M\pi^2(2\ell+1)^2/ 2\d^2}}\big\}-1.}
 \end{eqnarray}
%\eqno(2.7)
We note that
\begin{equation}\Theta_1(\d,M)=2\,\d /\sqrt{  {2M\pi } } +\mathcal O(1)
\end{equation}
when $\d\to \infty$, uniformly in any interval of   $M$.


 \vskip 3 pt  The  Theorem below provides precise estimates of $\P\big\{ \d | R_{M}\big\}$.
\begin{theorem}[\cite{W7}, Th.\,2.1]\label{radem.t2.1}
% Theorem 2.1.}
 Let $\a >\a^\prime>3/2$. There exist constants
$C$, $\d_0$  and  
$M_0$, depending on $\a,\a^{\prime}$ only, such that for any $M\ge M_0$, $M\ge \d\ge \d_0$:
\vskip 3 pt  
--- if $\d$ and $M$ are even,
\begin{eqnarray} \Big|\P\big\{ \d | R_{M}\big\} -{ 2 + 4 \cos^M {2\pi  \over \d}  \over \d}\Big| \le  M^{- \a^\prime },& \qquad 
\hbox{if} \ \  \d <
2\pi 
\sqrt{M\over {2\a\log M}} ,\cr
 \Big|  \P\big\{ \d | R_{M}\big\}-  {1\over \d}\Theta_1(\d,M)  \Big|\le C  {
\log^{5/2} M 
\over M^{3/2}},& \qquad \hbox{if} \ \d \ge  2\pi  \sqrt{M\over {2\a\log M}}. 
\end{eqnarray}
\vskip 3 pt   
--- if $\d$ and $M$ are odd,
\begin{eqnarray}  \Big|\P\{ \d|R_M\}-{1+ 2 (\cos {2\pi  \over \d})^M-2(\cos{\pi  \over \d})^M \over \d}\Big| 
\le M^{- \a^\prime }, & &  \hbox{if} \ \ \d <
2\pi 
\sqrt{M\over {2\a\log M}} ,\cr
  \Big|\P\{ \d|R_M\} - { \Theta_2(\d,M) \over \d}\Big|\le C   {\log^{5/2} M \over M^{3/2}},  & & \hbox{if} \ \d \ \ge  2\pi  \sqrt{M\over {2\a\log
M}}.
\end{eqnarray} 
\end{theorem}

\bigskip\par 
\subsection{An extremal divisor case.} Put  successively, 
$$ N_1 =1, \qquad N_k =\inf\{\, N>N_{k-1}  \, :\, N\   {\rm even \  and} \  \, N | R_{N^2}\, \},\qquad(k>1). $$ 
The fact that this random sequence  is well defined,  can be deduced from our result below. 
%Before stating it, it is  worth observing that for no
%sequence of positive reals
%$\{a_k, k\ge 1\}$ tending to infinity, the ratios 
%$${N_k\over a_k}$$
%can  converge in probability to a random variable that is positive almost surely. Indeed,
%according to a classical result, this would imply the validity of a randomly selected  central limit theorem, namely
%(see \cite{Bi}, Th.\,17.2 p.\,147) 
%$$ {R_{N_k}\over\sqrt{N_k}} \quad \buildrel{{\mathcal D}}\over{\Longrightarrow}\quad {\mathcal N}(0,1), $$
%which is naturally impossible. 

\vskip 3 pt The following fine asymptotic result    exhibits an exponential growth of the sequence
$(N_k)_k$. 

\begin{theorem}[Weber \cite{W6}]\label{NkRn} Put $s= 2 \sum_{j\in \Z}  e^{-2\pi^2j^2}$. For any $\tau>7/8$,
$$ \     \log N_k = {k\over s} + {\mathcal O}\big(  k^{\tau} \big) \qquad\qquad  \hbox{\it almost surely.} $$
\end{theorem}
 
  The result extends to Bernoulli sequences $\b= \{\b_i , i\ge 1\}$. Write $B_N  = \b_1+\ldots +\b_N$, $N=1,2,\ldots$ Since $
\e_i\buildrel{{\mathcal D}}\over{=}2\b_i-1 $, Theorem \ref{NkRn} gives the same estimate for the sequence 
 $$ M_1 =1, \qquad M_k =\inf\{\, M>M_{k-1}  \ :\ M\  \textstyle{even \  and} \quad M | \,2B_{M^2} \, \},\qquad(k>1) $$
 
The proof is intricate and  long.  We only give a sketch of the main steps. First for $N$ even,  one has the asymptotic estimate
\begin{equation}N\P\big\{ N | R_{N^2}\big\} = s + {\mathcal O}\Big( { \log^{5/2} N  \over N^2}\Big) . 
\end{equation}

  Next  it is   necessary to have at disposal asymptotic estimates of  
  %$\P\big\{ N | R_{N^2},  M | R_{M^2}\big\}$, and more precisely 
  the
correlation 
$$\P \{ N | R_{N^2},  M | R_{M^2} \} - \P \{  N |
R_{N^2} \}\P \{  M | R_{M^2} \}.$$
This is a considerably more   delicate task.
The next
statement is the crucial step towards the proof of Theorem \ref{NkRn}.
\begin{lemma}[\cite{W6}] \label{NkRnl2}
There exists an absolute constant
$C$, and for every $\e >0$, a constant $C_\e$ depending on $\e$ only, such that for any even positive  integers
$N\ge M$ large enough,
\vskip 3pt a) if $\  M\le N\le M+  N/(\log N)^{1/2}$ 
 \begin{align*} \Big|\P\big\{ N | R_{N^2},  M | R_{M^2}\big\}- \P\big\{  N |
R_{N^2}\big\}\P\big\{  M | R_{M^2}\big\}\Big|\le C\,{(\log N)^{1/2}\over N } {(\log M)^{1/2}\over M}
.\end{align*}
\vskip 3pt  b) if $  N>  M+  N/(\log N)^{1/2}$ 
\begin{align*}  \Big|\P\big\{ N | R_{N^2},  M | R_{M^2}\big\}-& \P\big\{  N |
R_{N^2}\big\}\P\big\{  M | R_{M^2}\big\}\Big|
\cr &\le C_\e\,\Big\{  {(\log M)^{5/2}(\log N)^{1/2}\over M^2N}+ {1\over N^2}(\log
N)^{3/4\, \,+\e} \Big\}.
\end{align*}
\end{lemma}

The next tool is Lemma 10 p.\,45  in Sprind\v zuk \cite{Sp}, a typical criterion of almost sure convergence, providing  in addition an estimate of the  remainder term. %This kind of criterion is of course particularly relevant in the study of divisibility in random walks.



\begin{lemma} \label{NkRnl4} 
Let $ \{f_l, l\ge 1\}$ be a sequence
 of non negative random variables, and let $\p=\{\p_l, l\ge 1\}$, $m=\{m_l, l\ge 1\}$ be two sequences of
non negative reals such that 
\begin{equation}\label{RNRM.2.49}
 0\le \p_l\le m_l \le 1\qquad (l\ge 1).
\end{equation}
 %\eqno(2.49)$$
Write 
$$ M_n =\sum_{1\le l\le n}  m_l.$$

\noi  Assume that $M_n\uparrow\infty$ and that the following condition is satisfied 
\begin{equation}\label{RNRM.2.50}\E\  \big| \sum_{i\le l \le j}(f_l-\p_l)\big|^2 \,\le
C\sum_{i\le l \le j} m_l ,
\qquad  0\le i\le
 j<\infty. 
 \end{equation}
 %\eqno(2.50) $$
    Then, for every $a>3/2$,   
\begin{equation}\label{RNRM.2.51} \sum_{1\le l\le n}  f_l \ \buildrel{a.s.}\over{=} \sum_{1\le l\le n}  \p_l +{\mathcal O}\Bigl( M_n^{1/2}  {\log}^{ a}
M_n \Big). 
\end{equation}
%\eqno(2.51)$$ 
 \end{lemma}
\subsection{Primality.}

 
\begin{theorem}[\cite{W5a}]\label{P-Rn} There exist two absolute constants $C_1$ and  $C_2$ such that for any integer $n\ge 3$,
$$ {  C_1  \over    \log
n }\le \P\{R_n\ \hbox{is prime} \}\le { C_2  \over    \log
n }. $$   
In fact, for every $\e>0$,
$${4    \over  \sqrt{ 2e\pi  }\log n}-{C _\e  \over  n^{  
    1/2-\e }}\le \P\{R_n\ \hbox{is prime} \}\le   \sqrt{{2
\over     \pi }  }\Big(\sum_{k=1}^{\infty}
{2^{ k }    }e^{-   2^{2k-1  } }\Big)  {1
\over \log     n  }+{C'_\e  \over  n^{  
    1/2-\e }},$$
where the constants $C_\e$, $C_\e '$ depend on $\e$ only.\end{theorem}
%\vskip 5 pt 
The constants obtained are certainly not the best, and the  limit  $\displaystyle{\lim_{n\to \infty}} (\log n)\,\P\{R_n\, \hbox{prime} \}$ may  exists.

\begin{proof}
Put  
$$I_k(n) =[-2^{ k}\sqrt n, 2^{ k}\sqrt n], \qq k=0,1,\ldots $$
next 
\begin{equation} \label{jkn}
  J_k(n)=  \begin{cases}  {\mathcal P}\cap  I_0(n) & \qq \hbox{if $k=0$},\cr         {\mathcal P}\cap\big(I_k(n) \backslash I_{k-1}(n) \big)  & \qq \hbox{if $k\ge 1 $},\end{cases}
\end{equation} 
and let  $k_n= \max\{ k: 2^{k-1}\le
a \sqrt{ 2\log \log n} \}$, where $a>0$ will be  chosen later on. 
Then 
$$\P\{R_n\ \hbox{ prime} \}=\sum_{k=0}^\infty \P\{R_n\in  J_k(n) \}=\sum_{k=0}^{k_n} \P\{R_n\in  J_k(n) \}+\rho_n . $$
%Recall that   $G(x) =e^{x^2}-1$ and  $(L^G(\P),\|.\|_{G,\P})$ is the associated Orlicz  space. Let  $g$ denote a standard Gaussian  random variable. We have $$\|{R_n\over\sqrt n} \|_{G,\P}
%=\big({2\over \pi}\big)^{1/2}\|{R_n\over \sqrt n} \big({\pi\over 2}\big)^{1/2}\|_{G,\P}\le \big({2\over \pi}\big)^{1/2}\|g\|_{G,\P}= \big({2\over \pi}\big)^{1/2}2 \big({2\over 3}\big)^{1/2}=   {4 \over \sqrt {3\pi} }:=c.$$
There exists a numerical constant $c$ ($c={4 \over \sqrt {3\pi} }$ is suitable) such that if $a>c\sqrt 2$, 
\begin{eqnarray*} \P\{  |R_n|> a\sqrt{2 n\log \log n}\}
&=&\P\{  |{R_n\over \sqrt n }|> a\sqrt{2  \log \log n}\}
\cr  & \le &e^{ -(2a^2/c^2) \log \log
n}\E e^{ {|R_n|^2/(
 c^2 n )}}  
 \cr  & \le & {2\over (\log n)^{2a^2/c^2}} .
 \end{eqnarray*} 
Thus
$$\rho_n\le \P\{|R_n|\ge a\sqrt{2n\log\log n}\} \ll {1\over  \log n  }. $$

Besides, 
$$\P\{R_n\in  J_k(n) \}= \P\{R_n\in  J_k(n) \}- {1\over  \sqrt{ 2\pi n}} \sum_{p\in  J_k(n)}e^{-  {p^2\over  2 n } }+{1\over  \sqrt{ 2\pi n}} \sum_{p\in 
J_k(n)}e^{-  {p^2\over  2 n } }.$$
By means of Theorem \ref{radem.t2.1}, 
$$\sum_{k=0}^{k_n}\Big|\P\{R_n\in  J_k(n) \} - {1\over  \sqrt{ 2\pi n}} \sum_{p\in  J_k(n)}e^{-  {p^2\over  2 n } }\Big|  \le   {C_\e\over n^{  
   3/2-\e }} \sum_{k=0}^{k_n}\#(
J_k(n)) \le C_\e   n^{  
   -1/2+\e }.$$
%$$\Big|\P\{R_n\in  J_0(n) \} - {1\over  \sqrt{ 2\pi n}} \sum_{p\in  J_0(n)}e^{-  {p^2\over  2 n } }\Big|  \le C_\e {\#( J_0(n))\over n^{  
%   3/2-\e }}\le C_\e  {\pi (n) \over n^{  
%   3/2-\e }}\le C_\e   n^{  
%   -1+\e } .$$
 
The following estimates are classical: (\cite{RS}, Th.\,2 and Cor.\,1, p.\,69, and  \cite{DR}, Lemma 5, p.\,439)
\begin{eqnarray*}x/\log x \le \pi(x)\le x/(\log x-3/2), & &\qq \hbox{for any $x\ge 67$,}   \cr 
 \pi(2x)-\pi(x)\le {x\over \log x}, &   & 
 \qq\hbox{for any $x>
1$}.
\end{eqnarray*}
Then $$  {4    \over  \sqrt{ 2e\pi  }\log n}\le{2\pi(\sqrt n)\over  \sqrt{ 2e\pi n}}\le {1\over  \sqrt{ 2\pi n}} \sum_{p\in  J_0(n)}e^{-  {p^2\over  2 n
} }\le   {2\pi(\sqrt n)\over  \sqrt{ 2\pi n}}\le   {4   \over ( \log n-3 ) \sqrt{ 2\pi
 }}.$$
And
\begin{eqnarray*}{1\over  \sqrt{ 2\pi n}}\sum_{k=1}^{k_n}  e^{-  {p^2\over  2 n } }\  \le \ {2\over  \sqrt{ 2\pi n}}\sum_{k=1}^{k_n} {2^{ k  }\sqrt
n\over \log (2^{ k } \sqrt n)}e^{-   2^{2k-1  } } \  \le \  \sqrt{{2
\over     \pi }  }\Big(\sum_{k=1}^{\infty}
{2^{ k }    }e^{-   2^{2k-1  } }\Big)  {1
\over \log     n  }
 .
 \end{eqnarray*}
We deduce
$$ {4    \over  \sqrt{ 2e\pi  }\log n}\le{1\over  \sqrt{ 2\pi n}}\sum_{k=0}^{k_n}  e^{-  {p^2\over  2 n } }\le  \sqrt{{2
\over     \pi }  }\Big(\sum_{k=1}^{\infty}
{2^{ k }    }e^{-   2^{2k-1  } }\Big)  {1
\over \log     n  }.$$
From there, it follows immediately that
$${4    \over  \sqrt{ 2e\pi  }\log n}-{C _\e  \over  n^{  
    1/2-\e }}\le \P\{R_n\ \hbox{ prime} \}\le   \sqrt{{2
\over     \pi }  }\Big(\sum_{k=1}^{\infty}
{2^{ k }    }e^{-   2^{2k-1  } }\Big)  {1
\over \log     n  }+{C'_\e  \over  n^{  
    1/2-\e }}.$$
Whence the claimed result.
\end{proof}
 %The study of $\P\{B_n\, \hbox{prime} \}$ is much more complicated. A bound of the type $\P\{B_n\ \hbox{ prime} \}\ge {C\over \log n}$ would imply


\bigskip\par
\bigskip\par








%Other results are also proved in \cite{W5a}. 
The following Theorem is complementing the above.
 

\begin{theorem}[\cite{W5a}]\label{P-RnRm} Let $m>n$ be two odd positive integers. Then for some absolute constant $C$, and every $\e>0$,
$$\P\{R_n\ {\rm and } \ R_m\ \hbox{are prime} \}\le  {C\over \sqrt{ m-n } \,\log n} +C_\e n^{-1/2+\e} + {C \log\log^+    (m-n  )\over (\log
n)^2}     .$$\end{theorem} 
 
  
\begin{proof}  Let  $m>n$ be two odd positive integers. We have
\begin{eqnarray}\label{P-RnRm.1}\P\{R_n\ {\rm and} \ R_m\hbox{ prime} \}&=&\P\{R_n\ {\rm and} \ R_n +(R_m-R_n)\ \hbox{prime} \}
\cr &=& \sum_{|p|\le  n}\P\{R_n=p \}\Big(\sum_{|v|\le m-n\atop
p+v\in {\mathcal P}} \P\{R_{m-n}=v \}\Big)
\cr &=&\sum_{|v|\le m-n}\Big(\sum_{|p|\le  n\atop
p+v\in {\mathcal P}}\P\{R_n=p \}\Big)\P\{R_{m-n}=v \}.
\end{eqnarray} 
%Since $m-n$ is even, only the case $v$ even may contribute. 
The summation term corresponding to $v=0$ equals to
\begin{eqnarray}\label{v=0} \sum_{|p|\le  n\atop
p \in {\mathcal P}}\P\{R_n=p \} \P\{R_{m-n}=0 \}& =&\P\{R_n  \, \hbox{prime}\} \P\{R_{m-n}=0 \} 
\cr &\le& {C\over \sqrt{ m-n } }\,\P\{R_n  \, \hbox{prime}\}   \,\le \,  {C\over \sqrt{ m-n } \,\log n}
\end{eqnarray} 
since $\P\{R_{m-n}=0 \}=2^{-(m-n)}{{m-n}\choose{{ m-n\over2}}}\sim{2\over \sqrt{2\pi(m-n)}}$.
\vskip 2 pt We next write for $|v|\le m-n$,
$$ \sum_{|p|\le  n\atop
p+v\in {\mathcal P}}\Big|\P\{R_n=p \}-{1\over  \sqrt{ 2\pi n}}   e^{-  {p^2\over  2 n } } \Big|\le C_\e n^{1-3/2+\e}= C_\e n^{-1/2+\e}.$$
Further,
$$  {1\over  \sqrt{ 2\pi n}} \sum_{|p|\le  n\atop
p+v\in {\mathcal P}}  e^{-  {p^2\over  2 n } } = {1\over  \sqrt{ 2\pi n}} \sum_{k=0}^\infty \sum_{p\in J_k(n)\atop
p+v\in {\mathcal P}} e^{-  {p^2\over  2 n } },$$
where $J_k(n)$ are defined in \eqref{jkn}. Recall  Siebert's estimate \cite{Si}: let $a\neq0$, $b\neq 0$ be integers with $(a,b)=1$, $2|ab$. Then for $x>1$,
 \begin{equation}\label{Siebert.est.}
\sum_{p\le x \atop   ap+b\in {\mathcal P} }\le 16\o{x \over (\log x)^2}  \prod_{p|ab\atop p>2}\, \Big({p-1\over p-2}\Big), \qq \o= \prod_{p>2}\Big(1-\frac{1}{(p-1)^2}\Big) .
\end{equation}
   Daboussi and   Rivat's observation  \cite{DR}, p.\,440, that
$$\frac{p-1}{p-2}= \Big(1-\frac{1}{p}\Big)^{-1}  \Big(1-\frac{1}{(p-1)^2}\Big)^{-1} $$   
and when $h$ is even
$$\prod_{p|h\atop p\ge 3}\Big(1-\frac{1}{p}\Big) = 2\frac{\p(h)}{h}, 
$$ 
%where   $\p $  is Euler's function, 
imply 
that for $x>1$ and  any positive even number $h$,
\begin{equation}\label{Siebert.est}
\sum_{p\le x \atop    p+h\in {\mathcal P} }1\le 8\prod_{p\atop (p,h)=1}\Big(1-\frac{1}{(p-1)^2}\Big){ x\over (\log x)^2} {h\over \p(h)} \le 8\,{ x\over (\log x)^2} {h\over \p(h)} .\end{equation}

In    \eqref{P-RnRm.1}, the  sums are indexed over   pairs $(p,p+v)$ with $|p|\le n$, $|v|\le m-n$.   By symmetry, %To treat the case when they have different signs, 
it     suffices to restrict to the case $p>0$, $v<0$, namely $(p,p-|v|)$.  
%The case $v =-2p$ provides the pair $(p,-p)$.
%, as there is no problem in applying   \eqref{Siebert.est} if $p$ and $v$ have same sign. 
\vskip 4 pt Consider three sub-cases: (a) $|v|< p$, (b) $|v|\ge 2p$, (c) $p\le  |v|<2p$.  
\begin{eqnarray*}
{\rm (a)}& &
%\hbox{Let $q=p-|v|$. Then $(q,q+|v|)$ satisfies} 
\hbox{$(p-|v|,p)$ with  $1\le p-|v|\le n$, $|v|\le m-n$.}
\cr {\rm (b)} & & \hbox{%$|v|-p=p+h$ with   
$(p, |v|- p)=(p,p+|v|-2p)$  with $1\le p\le n$, ${0\le |v|-2p\le m-n}.^{(*)}$}
   \cr {\rm (c)} & &
    %\hbox{Let $q=|v|-p$, $h=2p-|v|$.  Then $(q,q+h)$ satisfies}
    \hbox{$(|v|-p, p)=(|v|-p,|v|-p+2p-|v|)$ with $1\le |v|-p\le 2p-p=p\le n$,}     \cr & & \hbox{
      $1\le 2p-|v|\le 2|v|-|v|=|v|\le m-n$.}
     %  $0<|v|-p\le p\le |v|$.
\end{eqnarray*}
 ${}^*$(The case $2p-|v|=0$ in (b) is treated in  \eqref{v=0}.) %\vskip 2pt 
%(a)\quad We write $q=p-|v|\le n$, $q+|v|=p$. Let $(q,q+|v|)$ and notice that $1\le q\le n$, $|v|\le m-n$.
%\vskip 2pt
%(b)\quad Write   $|v|-p=p+h$ with $h=|v|-2p$. Let 
%$(p,p+h)$ and note that $p\le n$, $0<h\le m-n$.
%\vskip 2pt
%(c)\quad We have $0<|v|-p\le p\le |v|$. Let $(q,q+h)$ with $q=|v|-p$, $h=2p-|v|$, and note that $h\le 2|v|-|v|=|v|\le m-n$, $q\le 2p-p=p\le n$. 
%\vskip 2 pt
  Thus  the pairs $(p-|v|,p)$ are pairs $(\a, \a+h)$ of positive integers   all satisfying  $\a\le n$, $h\le m-n$. So that estimate \eqref{Siebert.est} remains applicable to this case, up to a multiplicative factor 4.%D'autre part
%$$  {|v|\over \p(|v|)}\le e^\gamma \log\log |v| + {2,50637\over \log\log |v|}\le C\log\log |v|.$$
%D'o\`u,
%$$\#\{ p\le x : p+v\in {\cal P} \}\le C{ x\over (\log x)^2} \log\log |v|.$$
\vskip 2pt

It follows that  $$\#\{ p\in J_k(n) : p+v\in {\mathcal P} \}\le \#\{ |p|\le 2^k\sqrt n : p+v\in {\mathcal P} \}\le C{ 2^k\sqrt n\over (\log 2^k\sqrt n)^2}
{|v|\over \p(|v|)}. $$ 
Hence,
$${1\over  \sqrt{ 2\pi n}} \sum_{k=0}^\infty \sum_{p\in J_k(n)\atop
p+v\in {\mathcal P}} e^{-  {p^2\over  2 n } }\le {C\over (\log n)^2}{|v|\over \p(|v|)} \sum_{k=0}^\infty 2^k e^{-2^{2k-3}}\le {C\over (\log
n)^2}{|v|\over \p(|v|)}.$$
 Consequently
$$\sum_{|p|\le  n\atop
p+v\in {\mathcal P}} \P\{R_n=p \}\le C_\e n^{-1/2+\e} + {C\over (\log
n)^2}{|v|\over \p(|v|)}.$$
By reporting,
$$\P\{R_n\ {\rm and} \ R_m\hbox{ prime} \}  =\sum_{|v|\le m-n}\Big(\sum_{|p|\le  n\atop
p+v\in {\mathcal P}}\P\{R_n=p \}\Big)\P\{R_{m-n}=v \} $$
$$ \le {C\over \sqrt{ m-n } \,\log n} +C_\e n^{-1/2+\e} + {C\over (\log
n)^2}\sum_{0<|v|\le m-n} {|v|\over \p(|v|)}\P\{R_{m-n}=v \} .$$
$$ \le {C\over \sqrt{ m-n } \,\log n} +C_\e n^{-1/2+\e} + {C\over (\log
n)^2}\ \E \, {|R_{m-n}|\over \p(|R_{m-n}|)} {\bf 1}_{\{|R_{m-n}|\ge 1\}} .$$
We have obtained 
\begin{eqnarray*}\P\{R_n\ {\rm and} \ R_m\hbox{ prime} \}&\le & {C\over \sqrt{ m-n } \,\log n} +C_\e n^{-1/2+\e} 
\cr & &\qq \quad + {C\over (\log
n)^2}\ \E \, {|R_{m-n}|\over \p(|R_{m-n}|)} {\bf 1}_{\{|R_{m-n}|\ge 1\}}.
\end{eqnarray*}
We have the following estimates concerning $\p$ (Rosser and Schoenberg \cite{RS},  inequality 3.40 and 3.41), for $h\ge 3$
\begin{eqnarray*}&& \sum_{k=1}^n {k\over \p(k)}={\mathcal O}(n)  ,
\cr & & {|v|\over \p(|v|)}\le e^\gamma \log\log |v| + {2,50637\over \log\log |v|}\le C\log\log |v|.
 \end{eqnarray*}
We deduce
$$\P\{R_n\ {\rm and} \ R_m\hbox{ prime} \}\le  {C\over \sqrt{ m-n } \,\log n} +C_\e n^{-1/2+\e} + {C \log\log^+    (m-n  )\over (\log
n)^2}     .$$
\end{proof}










  



\bigskip\par
 
\subsection{Primality mod($k$).}

Theorem \ref{P-Rn} can be extended to a   larger class of random variables under moderate integrability assumption (only the upper bound part is concerned).  Indeed, by using the Brun-Titchmarsh  theorem and the local limit theorem, a similar result is valid under a stronger form,   the constraint made on $k$ below, is weak but the bound obtained, although sharp, should admit  improvments (maybe the $\log\log$ term can be dropped).

\begin{theorem}[\cite{W5a}]\label{P-Sn} Let $X$ be a symmetric random variable such that   
 $\E X =0=\E X ^3$,  $\E X ^2=1$ and $\E |X| ^{3+\e }<\infty$  for some $\e >0$. Let $X_i$, $i\ge 1$ be independent copies of $X$ and for each positive integer $n$, $S_n= X_1+\ldots + X_n$. Let $k$ be some positive integer and  $\ell$ be an  odd integer such that
   $(\ell,k)=1$ and $
\sqrt{ n  \log\log n}/k\to
\infty$. Then for  every $B>1$,   
$$ \P\Big\{S_n \  \hbox{is prime},\,  S_n\equiv \ell \ ({\rm mod}\, k)\Big\}={\mathcal O}_B\Big\{   {\sqrt{   
\log\log n}\over \varphi(k)\log    n   } +  {  1\over (\log n)^{B} }\Big\}.
$$ 
\end{theorem}  








The proof uses the following local limit theorem due to Ibragimov, which characterizes the speed of convergence in the i.i.d. case. 
 
\begin{proposition}[\cite{P},  p.\,216]\label{llt.indep.ibra.g.moments} Let  $\{X_k, k\ge 1\}$ be a sequence of centered i.i.d. $\Z$-valued random variables, with common distribution function $F$, and unit variance. 
 Put $$r_n= \sup_{k}\Big|  \sqrt n \P\{S_n=k\}-{1\over  \sqrt{ 2\pi}}e^{-  {k^2\over  2 n } }\Big|  $$

Then   $ r_n ={\mathcal O}\big(n^{-(m+\d)/2}\big) $ 
for some integer  $m\ge 0$  and a real $\d$,  $\d\in ]0,1[$, if and only if  $$\sup\big\{h>0; \exists a\in
\Z:
\P\{X_1\in a+h\Z\}=1\big\}= 1,$$  and the two following conditions are realized:
\smallskip\par {\rm (i)} All moments of  $X_1$ up to the order  $m+2$ included exist and coincide with those of the standard normal law.  
\smallskip\par {\rm (ii)}
$$\int_{|x|\ge v}|x|^{m+2} F(dx)={\mathcal O}\big(v^{-\d}\big) \qq  \qq v\to \infty. $$
\end{proposition}
\begin{proof}[Proof of Theorem \ref{P-Sn}]
%The moments of $g$ satisfy
  % $ \E |g|^{2\ell}= {(2\ell)!\over 2^\ell \ell!} $,     $ \E |g|^{2\ell+1}=0 $,   
 %$\ell=0,1, \ldots$, it follows that 
% if  $X=\e$, Theorem \ref{llt.indep.ibra.g.moments} does apply with the only possible choice  $m=1$,   $\d\approx1$, since   $\E \e^4=1 $, whereas   $\E g^4=3$. And we have,
%\begin{equation}\label{}  \sup_{k}\Big|  \sqrt n \P\{R_n=k\}-{1\over  \sqrt{ 2\pi}}e^{-  {k^2\over  2 n } }\Big|   ={\mathcal O}_\e\big(n^{-{ 
%  1+\e}}\big).
 % \end{equation}
  %\eqno(3.1) 
% Now let $\{X_k, k\ge 1\}$ be  a sequence   of independent  integer-valued centered random variables such that $\E X  =\E X ^3=0$ and $\E X ^2=1$, and for some $0<\d<1$,
%$$\int_{|x|\ge v}|x|^{ 2} F(dx)={\mathcal O}\big(v^{-\d}\big) \qq {\rm as} \qq v\to \infty  .$$
 
 
  As there exists   $\e>0$ such that $\E |X|^{3+\e}<\infty $, $X$ satisfies the requirements of  Proposition \ref{llt.indep.ibra.g.moments} with 
$m=1$ and
 $0<\d<1$, $\d= \d(\e)$. Then
\begin{align*}
%\label{P-Sn(3.3)} 
\sup_{k}\big|  \sqrt n \P\{S_n=k\}-{1\over  \sqrt{ 2\pi}}e^{-  {k^2\over  2 n } }\big| ={\mathcal O}\big(n^{-(1+\d)/2}\big),\end{align*}
%\eqno(3.2)  $$
and
\begin{align}\label{P-Sn(3.3).} \big| \sum_{k\in J}\P\{S_n=k\}-{1\over  \sqrt{ 2\pi n}}\sum_{k\in J}
e^{-  {k^2\over  2 n } }\big|\le C  {\#\{J\}\over   n^{1+(\d/2)}}.
\end{align}
%\eqno(3.3)  $$

Besides, by means of a result of  Robbins and Siegmund (\cite{P}, p.\,314), if $\a>1$, $b>1$ are two arbitrary but fixed reals,
there exists a constant $C$ depending on   $\a$, $b$ and $X$, an integer  $n(\a)$ depending on  $\a$ only, such that for any $n\ge
n(\a)$,
\begin{align}\label{P-Sn(3.4)} \P\Big\{ \exists m\ge n : S_m>\sqrt{2\a(1+b)m  \log\log m}\Big\}\le C (\log n)^{-b} (\log\log n)^{-1/2}.
\end{align}
% \eqno(3.4) $${Sh.We1}visiblf
As  $X$ is symmetric, $\P\{|S_n|>u\}\le 2\P\{ S_n >u\}$; and thus it follows from  \eqref{P-Sn(3.4)} that 
\begin{align}\label{P-Sn(3.5)} \P\Big\{  |S_n|>\sqrt{2\a(1+b)n  \log\log n}\Big\}\le 2C (\log n)^{-b} (\log\log n)^{-1/2}.\end{align}
% \eqno(3.5) $$
 Let  $k$, $\ell$
be coprime integers. Put
$$J=\big\{2\le |p|\le \sqrt{2\a(1+b)n  \log\log n}: p\equiv \ell \ ({\rm mod}\,
k)\big\} .$$  
Using \eqref{P-Sn(3.3).} and \eqref{P-Sn(3.5)} we have
\begin{align}\label{P-Sn(3.6)} \bigg| \P\Big\{S_n \ \hbox{ prime},\,  S_n\equiv \ell \ ({\rm mod}\, k)\Big\}&- \sqrt{2\over       \pi n }\sum_{j\in J}e^ {-  { j^2
/  2n  }  }\bigg| \cr &\le  \  {C\#\{J\}\over   n^{1+(\d/2)}}+  {2C\over (\log n)^{ b} (\log\log n)^{ 1/2}}.
\end{align}
%\eqno(3.6)
Whence
\begin{align}\label{P-Sn(3.7)} \P\Big\{S_n \ \hbox{ prime},\, S_n\equiv \ell \ ({\rm mod}\, k)\Big\}\le\   {\#(J)\over      \sqrt{ \pi n} }   + C' (\log n)^{-b}
(\log\log n)^{-1/2}.\end{align}
%\eqno(3.7) $$
Using the notation  $\pi(x;\ell, k):= \#\big\{p\le x: p\equiv \ell \ ({\rm mod}\, k)\big\} $, we have
 $$    \#\big\{2\le  p \le
\sqrt{2\a(1+b)n  \log\log n}:  p\equiv \ell \ ({\rm mod}\, k) \big\}  =  \pi(\sqrt{2\a(1+b)n  \log\log n};\ell, k)  . $$ 
By means of  the Brun-Titchmarsh theorem (see for instance \cite[p.\,83]{Te1}), 
$$\pi(x+y;\ell, k)-\pi(x;\ell, k) \le (2+o(1)) {y\over \varphi(k)\log (y/k)},
$$
 as $y/k$ tends to infinity, where  $\varphi$ is Euler's totient function. \vskip 2 pt 
 
 Letting $x=2$ and $y=\sqrt{2\a(1+b)n  \log\log n} -2$,  we have in particular, as   $
\sqrt{ n  \log\log n}/k\to
\infty$, 
\begin{eqnarray}\label{P-Sn(3.8)} {\pi(\sqrt{2\a(1+b)n  \log\log n};\ell, k)\over \sqrt
  n}&\le& { (2+o(1)) \over \sqrt
  n}  {\sqrt{2\a(1+b)n  \log\log n}\over \varphi(k)\log  (\sqrt{2\a(1+b)n  \log\log n}/k)   }\cr &\le  &C''   {\sqrt{   
\log\log n}\over \varphi(k)\log    n   }.
\end{eqnarray}
%\eqno(3.8) 

By combining  \eqref{P-Sn(3.8)} with \eqref{P-Sn(3.7)}, we deduce that for all  $n$ large enough, if  $
\sqrt{ n  \log\log n}/k\to
\infty$ 
 $$ \P\Big\{S_n \  \hbox{ prime},\,  S_n\equiv \ell \ ({\rm mod}\, k)\Big\}\le  C''   {\sqrt{   
\log\log n}\over \varphi(k)\log    n   } + C'{  1\over (\log n)^{ b}(\log\log n)^{ 1/2} }.
%\eqno(3.9)
 $$

Consequently, for any   $B>1$, there exists a constant $C(B)$, such that, as $n$ is large and  $
\sqrt{ n  \log\log n}/k\to
\infty$ 
$$ \P\Big\{S_n \  \hbox{ prime},\,  S_n\equiv \ell \ ({\rm mod}\, k)\Big\}\le C(B)\Big\{   {\sqrt{   
\log\log n}\over \varphi(k)\log    n   } +  {  1\over (\log n)^{B} }\Big\}.
%\eqno(3.10)
 $$
 \end{proof}


































 
 





\section{Divisibility and primality in the Cram\'er random walk}\label{s9}
%Let $\mathcal P=\{p_i,i\ge 1\}$ denote   the sequence of consecutive prime numbers. 
Cram\'er's probabilistic model basically consists with  a sequence of independent random variables $\xi_i$, defined for
$i\ge 3$ by
\begin{equation}\label{C.xin}\P\{\xi_i= 1\}= \frac1{\log i}, \qq \quad \P\{\xi_i= 0\}= 1-\frac1{\log i}.
\end{equation}
% Let   $\mathcal C=\{ j\ge 1 \!: \xi_j=1\}$. Let also $S_n =\sum_{i=1}^n \xi_i$, $n\ge 1$ and note that obviously $S_x=\#\big\{ \nu\in \mathcal C \,:\,\nu\le x\big\}$  for all positive integers $x$. 
Let   $m_n=\E S_n=\sum_{j=3}^n \frac1{\log j}$, $B_n={\rm Var }S_n= \sum_{j=3}^n \frac1{\log j}(1-\frac1{\log j})$. 
Cram\'er's well-known twin prime conjecture, which seems largely accepted,   does not however assert that there are any primes in  the \lq primes\rq\,  of the model,  which are  the instants  of jump   
of the random walk  $\{S_n , n\ge 1\}$,  recursively defined  by 
  \begin{equation}\label{C.Pn} P_1= \inf\{ n\ge 3: \xi_n=1\},\qq   P_{\nu +1} =\inf\{ n>P_\nu: \xi_n=1\}  \qq \quad \nu\ge 1  .
\end{equation}
 We refer   to Granville \cite{G},   Pintz \cite{P}  and Montgomery and Soundararajan \cite{MS}.    Our recent paper Weber \cite{W8} offers a     probabilistic complementary study, in particular     of   the \lq primes\rq~ $P_\nu$      and of the primality properties of  $\{S_n, n\ge 1\}$.
%Before doing the proof, 
Recall    
%a few standard facts and  some new ones  concerning this model,  first,  
that  the standard limit theorems
%  (strong law of large numbers, central limit theorem, law of the iterated logarithm, local limit theorem) 
are fulfilled,  
   %   are satisfied by the naturally associated random walk  $\{S_n , n\ge 1\}$, 
   and  the    invariance principle   holds, derived from Sakhanenko's. The law of the iterated logarithm implies that 
\begin{equation}\label{cramer.pnt}
 \#\big\{ \nu\in \mathcal C : \nu\le x\big\} = \int_2^x \frac{\dd t}{\log t} + \mathcal O\big( \sqrt{x\log\log x}\big) ,
\end{equation}
with probability one.          Cram\'er's  easier prediction on  the order of magnitude of the remainder   in the \lq   prime number theorem\rq, cannot be separated from the fact that this one should  also be sensitive to the subsequence on which it is observed, as    shown  in \cite{W8}. Such a property  is hardly conceivable  and was nowhere emphasized.  
 In a work in progress, that  study is  transposed to other random models for  prime numbers. 


%\medskip\par \vskip 5 pt
 \subsection{Divisibility}The distribution of divisors of $S_n$ turns up to share similarities with the one of Bernoulli sums. We prove the following result, which was only stated  in   \cite{W8}.
 \begin{theorem}\label{Cramer.div.Sn}
\begin{align*}   
\P\{d|S_n\}   \, = \frac{1}{d}  + \frac{1}{d}\sum_{  1\le  j\le \frac{d}{\pi}( \log n)  ( { \a \over 2 n} )^{1/2}} &  e^{ 2i \pi m_n(\frac{j}{d})  - 2\pi^2  B_n\,(\frac{j}{d})^2  }  +\mathcal O \Big(   \frac{d^3 \a^3(\log n)^4}{n^{3/2}}   \Big)
.
\end{align*}
\end{theorem}
    
  
  Let $\p_k(t)$ be the characteristic function of $\xi_k$, $\p_k(t)= \E e^{2i\pi \xi_k t}$. Let also 
$$\Phi_n (t)=\E e^{2i\pi tS_n}= \prod_{k=1}^n\p_k(t) ,  $$
be the characteristic function of $S_n$.  We  prove the following estimate with explicit constants.
\begin{proposition}\label{Cramer.Phin} We have
$|\Phi_n(t)|\le \exp\big\{- 2B_n\sin^2\pi t  \big\} $.
Further
\begin{equation*}
  \Phi_n (t)= e^{ 2i \pi tm_n -  2   B_n  (\pi t)^2  + E_n(t)  },
    \qq \quad {\text with}\qquad |E_n(t)| \,\le\,  12\,m_n\, (\pi  |t|)^3.
\end{equation*}
In particular,  
\begin{equation*}
  \E e^{ iy\, \frac{S_n-m_n}{\sqrt {B_n}} }= e^{    - \frac{y^2}{2}    + E_n(y)  }, \qq \quad {\text and}\qquad |E_n(y)| \, \le  2 \, \frac{\sqrt {\log n}\,|y|^3}{\sqrt {n}}   .
 \end{equation*}

 \end{proposition} 
The proof crucially uses the Lemma below. 
  
  \begin{lemma}\label{lfp1} Let $m  $ be a positive real and $p$ be a real such that $0<p<1$. Let $\b$ be a random variable defined by ${\mathbb P}\{ \b=0\}= p$, $
{\mathbb P}\{
\b=m\}= 1-p=q$.  Let $\p(t) ={\mathbb E\,} e^{2i\pi t \b}$. Then
 we have the following estimates,

\smallskip
\hspace*{1.2mm}{\rm (i)} For all real $t$,
 $|{\varphi}(t)|\le \exp\big\{- 2pq\sin^2\pi t m\big\} $

\smallskip
{\rm (ii)} If $q|\sin \pi t m|\le 1/3$, then
 \begin{eqnarray*}  \p(t)&=& e^{  q      2i\pi mt   - 2qp (  \pi mt)^2 +E}
 , 
\end{eqnarray*}
with  
$|E|\, \le \,   12 q (\pi m|t|)^3$.
    \end{lemma}
      For   sake of completeness we include the proof.
   \begin{proof} (i) One verifies that $|\p(t)|^2=1- 4 pq \sin^2\pi mt.$ As moreover $1-\t\le e^{-\t}$ if  ${\vartheta} \ge 0$, we obtain $|\p(t)|^2\le e^{-4 pq \sin^2\pi mt}$. 

\vskip 3 pt  (ii) Let  $|u|\le u_0<1$. From the series expansion of $\log(1+u)$,  it follows that
$$\log(1+u)=u-\frac{u^2}{2}+R, \qq \qquad|R|\le |u|^3\,\sum_{j=0}^\infty \frac{|\theta|^j}{3+j}\le \frac{|u|^3}{3(1- u_0 )}.$$
 Then $ 1+u= \exp\{ u -\frac{u^2}{2} +B\}$,  with $|B|\le C_0\,|u|^3$ 
and  $C_0=\frac{1}{3(1- u_0 )}$.

 Writing that  $\p(t)= 1 +
q\big( e^{2i\pi mt} -1\big)=1 + u$,
  where $|u|= 2q|\sin\pi mt|\le  2q\big(1\wedge   \pi m|t| \big) $, we obtain 
  \begin{equation}\label{phi.u} \p(t)= 1+u=e^{q ( e^{2i\pi mt} -1 ) }\, e^{  -\frac{u^2}{2} +B}, \qq 
\quad |B|\le 8C_0\,q^3 |\sin\pi mt|^3.
\end{equation}
   
   In order to   estimate $u^2$, we let  $A(t)= e^{2i\pi mt} -1-  2i\pi mt+ \frac{(2 \pi mt)^2 }{2}$, and    write  $ (  e^{2i\pi mt} -1 )^2$ under the form 
\begin{align*} 
   &\   A(2t)-2A( t) -\big\{  -1-  4i\pi mt+ 8 ( \pi mt)^2\big\} +2\big\{ -1-  2i\pi mt+ \frac{(2 \pi mt)^2 }{2}\big\}+1
\cr &=
   A(2t)-2A( t) -     ( 2\pi mt)^2.
\end{align*}
Then
$ ( e^{2i\pi mt} -1 )^2+     ( 2\pi mt)^2 =    A(2t)-2A( t)$,
 and so   
  $\frac{u^2}{2}=-  2q^2(  \pi mt)^2+  \frac{q^2}{2}(A(2t)-2A( t)).$ 
 
\vskip 3 pt
Let $u_0=\frac{2}{3}$ so that $C_0=1$. We assumed $q|\sin \pi t m|\le 1/3$, thus $|u|\le 2/3$. We consequently  get with \eqref{phi.u}, 
\begin{eqnarray}\label{phi.bound} \p(t)&=&  e^{q ( e^{2i\pi mt} -1 )  -\frac{u^2}{2} +B }\, =\, 
e^{q ( e^{2i\pi mt} -1 ) +    2q^2(  \pi mt)^2-  \frac{q^2}{2}(A(2t)-2A( t))+B}
 %\cr &=&e^{q A(t) + q (    2i\pi mt- \frac{(2 \pi mt)^2 }{2})+    2q^2(  \pi mt)^2-  \frac{q^2}{2}(A(2t)-2A( t))+B}
\cr &=&e^{  q      2i\pi mt   - 2qp (  \pi mt)^2 +H+B}
, 
\end{eqnarray}
with $H=q A(t) -    \frac{q^2}{2}(A(2t)-2A( t))$. By using 
 the   estimate (Lemma 4.14 in Kallenberg \cite{Ka})
 %(Lemma 4.14 in Kallenberg \cite{Ka}), 
 \begin{equation}\label{expit}\Big|e^{ix} -\sum_{k=0}^n \frac{(ix)^k}{k!}\Big|\,\le\, \frac{2|x|^n}{n!}\wedge \frac{ |x|^{n+1}}{(n+1)!},
 \end{equation}
valid for any $x\in\R$ and $n\in \Z_+$, and    letting $\d(x)=x^2\big(1\wedge   \frac{|x|}{6} \big)$, we get 
\begin{eqnarray}\label{At.est}
 |A(t)|&\le&
 \d(2\pi m|t|)
.\end{eqnarray}
Using the rough bound  $\d(x)\le  \frac{|x|^3}{6} $, and   $|B|
\le\,8q^3\big(1  \wedge \pi m | t|\big)^3$, we get 
 \begin{equation}\label{HB.bound} |H|+|B|\, \le \, 2q|A(t)|+\frac{q^2}{2}|A(2t)| +|B|\, \le \, \big(\frac83 q+\frac43q^2+8q^3\big)(\pi m|t|)^3
\, \le \,  12 q (\pi m|t|)^3.
\end{equation}
 We conclude by inserting estimate  \eqref{HB.bound} into \eqref{phi.bound}.  \end{proof}
 
 \begin{proof}[Proof of Proposition \ref{Cramer.Phin}]
   
We apply Lemma \ref{lfp1}. Here we have $m=1$, $q=\frac{1}{\log k}$. As $|{\varphi}_k(t)|\le \exp\big\{- 2(1-\frac{1}{\log k})(\frac{1}{\log k})\sin^2\pi t  \big\}$, we get 
 $$|\Phi_n(t)|\le \exp\big\{- 2\sum_{k=3}^n(1-\frac{1}{\log k})(\frac{1}{\log k})\sin^2\pi t  \big\}.$$
Further condition $q|\sin \pi t m|\le 1/3$ in Lemma \ref{lfp1}  reduces for $\p_k(t)$ to 
$ \frac{1}{\log k}|\sin \pi t  |\le 1/3$.  Thus 
\begin{equation}
  \label{chf.xik}
  \p_k (t)= e^{ 2i \pi    ( \frac{1}{\log k}) t - 2\pi^2  (1-\frac{1}{\log k})(\frac{1}{\log k} ) \,t^2  + C_k(t)  },
\end{equation}
where
 $|C_k(t)|
\le\, \frac{12  (\pi  |t|)^3}{\log k}$. 
  Recalling that $m_n =\sum_{j=3}^n \frac1{\log j}$, $B_n= \sum_{j=3}^n \frac1{\log j}(1-\frac1{\log j})$, it follows that 
\begin{equation}
  \label{chf.sn}
\E e^{2i\pi tS_n}=  \Phi_n (t)= e^{ 2i \pi tm_n -  2\pi^2  B_n  t^2  + D_n(t)  },
\end{equation}
and
$|D_n(t)| \,\le\,  12\,m_n\, (\pi  |t|)^3  $. 
In particular,  
 \begin{equation}
  \label{chf.sn1}
 \E e^{ iy\, \frac{S_n-m_n}{\sqrt {B_n}} }= e^{    - \frac{y^2}{2}    + E_n(y)  },
 \end{equation}
with  $|E_n(y)| \,\le\,  \frac{3m_n}{2} \, (\frac{y}{\sqrt {B_n}})^3\,\le\,  2 \, \frac{\sqrt {\log n}|y|^3}{\sqrt {n}}  $. 
  This achieves the proof.
\end{proof}
 

\begin{proof}[Proof of Theorem \ref{Cramer.div.Sn}]  We note that 
\begin{eqnarray}\label{div.Sn}\P\{ d|S_n\}
&=& \E \Big(\frac{1}{d}  \sum_{j=0}^{d-1} 
e^{2i\pi S_n \frac{j}{d}}\Big) \ = \ 
\frac{1}{d}  +\frac{1}{d}  \sum_{j=1}^{d-1} \Phi_n \big(  \frac{j}{d}\big)
\cr &= &\frac{1}{d}+ \frac{1}{d}  \sum_{j=1}^{d-1} 
e^{ 2i \pi m_n(\frac{j}{d})  - 2\pi^2  B_n\,(\frac{j}{d})^2+\mathcal O (m_n(\frac{j}{d})^3 ) } 
.\end{eqnarray}

The    elliptic Theta function $$  \Theta (d,m_n,B)  =  \sum_{\ell\in \Z} e^{im\pi{\ell\over   d }-{B\pi^2\ell^2\over 2 d^2}}, 
 $$
 appears in the above estimate of $\P\{d|S_n\}$. 
Here $m=m_n$, $B=B_n$ and we have $m_n\sim B_n\sim \frac{n}{\log n}$. 
\vskip 3 pt 
 Let 
 $\a>\a'>3/2$. Let
 $$\p_n= ( \log n) \big( { \a \over 2 n}\big)^{1/2},
\qquad\qquad \tau_n= {\sin\p_n/2\over
\p_n /2}.
%\leqno(2.4)
$$ 
We assume
$n$ sufficiently large, say $n\ge
n_0$, for $\tau_n$ to be greater than $(\a^\prime/\a)^{1/2}$.  Consider two sectors
  $$A_n = ]0,\p_n[,  \qq\qq A'_n=[ \p_n, {\pi\over 2} [.$$ 
 %$$\sum_{ j\ge 1\, :\, { \pi j\over d}\in A_n}=\sum_{1\le  j\le \frac{d}{\pi}( \log n)  ( { \a \over 2 n} )^{1/2}}$$

 We write 
 \begin{align*}   \sum_{j=1}^{d-1} 
&e^{ 2i \pi m_n(\frac{j}{d})  - 2\pi^2  B_n\, (\frac{j}{d})^2 +\mathcal O (m_n(\frac{j}{d})^3 ) }\cr &= \Big( \sum_{ j\ge 1\, :\, { \pi j\over d}\in A_n}+\sum_{ j\ge 1\, :\, { \pi j\over d}\in A'_n}
\Big)
e^{ 2i \pi m_n(\frac{j}{d})  - 2\pi^2  B_n\,(\frac{j}{d})^2+\mathcal O (m_n(\frac{j}{d})^3 ) } . 
\end{align*}  

  Concerning the second sub-sum,
the inequality $\p_n
\le {\pi j\over d}<{\pi \over 2}$ implies that  
\begin{eqnarray*}  \sum_{ j\ge 1\, :\, { \pi j\over d}\in A'_n} e^{ - 2\pi^2  B_n\,(\frac{j}{d})^2}&\le&
\sum_{ j\ge 1\, :\, { \pi j\over d}\in A'_n} e^{ - 2c   \frac{n}{\log n}\,(\frac{\pi j}{d})^2}
\cr &\le&\sum_{ j\ge 1\, :\,{ \pi j\over d}\in A'_n} e^{ -  c\a   \,  \log n       }\le
dn^{-  \a'  }.
\end{eqnarray*}
Therefore 
\begin{align*}  \Big| \sum_{ j\ge 1: { \pi j\over d}\in A'_n}  e^{ 2i \pi m_n(\frac{j}{d})  - 2\pi^2  B_n\,(\frac{j}{d})^2+\mathcal O (m_n(\frac{j}{d})^3 ) }\Big|
  &\le    \sum_{ j\ge 1:{ \pi j\over d}\in A'_n}  e^{  
 - 2\pi^2  B_n\,(\frac{j}{d})^2}\big(1   + \mathcal O ( m_n\,d ) \big)\cr &\le 
Cm_nd \sum_{ j\ge 1:{ \pi j\over d}\in A'_n} e^{ - 2c   \frac{n}{\log n}\,(\frac{\pi j}{d})^2}
\cr &\le Cm_nd\sum_{ j\ge 1: { \pi j\over d}\in A'_n} e^{ -  c\a   \,  \log n       }
\cr &\le C\big( \frac{n}{\log n}\big)d^2
 n^{-  \a'  }\, \le \, C d^2\,n^{1-  \a'  }.
\end{align*}
Now, concerning the first sub-sum, 
%as $\{j\ge 1: { \pi j\over d}\in A_n\}=\{j:1\le  j\le \frac{d}{\pi}( \log n)  ( { \a \over 2 n} )^{1/2}\}$,
\begin{eqnarray*}& &  \sum_{ j\ge 1\, :\, { \pi j\over d}\in A_n}  e^{ 2i \pi m_n(\frac{j}{d})  -  2\pi^2  B_n\,(\frac{j}{d})^2+\mathcal O (m_n(\frac{j}{d})^3 ) } 
\cr &= &\sum_{  1\le  j\le \frac{d}{\pi}( \log n)  ( { \a \over 2 n} )^{1/2}}e^{ 2i \pi m_n(\frac{j}{d})  - 2\pi^2  B_n\,(\frac{j}{d})^2+\mathcal O (m_n(\frac{j}{d})^3 ) }
\cr &=& \sum_{  1\le  j\le \frac{d}{\pi}( \log n)  ( { \a \over 2 n} )^{1/2}}e^{ 2i \pi m_n(\frac{j}{d})  - 2\pi^2  B_n\,(\frac{j}{d})^2  }\big(\,1+ \mathcal O (m_n(\frac{j}{d})^3 )\,\big)
%\cr &= \sum_{  1\le  j\le \frac{d}{\pi}( \log n)  ( { \a \over 2 n} )^{1/2}}e^{ 2i \pi m_n(\frac{j}{d})  - 2\pi^2  B_n\,(\frac{j}{d})^2  }\Big(1+m_n \mathcal O ( \sum_{  1\le  j\le \frac{d}{\pi}( \log n)  ( { \a \over 2 n} )^{1/2}}(\frac{j}{d})^3 )\Big)
.
\end{eqnarray*}
We have 
\begin{equation*} m_n  \sum_{  1\le  j\le \frac{d}{\pi}( \log n)  ( { \a \over 2 n} )^{1/2}}(\frac{j}{d})^3\,\le\, C\, \frac{d^2 \a^2(\log n)^3}{n} .
\end{equation*}
Thus
\begin{eqnarray*}&&\sum_{  1\le  j\le \frac{d}{\pi}( \log n)  ( { \a \over 2 n} )^{1/2}}
 e^{ 2i \pi m_n(\frac{j}{d})  - 2\pi^2  B_n\,(\frac{j}{d})^2  }\big(1+ \mathcal O (m_n(\frac{j}{d})^3 )\big) 
\cr &=&\sum_{  1\le  j\le \frac{d}{\pi}( \log n)  ( { \a \over 2 n} )^{1/2}}e^{ 2i \pi m_n(\frac{j}{d})  - 2\pi^2  B_n\,(\frac{j}{d})^2  }  
\cr &&\quad +\mathcal O \Big( \sum_{  1\le  j\le \frac{d}{\pi}( \log n)  ( { \a \over 2 n} )^{1/2}}e^{    - 2\pi^2  B_n\,(\frac{j}{d})^2  }  \Big)\Big( m_n  \sum_{  1\le  j\le \frac{d}{\pi}( \log n)  ( { \a \over 2 n} )^{1/2}}(\frac{j}{d})^3 \Big)
\cr &=&\sum_{  1\le  j\le \frac{d}{\pi}( \log n)  ( { \a \over 2 n} )^{1/2}}e^{ 2i \pi m_n(\frac{j}{d})  - 2\pi^2  B_n\,(\frac{j}{d})^2  }  
\cr &&\quad +\mathcal O \Big(   \frac{d^2 \a^2(\log n)^3}{n} \sum_{  1\le  j\le \frac{d}{\pi}( \log n)  ( { \a \over 2 n} )^{1/2}}e^{    - 2\pi^2  B_n\,(\frac{j}{d})^2  }  \Big)
\end{eqnarray*}
Whence, 
\begin{align*}  \sum_{ j\ge 1\, :\, { \pi j\over d}\in A_n} &e^{ 2i \pi m_n(\frac{j}{d})  -  2\pi^2  B_n\,(\frac{j}{d})^2+\mathcal O (m_n(\frac{j}{d})^3 ) } 
\cr &=\sum_{  1\le  j\le \frac{d}{\pi}( \log n)  ( { \a \over 2 n} )^{1/2}}e^{ 2i \pi m_n(\frac{j}{d})  - 2\pi^2  B_n\,(\frac{j}{d})^2  }  
  +\mathcal O \Big(   \frac{d^3 \a^3(\log n)^4}{n^{3/2}}   \Big) .
\end{align*}
Combining both estimates gives
\begin{align*}   
\P\{d|S_n\} \,=\, \frac{1}{d}  + \frac{1}{d}\sum_{  1\le  j\le \frac{d}{\pi}( \log n)  ( { \a \over 2 n} )^{1/2}} &  e^{ 2i \pi m_n(\frac{j}{d})  - 2\pi^2  B_n\,(\frac{j}{d})^2  }  +\mathcal O(d^2 n^{-\a'})
 \cr &\qq +\mathcal O \Big(   \frac{d^3 \a^3(\log n)^4}{n^{3/2}}   \Big)
 \cr \, = \frac{1}{d}  + \frac{1}{d}\sum_{  1\le  j\le \frac{d}{\pi}( \log n)  ( { \a \over 2 n} )^{1/2}} &  e^{ 2i \pi m_n(\frac{j}{d})  - 2\pi^2  B_n\,(\frac{j}{d})^2  }  +\mathcal O \Big(   \frac{d^3 \a^3(\log n)^4}{n^{3/2}}   \Big)
.
\end{align*}
This achieves the proof.
 \end{proof} 

 
Now  as $\sum_{j\ge d/2} e^{-A\frac{j^2}{2d^2}}\le e^{-cA }$, $c$ an absolute constant, for all $d\ge 2$ and $A\ge 2$, we have
\begin{align*}   
 \sum_{  j> \frac{d}{\pi}( \log n)  ( { \a \over 2 n} )^{1/2}}    e^{    - 2\pi^2  B_n\,(\frac{j}{d})^2  }& \le  \sum_{  j> \frac{d}{\pi}( \log n)  ( { \a \over 2 n} )^{1/2}}    e^{    - A(\frac{n}{\log n})(\frac{j}{d})^2  } 
 \cr &\le e^{    - A'(\frac{n}{\log n})(\frac{1}{d})^2(d( \log n))^2  ( { \a \over 2 n} ) )  } \cr & \le e^{    - A''\a \log n } =n^{- A''\a}.
\end{align*}
 
 We therefore also have
\begin{corollary}
\begin{eqnarray*}\Big|\P\{d|S_n\}-\frac{\Theta (d,m_n,B_n)}{d}\Big|\le C\ \frac{d^3 \a^3(\log n)^4}{n^{3/2}}  .
\end{eqnarray*}
\end{corollary}
 %\begin{eqnarray*}  \sum_{ j\ge 1\, :\, { \pi j\over d}\in A_n} e^{ - 2\pi^2  B_n\,(\frac{j}{d})^2}&\le& \sum_{ 1\le j  < d ( \log n)  ( { \a \over 2 \pi^2 n} )^{1/2} } e^{ - 2\pi^2  B_n\,(\frac{j}{d})^2} .
%\end{eqnarray*}



%\begin{eqnarray}\label{ } 
%  \sum_{ j\ge 1\, :\, { \pi j\over d}\in A_n}  \Phi_n \big(  \frac{j}{d}\big)
%&= & \sum_{ j\ge 1\, :\, { \pi j\over d}\in A_n}  
%e^{ 2i \pi m_n \frac{j}{d}  - 2\pi^2  B_n\,(\frac{j}{d})^2}\Big(1   + \mathcal O\big( m_n\,\sum_{ j\ge 1\, :\, { \pi j\over d}\in A_n} (\frac{j}{d})^3\big)  \Big) 
%\cr &= &   \sum_{j=0}^{d-1} e^{ 2i \pi m_n 
%\frac{j}{d}  - 2\pi^2  B_n\,(\frac{j}{d})^2}\big(1   + \mathcal O ( m_n\,d ) \big)\end{eqnarray}
We end  the paper with recent results related to   primality in the Cram\'er model.









\subsection{Quasiprime and prime property.}
We   obtained in   \cite{W8} -without assuming RH- a  sharp estimate  of $\P \{S_n\ \text{prime}\, \}$, for almost all $n$, namely for all $n$, $n\to\infty$ through a set $\mathcal S$  of natural density $1$. 
%We record the following result.  
\begin{theorem}\label{Snprime.M} 
{\rm (i)} For any constant $b>1/2$, 
\begin{equation}\label{Snprime.i}   \P \{S_n\ \text{prime}\, \}
\, =\, \frac{ 1 }{ \sqrt{2\pi  B_n  } }\,\int_{m_n-\sqrt{ 2bB_n\log n}}^{m_n+\sqrt{ 2bB_n\log n}}
 e^{- \frac{(t- m_n)^2}{    2  B_n   } } \, \dd \pi(t)  
    + \mathcal O\Big( \frac{(\log n)^{3/2 }}{\sqrt n} \Big),
  \end{equation}
as $n\to\infty$, where   $\pi(.)$ is the prime counting function.
  \vskip 3 pt {\rm (ii)} There exists a set of integers $\mathcal S$ of density 1, such that  \begin{equation}\label{abel.base,}  
 \P \{S_n\ \text{prime}\, \}
\, =\, \frac{ (1+ o( 1) )}{ \sqrt{2\pi  B_n  } } 
\int_{m_n-\sqrt{ 2bB_n\log n}}^{m_n+\sqrt{ 2bB_n\log n}}
  \,  e^{- \frac{(t- m_n)^2}{    2  B_n   } }\, \dd \pi(t),  \end{equation} 
as $n\to \infty$,  $n\in \mathcal S$. Further,
   \begin{equation}\label{abel.base,,}  
\liminf_{  \mathcal S\ni n\to\infty}\ (\log n)\P \{S_n\ \hbox{prime}  \}
\ge  \frac{1  }{  \sqrt{2\pi e}\, } .
   \end{equation}
 \end{theorem} 
%The proof uses a result of  Selberg  \cite{Se}.        \vskip 8 pt
% \subsection{Quasi-primality of ${S_n}$.}    \noi   
The proof  allows one (Remark \ref{rem.Snprime.M}) to treat other cases, so we   included it. \begin{proof}[Proof of Theorem \ref{Snprime.M}.]
(i) By Lemma 7.1 p.\,240 in \cite{P}, for $0\le x\le B_n$
\begin{eqnarray*} \P \{|S_n-m_n|\ge x \}&=& \P \{S_n-m_n\ge x \}
+\P \{-(S_n-m_n)\ge x \}
\cr &\le &2\,\exp\Big\{-\frac{x^2}{2B_n}\Big(1-\frac{x}{2B_n}\Big)\Big\},
\end{eqnarray*}
noticing that   $\{-\xi_j\}_j$ also satisfies the conditions of Kolmogorov's Theorem. Let $b>b'>1/2$.
Then for all sufficiently large $n$, since $\log  B_n\sim \log n$,
\begin{equation} \P \{|S_n-m_n|\ge  \sqrt{2bB_n\log   n} \}\le
2\,  n^{-b'}
.
\end{equation}

 We have \begin{align*}   \Big|  \P \{S_n \in\mathcal P \}&  -  \P \{S_n \in \mathcal P\cap[m_n- \sqrt{ 2bB_n\log n},m_n+\sqrt{ 2bB_n\log n}] \} \Big| 
\cr &\le    \P \{|S_n-m_n|\ge \sqrt{ 2bB_n\log n} \}
\cr &\le   n^{-b'}.
   \end{align*}
Further,
\begin{eqnarray*}& & \Big|  \P \{S_n \in \mathcal P\cap[m_n-\sqrt{ 2bB_n\log n},m_n+\sqrt{ 2bB_n\log n}] \}
\cr & &-\sum_{\k\in\mathcal P\cap[m_n-\sqrt{ 2bB_n\log n},m_n+\sqrt{ 2bB_n\log n}]} \frac{ e^{- \frac{(\k- m_n)^2}{    2  B_n   } } }{ \sqrt{2\pi  B_n  } } 
\Big|
 \cr &\le & \sum_{\k\in\mathcal P\cap[m_n-\sqrt{ 2bB_n\log n},m_n+\sqrt{ 2bB_n\log n}]}\Big|\P \{S_n =\kappa \} -\frac{ e^{- \frac{(\k- m_n)^2}{    2  B_n   } } }{ \sqrt{2\pi B_n   } }     \Big|   
  \cr     & \le &    C\,\#\big\{\mathcal P\cap[m_n-\sqrt{ 2bB_n\log n},m_n+\sqrt{ 2bB_n\log n}] \big\}\cdot
 \frac{(\log n)^{3/2}}{n} 
  \cr     & \le & 
  C\, \sqrt b\ \frac{(\log n)^{3/2 }}{\sqrt n} 
.
  \end{eqnarray*}
Therefore
\begin{equation} \label{3} \Big|  \P \{S_n \in \mathcal P  \}
 -\sum_{\k\in\mathcal P\cap[m_n-\sqrt{ 2bB_n\log n},m_n+\sqrt{ 2bB_n\log n}]} \frac{ e^{- \frac{(\k- m_n)^2}{    2  B_n   } } }{ \sqrt{2\pi  B_n  } } 
\Big|
     \, \le \, 
   C\,\sqrt b\  \frac{(\log n)^{3/2 }}{\sqrt n} 
.
  \end{equation}
 By    expressing the inner sum as a Riemann-Stieltjes integral  \cite[p.\,77]{A}, we get

\begin{equation}\label{abel.base}   \P \{S_n \in \mathcal P  \}
\, =\, \int_{m_n-\sqrt{ 2bB_n\log n}}^{m_n+\sqrt{ 2bB_n\log n}}
  \,\frac{ e^{- \frac{(t- m_n)^2}{    2  B_n   } } }{ \sqrt{2\pi  B_n  } }\, \dd \pi(t)
    + \mathcal O\Big( \frac{(\log n)^{3/2 }}{\sqrt n} \Big)
.
  \end{equation}

\vskip 5 pt (ii) We note that 
  \begin{eqnarray} \label{int.min}
  \int_{m_n-\sqrt{ 2bB_n }}^{m_n+\sqrt{ 2bB_n }}
  \,\frac{ e^{- \frac{(t- m_n)^2}{    2  B_n   } } }{ \sqrt{2\pi  B_n  } } \dd \pi(t)&\ge & L \  \frac{ \pi(m_n+\sqrt{ 2bB_n })-\pi(m_n-\sqrt{ 2bB_n })}{ \sqrt{   B_n  }},
\end{eqnarray}
with $L=\frac{e^{-b }}{  \sqrt{2\pi }}$. 
\vskip 3 pt
 We use  Theorem 4 in  Selberg  \cite{Se}. Let    $\Phi(x)$ be positive and increasing and such that $\frac{\Phi(x)}{x} $  decreasing for $x>0$. Further assume that 
 \begin{eqnarray} \label{phi.ab}
{\rm (a)}\quad \lim_{x\to\infty}  \frac{\Phi(x)}{x}= 0 \qq\qq {\rm (b)}\quad\liminf_{x\to\infty}\frac{\log \Phi(x)}{\log x}>\frac{19}{77}.
\end{eqnarray}
Then there exists a   set $\mathcal S$   of positive reals of density one such that
 \begin{equation}  \label{phi.ab.enonce}\lim_{\mathcal S\ni x\to \infty}\frac{\pi(x+\Phi (x))-\pi(x)}{({\Phi(x)}/{\log x})}  =1.
\end{equation}
 Let $\Phi(x) = \sqrt{2b  x}$. Then the   requirements in \eqref{phi.ab} are fulfilled, and so \eqref{phi.ab.enonce} holds true. 
  Now let $C$ be some possibly large but fixed positive number, as well as some positive real $\d<1/2$.    
    \vskip 5 pt   By \eqref{phi.ab.enonce}, the  set of $x>0$, call it $\mathcal S_\d$, such   that 
    \begin{equation}\label{pi.phi.d}
\pi(x+\Phi(x))-\pi(x) \, \ge  \,(1-\d)\, \frac{ \Phi(x)}{\log x}.
 \end{equation}
  has    density 1. 
  Note that   $\mathcal S_{\d'}\subseteq \mathcal S_{\d''}$ if  $\d'\le \d''$.   
 \vskip 3 pt Pick  $x\in\mathcal S_\d$ and let $\D(x)= \pi(x+\Phi(x))-\pi(x)$. Note that if $|y-x|\le C$,   $\big|\Phi(y)-\Phi(x)\big|=o(1)$ for $x$   large.
 Thus  for every $y\in [x-C,x+C]$, $\big|\D(y)-\D(x)\big|\, \le C'$, and so 
$$\D(y)   \, \ge  \,(1-\d)\, \frac{ \Phi(x)}{\log x}-C'.$$
 the constant $C'$ depending on $C$ only. 
 As $\big|\frac{\Phi(x)}{\log x}- \frac{\Phi(y)}{\log y}\big|\le 
\frac{C}{2\sqrt x \log x}$, we have 
\begin{equation*}
\D(y)   \, \ge  \,(1-\d)\, \frac{ \Phi(y)}{\log y}-C' -\frac{C}{2\sqrt x \log x}.
 \end{equation*} 
  
Thus every $y\in [x-C,x+C]$ also satisfies \begin{equation}\label{pi.phi.d.a}
\D(y)   \, \ge  \,(1-2\d)\, \frac{ \Phi(y)}{\log y},
 \end{equation} 
 if  $x $ is large enough.
 
\vskip 3pt   
Let $\nu=\nu(x)$ be the  unique integer   such that $m_{\nu-1}< x\le m_{\nu}$. As $m_{\nu}-m_{\nu-1}=o(1)$,  $\nu \to \infty$,  it follows that $m_{\nu}\in [x-C,x+C]$ provided that   $x$ is large enough, in which case we have by   \eqref{pi.phi.d.a},
  \begin{equation}\label{pi.phi.d.b}
 \frac{\pi(m_{\nu}+\Phi(m_{\nu}))-\pi(m_{\nu})}{\Phi(m_{\nu})}  \ge  \frac{1-2\d}{\log m_{\nu}}.
 \end{equation} 
   Let $X\ge 1$ be a large positive integer and $\e$  a small positive real. The number $N(X)$ of intervals $]\m-1,\m]$, $\m\le X$ such that $\mathcal S_\d \cap ]\m-1,\m]\neq \emptyset$ verifies $N(X)/X\sim 1$, $X\to \infty$, since $\mathcal S_\d$ has density 1. 
 
 Given such an $\m\le X$, pick $x\in \mathcal S_\d \cap ]\m-1,\m]$. 
 We know (recalling that $m_{\nu}-m_{\nu-1}=o(1)$,  $\nu \to \infty$) that some $m_{\nu}$, $\nu=\nu(x)$ belongs to   $ ]\m-1-\e,\m+\e]$, and that \eqref{pi.phi.d.b} is satisfied.   The union of these intervals $[\m-1-\e,\m+\e ]$ is contained in $[1-\e , X+\e ]$. It follows that the number of $\nu$ such that \eqref{pi.phi.d.b} is satisfied, forms a set of density 1. 
  
 \vskip 4 pt  We now use an induction argument in order to replace $2\d$ in  \eqref{pi.phi.d.b} by a quantity $\e(\nu)$ which  tends to 0 as $\nu$ tends to infinity  along some other set of density 1, which we shall build explicitly.    Let $\mathcal T_n$ be the set   of $\nu$'s of density 1, corresponding to $\d=\frac1n$, $n\ge 3$. Let $X_3$ be large enough so   that $\#\{\mathcal T_3\cap [1, X ]\}\ge X (1-1/3)$ for all $X\ge X_3$. Next let  $X_4>X_3$ be sufficiently large  so   that $\#\{\mathcal T_4\cap [X_3, X ]\}\ge X (1-1/4)$ for all $X\ge X_4$. Like this we manufacture an increasing sequence $X_j$,   verifying for all $j\ge 3$,
  $$\#\{\mathcal T_j\cap [X_{j-1}, X ]\}\ge X (1-1/j), \qq \qq {\rm for\ all}\  X\ge X_j .$$
  The resulting set 
  $$ \mathcal T= \bigcup_{j=3}^\infty \mathcal T_j\cap [X_{j-1}, X_j ]$$
  has density 1 and further we have the inclusions 
  $$\mathcal T\cap [X_{l-1}, \infty )=\bigcup_{j=l}^\infty \mathcal T_j\cap [X_{j-1}, X_j ]\subset \mathcal T_l\cap\bigcup_{j=l}^\infty   [X_{j-1}, \infty )=\mathcal T_l\cap [X_{l-1},\infty),\qq\quad l\ge 4,$$
as   the sets $\mathcal T_j$ are decreasing with $j$ by definition.   

\vskip 3 pt We finally have  by \eqref{int.min},
  \begin{equation}\label{pi.phi.d.b.1}
 \frac{\pi(m_{\nu}+\Phi(m_{\nu}))-\pi(m_{\nu})}{\Phi(m_{\nu})}  \ge  \frac{1-\e(\nu)}{\log m_{\nu}},
 \end{equation}
along $\mathcal T$, for some sequence of reals $\e(\nu)\downarrow 0$ as $\nu\to \infty$.
 
\vskip 3 pt Therefore 
  \begin{eqnarray}\label{base} 
  \int_{m_\nu-\sqrt{ 2bB_\nu }}^{m_\nu+\sqrt{ 2bB_\nu }}
  \,\frac{ e^{- \frac{(t- m_\nu)^2}{    2  B_\nu   } } }{ \sqrt{2\pi  B_\nu  } } \dd \pi(t)&\ge & L \  \frac{ \pi(m_\nu+\sqrt{ 2bB_\nu })-\pi(m_\nu-\sqrt{ 2bB_\nu })}{ \sqrt{   B_\nu  }}\cr & \ge &  L\, \frac{1-\e(\nu)}{\log m_{\nu}},
\end{eqnarray}
for all $\nu \in\mathcal T$, recalling that $L=\frac{e^{-b }}{  \sqrt{2\pi }}$.

Also  \begin{equation}\label{abel.base,proof}  
\liminf_{  \mathcal T\ni \nu\to\infty} (\log \nu)\P \{S_\nu\ \hbox{prime}  \}
\ge  \frac{1  }{ \sqrt{2\pi e}\, } .
 \end{equation}
  It further follows from \eqref{abel.base} that  
\begin{equation}\label{abel.base,,proof}  
 \P \{S_\nu\ \hbox{prime}  \}
\, =\, \big(1+ o( 1)\big)
\int_{m_\nu-\sqrt{ 2bB_\nu\log \nu}}^{m_\nu+\sqrt{ 2bB_\nu\log \nu}}
  \,\frac{ e^{- \frac{(t- m_\nu)^2}{    2  B_\nu   } } }{ \sqrt{2\pi  B_\nu  } }\, \dd \pi(t),
  \end{equation}
   all $\nu\in \mathcal T$. This completes the proof of Theorem \ref{Snprime.M}.
   \end{proof}
\begin{remark}\label{rem.Snprime.M}\rm 
The interested reader will have  observed that the property $m_\nu-m_{\nu-1}= o(1)$ is not fully used, in fact it suffices that  $m_\nu-m_{\nu-1}= \mathcal O(1)$, selecting at the beginning of the proof $C$ sufficiently large.  Thus the proof of Theorem \ref{Snprime.M} also adapt to the Bernoulli case. In particular there exists a set of integers $\mathcal U$ of density 1 such that
\begin{equation}\label{bernoulli.d1}  
\liminf_{  \mathcal U\ni \nu\to\infty} (\log \nu)\P \{B_\nu\ \hbox{prime}  \}
>0 .
 \end{equation}
\end{remark}
\vskip 4  pt 

 
 

  Let $\Pi_z=\prod_{p\le z}p$. According to Pintz \cite{Pi}, an integer $m$ is   $z$-quasiprime, if $(m, \Pi_z)=1$. Let $S'_n= \sum_{j= 8}^n \xi_j$, $n\ge 8$. Note that the introduction of $S'_n$ in place of $S_n$ is not affecting Cram\'er's conjecture. The probability that $S'_n$   be $z$-quasiprime is   for all $n$ large enough is studied in \cite{W8}, Theorem 2.6.
   The proof is based on a randomization argument.
     
 \begin{theorem}\label{cramer.quasi.prime.P}
 We have for any $0<\eta<1$, and all $n$ large enough  and      $ \zeta_0\le  \zeta\le \exp\big\{ \frac{c\log n}{\log\log n}\big\}$, 
\begin{eqnarray*}  \P\big\{  S'_n\hbox{\ $\zeta$-quasiprime}\big\}   \,\ge  \, (1-\eta)   \frac{ e^{-\gamma} }{  \log
\zeta }. \end{eqnarray*}
 where $\gamma$ is  Euler's constant and $c$ is a positive constant.\end{theorem} 

\vskip 3 pt  
 \subsection{Primality of ${P_n}$.} 
 %In the   next Theorem, we test which infinite sequences of primes  are ultimately avoided  by   the \lq primes\rq~  $ P_\nu$, with probability 1. 
We also showed in \cite{W8} that when the \lq primes\rq~  $ P_\nu$   are observed   along    moderately growing subsequences,  then with probability 1, they  ultimately avoid       any given infinite set of primes satisfying a reasonable  tail's  condition.  We also test which infinite sequences of primes  are ultimately avoided  by   the \lq primes\rq~  $ P_\nu$, with probability 1. 

\begin{theorem}\label{mathfrak}  Let $\mathcal K$ be an increasing sequence of naturals such that  
the series  $\sum_{k\in \mathcal K}  k^{-\b} $ converges for some   $\b\in ]0, \frac12[$.
Let $\mathfrak P$  be an increasing sequence of primes such that for some  $b>1$,
 $$\sup_{k\in\mathcal K} \frac{\#\{\mathfrak P\cap [k,  b k] \}}{  k^{\frac12-\b}}<\infty. $$
Let also  $\{\D_k,k\ge 1\}$ be the instants of jump of the Bernoulli sequence $ \{B_k,k\ge 1\}$.
%  defined in \eqref{delta.Delta}.
Then \begin{eqnarray*}
   \P\big\{   \D_k\notin\mathfrak P, \quad k\in \mathcal K\ \hbox{ultimately}\big\}=1.
    \end{eqnarray*}
Further, 
\begin{eqnarray*}
   \P\big\{   P_\nu\notin\mathfrak P, \quad \nu\in \mathcal K\ \hbox{ultimately}\big\}=1.
    \end{eqnarray*}
Moreover (case $\b=1/2$), let  $\mathfrak P$ be such that 
$\sum_{p\in \mathfrak P , \, p>y} p^{-1/2} = \mathcal O \big(y^{-1/2}\big)$, and $\mathcal K$ be   such that
  $\sum_{k\in \mathcal K}  k^{-1/2}<\infty$.
Then 
\begin{eqnarray*}
   \P\big\{   P_\nu\notin\mathfrak P, \quad \nu\in \mathcal K\ \hbox{ultimately}\big\}=1.
    \end{eqnarray*}
    \end{theorem} 
%The study made in  \cite{W8} is   at least   partly transposable to other models for the prime numbers.


 
 
 
 
 
 
 
 
 
 
 
 
 
 
 
 
 
  
  
  \section{Concluding Remarks.}\label{s10}
At the period we were interested in studying the divisors properties of the Bernoulli random walk, the topic, the study of arithmetical properties of binomial random walks was not much attracting attention, and we  worked   isolated.
%, to say the least, and we have worked isolated. %A common critic raised was that such study can be deduced from the uniform model, and so  does not bring anything new. 
\vskip 2 pt
The study of $\E f(B_n)$ for instance is quite tractable because there are exact formulas for $\P\{B_n=k\}$ or $\P\{d|B_n=k\}$,  $\P\{P^-(B_n)>\zeta \}$, $\P\{  B_n \,{\rm prime} \}$, \ldots, and standard tools from 
%(such as summation by parts, 
Fourier analysis    can be used to obtain results on $\E f(B_n)$ from standard number theoretic results; but this was not much taken into account. The used methods are intrinsic to the Bernoulli model. See   Remark \ref{unif.est.theta.alpha.rem}. %We don't know for instance any other estimate from the type given in Theorem \ref{estPdlBn}. 
%Quite importantly also, there is, 
As  indicated in subsection \ref{sub.bin.way.back}, there is a way back to the non random setting along Pascal matrices; 
%No doubt that this apparently overlooked 
%   thus      should draw more attention. 
 %It was ignored that 
thus a study of the Bernoulli case transfers to the  arithmetical one using inverse of Pascal matrices.  We found Call and Velleman's result only recently.
%\vskip 2 pt 
%Recall that any arithmetical function $f(n)$ can be represented in the form
%$$ f(n) =\sum_{d|n} \Phi(d), \qq \qquad \Phi(n)= \sum_{d|n}\mu\Big(\frac{n}{d}\Big) f(d).$$
 For    any arithmetical function $f(n)$
 %by Proposition \ref{way.back}, 
 we have the reverse formula
 $$ f(i)=\sum_{j=1}^i  (-1)^{i-j}2^j{{i-1}\choose{j-1}}\E f(B_j),\qq\quad i\ge 2.
$$ 
%$$\E f(n) =\sum_{d|n} \P\{d|B_n\}\,\Phi(d).$$
Now the research in this area has received a recent growing interest.
%, 
%considered, its interest was nearly called into questions around the author. 
We steadily found the subject attractive because of the combined interaction 
%between
 of Arithmetic and Probability  which can just be promising, and are convinced  that much remains to discover, and to do. The topic has large potential of development.
 %Now the topic has received a recent growing interest. This is a satisfaction as we were always convinced about its large potential of development, also 
 %convinced about the fact that much remains to discover, and to do. 
 %We refer to the survey of Fernandez and Fernandez \cite{FF}. 
\vskip 2 pt
Random walks are also an important tool of investigation  in Analytic  Number Theory.  This topic is the object of the recent work Weber \cite{W15}.     
  \appendix
   
%   \label{appendix} 
  

  
 %%%%%%
%%%%%%
\section{Erd\"os-R\'enyi's model}\label{s5.a}
%%%%%%%%
%%%%%%%%

    Erd\"os and R\'enyi \cite{ER} have   investigated models for integers involving  infinite sequences of  independent   random
variables $\b_n$, with $\P\{\b_n=0\} =1-\P\{\b= 1\}=\varrho_n
$. They thorougly studied
 their additive properties in relation with classical problems of   additive representation  of integers such as Waring problem. See
Halberstam and Roth \cite{HR}, see also Landreau \cite{L}.   

\section{Kubilius' model} The Bernoulli model  only involves  one
probability space: the infinite product of probability spaces      $(\{0,1\}^\N, \mathcal B(\{0,1\}^\N),
\m^\N )$, where  $\m= \frac{1}{2}(\d_{0}+\d_{1})$,  $\d_x$ denoting the Dirac mass at point $x$. By construction,  the Kubilius model  relies upon all
probability spaces
$$\hbox{$(\{ 1,2, \ldots ,n \}, \m_n)$}$$
  $\m_n$ being  the normalized counting measure.
This is an obstacle for the application of general criteria of  almost sure convergence, typically G\'al-Koksma, Stechkin or M\'oricz theorems. 
   However,
     a bridge exists in Kubilius theory with some standard infinite  product of probability  spaces, namely with some inherent random walk built up from
independent non identically distributed random variables. 
%This can be presented as follows.

  Let
$\{Y_p, p\ge 1\}$ be a sequence of independent binomial random variables such that
$\P\{Y_p=1\}=1/p$ and
$\P\{Y_p=0\}=1-1/p$. We can view $Y_p$ as modelling whether   an integer taken at random is divisible by $p$ or not. Let 
$$ T_n=\sum_{p\le n}  Y_p.  $$
By Mertens estimate
 $s_n^2=\E (T_n-\E T_n)^2= \sum_{p\le n} {1\over p}-{1\over p^2}=\log\log n +\mathcal O(1)$.  
 The sequence $\{T_n, n\ge 1\}$ is known to asymptotically behave as the truncated prime divisor function
$$\o(m,t)=\#\{ p\le t : p|m\} ,  $$
  at least when $t$ is not too close to $m$.  More precisely, let
$$\o_r(m)=\big( \o(m,1), \ldots, \o(m,r)\big), $$
where $r$ is some integer with $2\le r\le x$, and put $u={\log x\over \log r}$. The bridge mentionned before is given by the following sharp estimate.
\begin{theorem} \label{km}  Given $c<1$ arbitrary, we have uniformly in
$x,r
$ and
$Q\subset \Z^r$,
$$  \frac{\#\{ m\le x: \o_r(m)\in Q\}}{ x}=\P\big\{(T_1, \ldots, T_r)\in Q\big\}+ \mathcal O\big( x^{-c}+ 
e^{-u\log u}\big).
$$
\end{theorem} 
   We refer   to Elliott \cite{E} p.119-122 and   Tenenbaum \cite{Te} where a thorough study is presented.
 \vskip 2pt  The error term $e^{-u\log u}$  
naturally    brings  restrictions in   applications   to asymptotic problems.  To make it small,
it requires     if
$r=r(x)$ that
$  r(x)=\mathcal O _\e (x^\e)$   for all $\e>0$.   This amounts to truncate  the prime divisor function $\o(m )$  at level
$\mathcal O _\e (x^\e)$, which is satisfactory as long as $m\ll x$. However, these integers have a negligible contribution on the size
of    the left-term in Theorem  \ref{km}.   The  model is therefore mostly adapted to the analysis of the distribution of   small prime divisors
of an integer (see notably \cite{E} p.122). 
 %%%%%%%%%%%%%%%%%%%%%%%%%%%%%%%%%%%%%%%%%%%%%%%%%%%%%%%%%%%%%%%%%%%%%%
 \section{Proof  of Theorem \ref{romult1}.}   
 
  The following lemma  is   classical.  
\begin{lemma} \label{romult}Let $f\in \Z(X)$ and put $\varrho_f(d)=\#\big\{0\le y<d: d|f(y)\big\}$. Then $\varrho_f$ is a
multiplicative function.
\end{lemma}
 
%\begin{comment}
\vskip 4 pt 
\noi   {\it Comment:}  One way to \lq expedite\rq \, the proof is to say that it is \lq\lq well-known and trivial\rq\rq, or to say that \lq\lq it is proved using the Chinese Remainder Theorem\rq\rq. We have here a proof ready and  are of mind we can't get out of it so. Further the present paper also address to probabilists, who are not necessarily    acquainted  with  arithmetical arcana.   
%\end{comment}   \vskip 3 pt 
 We provide a simple and complete proof. We couldn't find it in the existing literature,  where however,  different but sometime sketchy proofs are  available.
%A proof is given in Appendix for the sake of completeness. 
  \begin{proof}
  %[\bf Proof of Lemma \ref{romult}]
  Write $f(x)= a_0+a_1x+\ldots + a_nx^n$, $a_j\in \Z$, $0\le j\le n$. Let $d=d_1d_2$ with $(d_1,d_2)=1$. We first
show that
$\varrho_f(d_1)\varrho_f(d_2)\le
\varrho_f(d )$. Let
$(y_1, y_2)$ be such that $0\le y_i<d_i$ and   $d_i|f(y_i)$, $i=1,2$. 
\vskip 3 pt
Now as $(d_1,d_2)=1$,   there exists an 
%unique
 integer $y$, $0\le y<d$ such that (\cite[Th.\,59]{HW}),
\begin{equation}\label{equiv}\hbox{$y\equiv y_i\  {\rm mod}(d_i)$,\qq$i=1,2$.}
\end{equation}
% (\cite{HW}, Th. 59) ,
% \bibitem{HW} G. H. Hardy et E. M. Wright,      \emph{An introduction to the theory of numbers}, fifth edition, Oxford at the Clarendon Press, Oxford, (1979).
Writing that  $y=y_1+\ell d_1$ we have
 \begin{eqnarray*}f(y)&=& a_0+a_1(y_1+\ell d_1)+\ldots + a_n(y_1+\ell d_1)^n \cr & =&a_0+(a_1 y_1+d_1 A_1  ) +\ldots + (a_n y_1^n +  d_1A_n)=
f(y_1)+ d_1 F ,
\end{eqnarray*}
where $A_1, \ldots, A_n$, $F $ are integers. As   $d_1|f(y_1)$ by assumption, it follows that   $d_1|f(y)$. Writing now that  $y=y_2+m d_2$, we similarly get  $d_2|f(y)$, whence $d|f(y)$. 
\vskip 3 pt
Now let $(y'_1, y'_2)$ be such that
$$\hbox{$0\le y'_i<d_i$ and   
$d_i|f(y'_i)$, \qq $i=1,2$.}$$ 

Let also $y'$ integer, be 
%the   unique  integer
  satisfying 
$$\hbox{$0\le y'<d$, \qq $y'\equiv
y'_i\  {\rm mod}(d_i)$,\quad $i=1,2$,}$$ 
   thereby implying that  $d|f(y')$ by the above argument. 
\vskip 3 pt
 
 We have the implication: $y=y'\Rightarrow (y_1, y_2)=(y'_1, y'_2)$. Indeed as $y=y'$,  we have
$$y= y_1+\ell d_1=y_2+k d_2=y'_1+\ell' d_1=y'_2+k' d_2.$$
And so $y_1-y'_1=  (\ell'-\ell) d_1$. Since $0\le y_1, y'_1<d_1$, this implies $y_1=y'_1$. Similarly we get $y_2=y'_2$, so that $(y_1,
y_2)=(y'_1, y'_2)$. 
Therefore $\varrho_f(d_1)\varrho_f(d_2)\le
\varrho_f(d )$. 

\vskip 3 pt Conversely, let $0\le y<d$ be such that $d|f(y)$. Let $y_1, y_2$, $0\le y_i<d_i$ be such that  $y_i\equiv y \, {\rm
mod}(d_i)$, $i=1,2$. Then, in the same fashion
\begin{eqnarray*}f(y_1)&=& a_0+a_1(y +\ell d_1)+\ldots + a_n(y +\ell d_1)^n \cr &=&a_0+(a_1 y +d_1 B_1  ) +\ldots + (a_n y ^n +  d_1B_n)\cr & =&
f(y )+ d_1 G.
\end{eqnarray*}
Thus $d_1|f(y_1)$ and  similarly $d_2|f(y_2)$. 

\vskip 2 pt Now let   $0\le y'<d$ be such that $d|f(y')$, and let $(y'_1, y'_2)$ be the corresponding
pair of integers. We shall prove the implication $  (y_1, y_2)=(y'_1, y'_2) \Rightarrow y=y'$. 

Write $y_1=y+\ell d_1$,
$y'_1=y'+\ell' d_1$. If $y_1=y'_1$, then $y-y'=(\ell'-\ell)d_1 $ so that $d_1|y-y'$. Similarly $y_2=y'_2$ implies $d_2|y-y'$. Thus
$d |y-y'$. As $0\le y,y'<d$ this implies that $y=y'$. Hence the requested implication.
% $  (y_1, y_2)=(y'_1, y'_2) \Rightarrow y=y'$. 
We deduce  that
$ 
\varrho_f(d )\le \varrho_f(d_1)\varrho_f(d_2) $. The proof is now complete. 
\end{proof}

\vskip 7 pt
%\noindent {\bf Acknowledgments} 
%\vskip 3 pt  I thank  Profs. Dubickas, Laurin\v cikas,   and M\"uller for helpful comments on  this problem, and Prof.  McKee for   useful remarks on the papers \cite{McK1}, \cite{McK}.  
 
 
%   \section{Proof.}\label{s1}    
   %The proof is based on the lemma below.
     
     
     
    Before passing to the proof of Theorem \ref{romult1}, recall some related results. In McKee \cite{McK},   counting functions   associated with   quadratic forms $y^2 + by+c$ are investigated, under the restriction $\D= b^2 -4c<0$. As explained in section 2,  the number of solutions to the  congruence
 $$y^2 + by+c\equiv 0 \ {\rm (mod}\, d{\rm )}, \quad 0\le y<d,$$
 equals the number of solutions to the  congruence
 $y^2 \equiv 0 \ {\rm (mod}\, 4d{\rm )}$, $b\le y<b+2d$.  The condition $\D <0$ is not satisfied in the case we consider.   
   \vskip 2 pt
   The case of irreducible polynomials is mainly adressed in the literature and we couldn't find anything corresponding to the solved problem   in this work.    That question was already considered in Weber \cite{W4}.  For results concerning summatory functions associated with counting functions of the number of divisors of irreducible polynomials, we refer to Broughan \cite{Br}, Hooley \cite{Ho}, McKee \cite{McK1} and Scourfield \cite{S}, among  other contributors.  
 % We note that we can also obviously write $\varrho_f(d)$ 
 %in Lemma \ref{romult}
% as  $\varrho_f(d)=\#\big\{1\le y\le d: d|f(y)\big\}$. 
\begin{proof}[\bf Proof of Theorem \ref{romult1}] 
The proof is  long and technical because  of many sub-cases entering into consideration, which is natural when studying this kind of question, according to sayings of patented number theorists colleagues. It can be used to treat similar questions. First consider the   case  $k>0$, which is the
main case. Recall that
\begin{eqnarray*}
%\label{rho}
\varrho_k(D)=\#\big\{1\le y\le  D: \  D| y^2+ky \big\} ,\qq \quad k=0,1,\ldots
\end{eqnarray*} 
The function $\varrho_k(D)$ being multiplicative,
it suffices to compute $\varrho_k(p^r)$. Observe to begin that 
$$\varrho_k(p^r)=2, \quad r=1,2,\ldots\qq\quad \hbox{if $p\not| k$}. $$ 
 Indeed, if $(y,p)=1$, then $p^r|y+k$ and there is only one solution given by
$y= \tilde k$ or $y=p^r-\tilde k$ where $\tilde k$ is the residue class of $k$ modulo $p^r$. If $y=p^s Y$, $(Y,p)=1$, $1\le s<r$, then $p^r| y(y+k)\Leftrightarrow p^{r-s} |(p^sY+k)$. And so $p|k$,
which was excluded. There is thus no solution of this kind. Finally, it remains one extra solution $y=p^r$. Thus $\varrho_k(p^r)=2$.
\vskip 5pt     We now assume that $p|k$. We can range the solutions $y$ of the equation $p^r|y(y+k)$ in disjoint classes of type $y=p^sY$,
with
$(Y,p)=1$. When
$r=1,2$ or
$3$, there is a direct simple computation.  
\smallskip\par   {\it Case $r=1$.}  We have $\varrho_k(p)=\#\big\{1\le y\le p: \,  p| y(y+k)  \big\}$. 
 Since $p |k$, it follows that  $p| y(y+k)\Leftrightarrow p|y^2$. As $y\le p$, the prime divisors of $y$ are less or equal to $p$ and    so are the prime divisors of $y^2$. There is thus only one solution: $y=p$.
Whence $    \varrho_k(p)=1$.
 \smallskip\par    {\it Case $r=2$.}   We have $\varrho_k(p^2)=\#\big\{1\le y\le p^2: \,  p^2| y(y+k)  \big\}
$. As $p\| k$, 
   If $(y,p)=1$, then $p^2| y(y+k)\Leftrightarrow p^2|y+k$ and so $p|y$, which is
impossible. There is no solution of this type. 
 If $y=pY$, with  $(Y,p)=1$, then  $p^2| y(y+k)\Leftrightarrow p|(pY+k)$, which  is satisfied. 
The number of solutions is $\#\{ Y\le p:
(Y,p)=1\}=p-1$. The case $y=p^2Y$, with  $(Y,p)=1$ coincides with $y=p^2$.
 Therefore,  
 $    \varrho_k(p^2)=   p-1+1=p$.
\smallskip\par {\it Case $r=3$.}   We have $\varrho_k(p^3)=\#\big\{1\le y\le p^3: \  p^3| y(y+k)  \big\}. 
$
     If $(y,p)=1$, then $p^3| y(y+k)\Leftrightarrow p^3|y+k$ and so $p|y$, which is
impossible. There is no solution of this type. If $y=pY$, with  $(Y,p)=1$, then  $p^3| y(y+k)\Leftrightarrow p^2|(pY+k)$. If
$v_p(k)\ge 2$, this implies that $p|Y$, which is impossible, and so there is no solution of this type in that case. If $v_p(k)=1$, write
$k=p K$, we get $p|Y+ K$.  The number of solutions is
$\#\{ Y\le p^2 :Y\equiv -K\, {\rm mod}(p)\}=p $.  If $y=p^2Y$, with  $(Y,p)=1$, then  $p^3| y(y+k)\Leftrightarrow p |(p^2Y+k)$, which is
always realized. The number of solutions is $\#\{ Y\le p: (Y,p)=1\}=p-1$. 
  The extra case is
$y=p^3 $.
 Therefore   
\begin{equation}    \varrho_k(p^3)=\begin{cases} 
                                   2p\quad &{\rm if}\ \quad v_p(k)=1
 \cr
                                   p  \quad &{\rm if}\ \quad v_p(k)\ge 2.   
            \end{cases}
\end{equation}
An additional   sub-case thus appears. 
Summarizing
\begin{equation}\varrho_k(p^r)=\begin{cases}  
 1\ &{\rm if} \    r=1 , 
 \cr
 p\ &{\rm if}\  r=2 ,   
  \cr
                                   2p  \quad &{\rm if}\   r=3 , \   v_p(k)=1 . 
\cr
                                   p  \ &{\rm if}\   r=3 , \   v_p(k)\ge 2 . 
\end{cases}  
\end{equation}

 \medskip {\it Case $r\ge 4$.} Put 
 $r' =\lfloor{r\over 2}\rfloor$.  Then
$r' 
\ge 2
$. We have
$\varrho_k(p^r)=\#\big\{1\le y\le p^r:
\  p^r| y(y+k) 
\big\}. 
$  
  If $(y,p)=1$, then $p^r| y(y+k)\Leftrightarrow p^r|y+k$ and so $p|y$, which is
excluded and there is no solution. 

Apart from the trivial solution $y=p^r$, the other possible solutions are of type $y=p^s
Y$,  
$(Y,p)=1$,
$1\le s<r$; and we shall distinguish three cases:   
$$\hbox{(i) $r'<s<r$, \qq (ii) $s=r'$,\qq (iii) $1\le s< r'$.}$$
 
\noi (i) Since $r'<s<r$, then $ r/2 \le  s $, and so $1\le r-s\le s$.  Further 
$p^r| y(y+k)$ means $p^{r-s}|  Y(p^sY+k)$ or
$p^{r-s}|  p^sY +k  $,   which is possible if and only if
$p^{r-s}|  k
$, namely $ r-s \le v_p(k)
$.  Thus 
$$\max(r'+1,r-v_p(k)) \le s<r.$$
 We have $Y\le p^{r-s}$, $(Y,p^{r-s})=1$. Their number
  is   $\phi(p^{r-s})$ where $\phi$ is Euler's function, and since $\phi(p^{r-s})=p^{r-s}(1-{1\over p}) $, 
the corresponding number of solutions is 
\begin{eqnarray*} \sum_{\max(r'+1,r-v_p(k)) \le s\le r -1 } p^{r-s}(1-{1\over p})   &=&   (1-{1\over p})\sum_{1\le  v  
\le (r-r'-1)\wedge v_p(k)}  
p^{v}
\cr &  =&  p\, {p^{(\lfloor {r-1\over 2}\rfloor )\wedge v_p(k)}-1\over p-1}(1-{1\over p}) 
\cr &  =&  
p^{\lfloor {r-1\over 2}\rfloor\wedge v_p(k)} -1   . 
\end{eqnarray*}
   (ii)      We consider   solutions of type    $y=p^{r'}Y$, $(Y,p)=1$. 
\smallskip\par  ---  If $r$ is odd, $r=2r'+1$, then  $p^{2r'+1}| p^{ r' }Y(p^{ r' }Y+k)$ means
$p^{ r'+1}|   (p^{ r' }Y+k)$. So $p^{ r'}|k$ and thereby $v_p(k)\ge r'$.
  If $v_p(k)< r'$, there is   no solution.  If $v_p(k)>  r'$, it implies that  $p|Y$ which is impossible and   there is   again no  
solution. 
 The remainding case $v_p(k)=  r'$ is the only one providing solutions. Write $k=p^{ r' }K $, $(K,p)=1$, then $p|Y+K$. Since
$(K,p)=1$, the solutions are the numbers
$Y$ such that
$1\le Y\le  p^{ r' +1}$ and 
$Y\equiv -K\, {\rm mod}(p)$. Let $1\le \kappa<p$ be such that $K  \equiv \kappa\, {\rm mod}(p)$.
The number of solutions is
\begin{equation} \#\big\{Y\le  p^{ r' +1} : Y\equiv -\kappa\  {\rm mod}(p)\big\} =\#\big\{ - \kappa  +jp: 1\le j\le p^{ r' }\big\}= p^{
r'  }.
\end{equation}
   --- If $r$ is even, $r=2r'$, then  $p^{2r' }| p^{ r' }Y(p^{ r' }Y+k)$ reduces to $p^{ r' }|  (p^{ r' }Y+k)$, so $p^{
r'}|k$. If $v_p(k)< r'$, there is no solution.  If $v_p(k)\ge   r'$,  write $k=p^{v_p(k) }K$, $(K,p)=1$, this is always realized and   the
number of solutions is
$$\#\big\{Y\le  p^{ r' } :  (Y,p)=1\big\} =\phi(p^{ r' })=p^{ r' }(1-{1\over p}) .$$ 
 \par \noi (iii) We consider the last type of  solutions: $y=p^sY$, $(Y,p)=1$,    $1\le s< 
r'$.  Notice first,  since $s<r'$ that  $s<r/2$, and so $r-s>r/2>s$.
As
$p^r| y(y+k)$ means 
 $p^{r-s}|   Yp^s+k $, we deduce that $p^s|k$, namely $ s\le v_p(k) $. 
 If $  v_p(k)<s $, there is no solution.  
  If $   v_p(k)>s $, then $p|Y$ which is impossible.  

If $s= v_p(k) $, which requires $v_p(k)<r'$, write $k=p^{v_p(k)}K$, $(K,p)=1$. Then
we get the equation $p^{r-2v_p(k)}|   Y +K $; so  $Y  
\equiv -K \, {\rm mod}(p^{r-2v_p(k)})$.  
  Let $1\le \kappa<p^{r-2v_p(k)}$ be such that $K  \equiv \kappa\, {\rm mod}(p^{r-2v_p(k)})$.
The number of solutions is
\begin{eqnarray*} \#\big\{Y\le p^{r-v_p(k)}:    Y   &\equiv& -\k \  {\rm mod}(p^{r-2v_p(k)})\big\}
\cr &=&\#\big\{  -\kappa +jp^{r-2v_p(k)}:  1\le j\le p^{v_p(k)}\big\}
\cr &=&p^{v_p(k)}.
\end{eqnarray*}
  \vskip 7pt
 Summarizing the case $r\ge 4$, if $r$ is odd, $r=2r'+1$ 
\begin{equation}     \varrho_k(p^r) = \begin{cases}   2p^{v_p(k)} & \quad\hbox{\rm if   $1\le v_p(k)\le r'$},
   \cr 
        p^{r'} & \quad\hbox{\rm if   $v_p(k)> r'$}.
   \end{cases}
\end{equation}
And if $r$ is even, $r=2r'$  
\begin{equation}       \varrho_k(p^r) =  \begin{cases}    
   2p^{v_p(k)} & \quad\hbox{\rm if  $1\le v_p(k)< r'$},
   \cr
 p^{r'} & \quad\hbox{\rm if   $v_p(k)\ge r'$} .
\end{cases}
\end{equation}
This remains true if $r=1,2,3$. Observe that ($v_p(k)< r'$, $r$ even) or ($v_p(k)\le  r'$, $r$ odd) are equivalent to $v_p(k)<{r\over
2}$.  
\smallskip\par Therefore, for
all $r\ge 2$ 
 
\begin{equation}        \varrho_k(p^r) =  \begin{cases}    2  & \quad {\rm if}  \   p\not| k  ,
   \cr  2p^{  v_p(k)} & \quad\hbox{\rm if}\  1\le v_p(k)<{r\over
2} ,
   \cr 
        p^{\lfloor{r\over
2}\rfloor}   \ & \quad\hbox{\rm if} \  v_p(k)\ge {r\over
2} .
     \end{cases}
\end{equation}

Consequently \begin{equation} \label{valro}\varrho_k(D)= \prod_{p|D  }\varrho_k(p^{  v_p(D) })=  \prod_{ v_p(k)<{v_p(D)/
2}}(2p^{v_p(k)})\cdot\prod_{   v_p(k)\ge {v_p(D)/
2} }  p^{\lfloor{v_p(D)\over
2}\rfloor} .
\end{equation} 
Recall that 
$$  D_{1/2} \,=\,\prod_{p|D} p^{\lfloor {v_p(D)\over 2}\rfloor }.$$

We consider again several cases:
\vskip 3 pt 
--- If $p$ is such that $v_p(D)=1$, then it produces in the writing of $\varrho_k(D)$ only a factor 1, since $v_p(k)<1/2$, so $v_p(k)=0$   and $\lfloor{v_p(D)\over
2}\rfloor=0$. 
\vskip 2 pt 
--- Now let $p$ be such that $v_p(D)\ge 2$, and compare $v_p(D)/2$ with $\lfloor{v_p(D)\over
2}\rfloor$. If $v_p(D)$ is even, both quantities are equal. If $v_p(D)$ is odd, $v_p(D)= 2z+1$, then 
$v_p(D)/2=z + 1/2$,  
$\lfloor{v_p(D)\over
2}\rfloor= z$. 
\vskip 2 pt
Thus in the latter case:
\vskip 2 pt
 
a) If $v_p(k)<{v_p(D)/2}$, namely $v_p(k)\le z$, we get the factor $2p^{v_p(k)}= 2p^{\min(v_p(k),z)}=2p^{  v_p((k,D_{1/2})) }$. If $v_p(k)\ge{v_p(D)/2}$, namely $v_p(k)> z$, we have the factor $p^z=p^{\min(v_p(k),z)}=p^{  v_p((k,D_{1/2})) }$. 

 
\vskip 2 pt
 
b) In the case $v_p(D)$ is even, if $v_p(k)<{v_p(D)/2}$, namely $v_p(k)< z$, we have the factor $2p^{v_p(k)}=2p^{\min(v_p(k),z)}=2p^{  v_p((k,D_{1/2})) }$. If $v_p(k)\ge {v_p(D)/2}= z$, we have the factor $p^{z}=p^{\min(v_p(k),z)}=p^{  v_p((k,D_{1/2})) }$.

\vskip 2 pt 

  Therefore the  contribution of 
 a prime $p$ such that $v_p(D)\ge 2$, is   either $2p^{  v_p((k,D_{1/2})) }$ if $p$ is part of the first product, or $p^{  v_p((k,D_{1/2})) }$ if $p$ is part of the second. 
 
 A prime number $p$ is part of the first product, if and only if the inequation $1\le v_p(k)<v_p(D)/2$ is fulfilled.
 %, which is equivalent to saying that $p|k$, and $p^3|D$. 
 %As $p^3|D\Leftrightarrow p^2|p^{\lfloor {v_p(D)\over 2}\rfloor }$, 
 We 
therefore get the condensed form, 
%\begin{equation} \label{valro1}\varrho_k(D) \,=\, 2^{\#\{p\,:\, p|k\,,\, p^3|D \}}\, (k, D_{1/2}).
%\end{equation} 
\begin{eqnarray} \label{valro1}\varrho_k(D)= \prod_{p|D  }\varrho_k(p^{  v_p(D) })&= & \prod_{ v_p(k)<{v_p(D)/
2}}(2p^{v_p(k)})\cdot\prod_{   v_p(k)\ge {v_p(D)/
2} }  p^{\lfloor{v_p(D)\over
2}\rfloor}\cr  &=&2^{\#\{p\,:\, 1\le v_p(k)<v_p(D)/2 \}}\,\prod_{p} p^  {\min (v_p(k),v_p(D)/2)} 
\cr  &=&2^{\#\{p\,:\, 1\le v_p(k)<v_p(D)/2 \}}\, (k, D_{1/2}).
\end{eqnarray} 




\vskip 5 pt Finally,   consider the case  $k=0$, namely       $\varrho_0(p^r)=\#\big\{1\le y\le p^r: \  p^r| y^2   \big\}  $. Notice that 
$\varrho_0(p)=\#\big\{1\le y\le
p: \  p| y   \big\}=1 $. Let
$r>1$, and write   $y=p^sY$, $(Y,p)=1$ and $1\le s\le r$. If $s=r$, there is the trivial unique solution $y=p^r$. If $1\le s<r$,
 $p^r| y^2\Rightarrow 2s\ge r$, in which case, the number of solutions is  
$$\#\big\{Y\le p^{r- s}: (Y,p)=1\big\}=\phi(p^{r- s})=p^{r- s}(1-{1\over p}).  $$
 Consequently
 $\varrho_0(p^r)= 1+\sum_{r/2\le s< r}p^{r- s}(1-{1\over p})$ . 
If $r$ is even, $r=2r'$
$$\varrho_0(p^r)= 1+\sum_{\s =1}^{r' } p^{\s}(1-{1\over p})= 1+p\sum_{u =0}^{r'-1 } p^{u}(1-{1\over p})= 1+p\Big({p^{r' }-1\over p-1}\Big)
 ( {p-1\over p})=p^{r' } ,$$
whereas if $r$ is odd, $r=2r'+1$,
 $\varrho_0(p^r)= 1+\sum_{\s =1}^{r'} p^{\s}(1-{1\over p})=p^{r' }$.
Thus 
$$\varrho_0(p^r)=p^{\lfloor {r\over 2}\rfloor }. $$
It follows that 
$$\varrho_0(D)=\prod_{p|D} p^{\lfloor {v_p(D)\over 2}\rfloor }\, =\, D_{1/2}.  $$
The proof is now complete. 
\end{proof}


 \vskip 7 pt
 \noindent {\bf Acknowledgments} 
 \vskip 3 pt  I    thank  Sanying Shi   for  exchanges     around divisibility  in Bernoulli random walk.  
 





 %%%%%%%%%%%%%%%%%%%%%%%%%%%%%%%%%%%%%%%%%%%%%%%%%%%%%%%%%%%%%%%%%%%%%%%%%%%%%%
 
\begin{thebibliography}{99}
 \bibitem{A} T. M. Apostol, \emph{Introduction to analytic number theory}. Springer-Verlag, New York-Heidelberg, (1976), Undergraduate Texts in Mathematics.
 
  \bibitem{A1} T. M. Apostol, \emph{Modular Functions and Dirichlet series in number theory}, Second Edition. Springer-Verlag, New York-Heidelberg, (1990), Graduate Texts in Mathematics.
  
   \bibitem{B}   R.  Bellman,  \emph{A brief introduction to Theta functions}, (1961), Holt, Rinehart and  Winston, New-York.

\bibitem{Br} K. A. Broughan, (2002) Restricted divisors sums. \emph{Acta arithmetica} {\bf 101.2},
105--114.
 
    \bibitem{BE} J. Bureaux and N. Enriquez,   The probability that two random integers are coprime, \emph{Math. Nachr.} {291} (2018) no. 1, 24--27.  

  \bibitem{Bi} P. Billingsley,   \emph{  Convergence of probability measures}, (1968) John Wiley \& Sons, Inc.,
 New York-London-Sydney.  
 
 \bibitem{CV}  G.\,S. Call and D.\,J. Velleman, Pascal's matrices, \emph{Amer. Math. Monthly} {\bf 100},
  (1993) 372--376.
 
 \bibitem{CFF}  J. Cilleruelo, J.\,L. Fern\'andez and P. Fern\'andez, Visible points in random walks, \emph{Eur. J. Combin. } {\bf 75} (2019),
  92--112. 
  
   \bibitem{C}  H. Cram\'er, (1936) On the order of magnitude of the difference between consecutive primes, \emph{Acta
Arithmetica} {\bf 2},
  23-46.  

\bibitem{DR} H. Daboussi and J. Rivat, (2000) Explicit upper bounds for exponential sums over primes, \emph{Mathematics of computation} {\bf 70}, Number 233, 431--447.


\bibitem{DS}   P. Diaconis  and Ch. Stein,  Some tauberian theorems related to coin tossing,  \emph{Ann. Prob.},   {\bf 6} No3  (1978) 483-490.  

%\bibitem{DW} Dvoretsky, A. and  Wolfowitz  J., Sums of random integers reduced mod $m$,  \emph{Duke Math. J.}  {\bf 18} No 2 (1951), 501--507.

 
  \bibitem{E}  P. D.  Elliott,   (1980) \emph{Probabilistic number theory
II}, Springer, New York.  
 
 
 \bibitem{ER} P. Erd\"os and A. R\'enyi, Additive properties of random sequences of positive integers, \emph{Acta Math.} {\bf 6}, (1960), 83--110.


 \bibitem{EZ} P. Erd\"os and S. K. Zaremba, (1972) The arithmetical function $\sum_{d|n} \frac{\log d}{d}$,
\emph{Demonstratio Math.} {\bf 6}, Part. 2,  575--579.


\bibitem{F}  {W. Feller}, (1968)   \emph{An Introduction to Probability Theory  and its Applications}, {\bf 1}, 3rd ed. Wiley,  New York. 


  \bibitem{FF} J.\,L. Fernandez and  P.  Fernandez, 
  Divisibilities properties of random samples of integers, RACSAM (2021), 115: 26. 
  
  
 % \bibitem{FP}   Freiman G. A., Pitman J. (1994)  {\sl Partitions into distinct large parts},  J. Austral. Math. Soc. (Series A)  {\bf57}, 386--416.
 

 \bibitem{HR} {  H. Halberstam and  R. F. Roth},   \emph{Sequences}, Springer-Verlag, (1983).
 
 
 \bibitem{Hu}  M. N. Huxley, On the differences between consecutive primes, \emph{Invent. Math.} {\bf 15} (1972), 164--170.

 \bibitem{G}  A. Granville, Harald Cram\'er and the Distribution of Prime Numbers, \emph{Scand. Actuarial J.}    No. 1, (1995) 12--28.



 \bibitem{H1}  G.H. Hardy  \emph{Divergent series}, (1963), Oxford at the Clarendon
Press.
  
  \bibitem{H}  G.H. Hardy,   An inequality for Hausdorff means, \emph{J. London Math. Soc.} {\bf 18}  (1943), 46--50.

  \bibitem{HW} {G. H. Hardy and E. M. Wright},     \emph{An introduction to the theory of numbers}, fifth edition, Oxford at the Clarendon Press, Oxford, (1979).
 
 \bibitem{HB}  D. R. Heath-Brown, The number of primes in a short interval, \emph{J. reine angew. Math.} {\bf 389} (1988), 22--63. 
 
\bibitem{HT}  A. Hildebrand and G. Tenenbaum, Integers without large factors, \emph{J. Th\'eorie des nombres, Bordeaux} {\bf 5} (1993), 1--74.  

\bibitem{Ho} C. Hooley, (1963) On the number of divisors of quadratic polynomials. \emph{Acta Mathematica} {\bf 110},
97--114.

\bibitem{McK1}  J. McKee,  (1999)  The average number of divisors of an irreducible quadratic polynomial,  \emph{Math.
Proc. Cambridge Philos. Soc.} {\bf 126}, 17--22.


\bibitem{McK}  J. McKee,  (1996) A note on the  number of divisors of quadratic polynomials,
 in  \emph{Sieve Methods, Exponential Sums, and their Applications in Number Theory}, edited by Greaves, Harman and Huxley, CUP 1996, 275--287.

 \bibitem{S} E. J. Scourfield, (1961) The divisors of a quadratic polynomial, \emph{Proc. Glasgow Math. Soc.} {\bf 5}, 8--20.


 \bibitem{J}   W.B. Jurkat,       Ein funktionentheoretischer Beweis f\"ur $\mathcal O$-Taubers\"atze bei den Verfahren non Borel und Euler-Knopp, \emph{Arch. Math.} {\bf 7}  (1956),
278--283.


\bibitem{Ka} O. Kallenberg,  \emph{Foundations of modern probability theory}, (1997)  Springer Verlag New-York.

  \bibitem{L}  B. Landreau,      \'Etude probabiliste des sommes de puissances $s$-i\`emes, \emph{Compositio} {\bf 99} n01  (1995), 1--31.\bibitem{LN}    W. J. Leahey, and J. E. Nymann,  On the probability that an integer chosen according to the binomial  distribution be $k$-free,  \emph{ The Rocky Mountain J. of Math.} {\bf 7}, No 4 (1977)  769--774.
 
 % \bibitem{Ma}  H. Maier,  (1984)    Primes in short intervals, \emph{Michigan Math. J.} {\bf 32}, 221--225.


% \bibitem{Me}  M. Merca,   A note on Cosine Power Sums, \emph{J. Integers Sequences} {\bf 15},  (2012) Article 12.5.3.

%\bibitem{LW} M.  Lifshits and  M. Weber (2009)  Sampling the Lindel\" of conjecture with the Cauchy random walk,   \emph{Proc. London Math. Soc.}
%{\bf 98} (3),    241--270. 


 \bibitem{MC}  P. J. McCarthy,  {\it Introduction to Arithmetical Functions}, (Universitext), (1986), Springer-Verlag New-York Inc.

  \bibitem{di}  MacDiarmid,  C.,  (1998).  Concentration, \emph{Prob. Methods for Algorithmic Discrete Math.}, 195--248, Algorithms Combin. {\bf
16}, Springer, Berlin.

\bibitem{M} D. MacDonald,    A local limit theorem for large deviations of sums of independent, non-identically distributed
random variables, \emph{Annals of Prob.} {\bf 7} (1979) no.\ 3, 526--531.

\bibitem{MT}   K. Matom\"aki and J. Ter\"av\"ainen, A note on zero density results implying large value estimates for Dirichlet polynomials,   arXiv:2403.13157v1, (2024).

%\bibitem{M1}  M.  Mikol\'as, (1949) Farey series and their connection with the prime problem {(I)}, \emph{Acta Sci. Math. Szeged} {\bf 13},   93--117.
 
% \bibitem{M2}  M.  Mikol\'as, (1949) {\rm   Farey series and their connection with the prime
%problem (II)}, \emph{Acta Sci. Math. Szeged,} {\bf 14},   5--21.
%\bibitem{Mi}  D.S.  Mitrinovi\'c,   \emph{Analytic inequalities}, Springer Verlag {\bf 165}, (1970).

   \bibitem{Mo}    T. T.  Moh,  On a general Tauberian theorem,  \emph{Proc. Amer. Math. Soc.} {\bf 36}, No 1 (1972) 167--172.

 \bibitem{MV} {H. L.  Montgomery, R. C.  Vaughan} (2007) \emph{Multiplicative Number Theory I. 
 Classical Theory}, Cambridge studies in advanced math, Cambridge University Press, %Cambridge, United Kingdom.
 \bibitem{MS} {H. L.  Montgomery, K. Soundararajan}, Primes in Short Intervals \emph{Commun. Math. Phys.} {\bf 252}, 589--617, (2004).

% \bibitem{QG} J. M. Quoniam, and M. G. Greening, A trigonometric summation, \emph{Amer. Math. Monthly}, {\bf 75} (1968), 405--406.

\bibitem{P}  V. V.   Petrov,     \emph{Sums of Independent Random Variables}, Ergebnisse der Math. und ihre
 Grenzgebiete  {\bf 82}, (1975) Springer. 

   \bibitem{Pi} J. Pintz,    Cram\'er vs. Cram\'er,  On Cram\'er's probabilistic model for primes, \emph{Funct. Approx. Comment. Math.} {\bf 37}, part 2, (2007) 361--376.


\bibitem{Po} A. G. Postnikov,
\emph{Introduction to analytic number theory}, AMS Translation of mathematical monographs {\bf
68}. (First publ. in Russian in 1971.) Amer. Math. Soc., 1988.



\bibitem{R}  A. R\'enyi,   \emph{Foundations of probability}, (1970)
  Holden-Day series in probability and statistics.
 
%\bibitem{Ri} J. Riordan, \emph{Combinatorial Identities}, John Wiley \& Sons, (1968).
 
  \bibitem{RS} J. B. Rosser, L. Schoenfeld, (1962) Approximate formulas for some functions of prime numbers, \emph{Illinois J. Math.} {\bf 6}(1), 64--94.
  
 \bibitem{Ro}  Y.  A. Rozanov,  
 On a local limit theorem for lattice distributions,
 \emph{Theor. Prob.   Appl.}, {\bf 2} (1957),  no.\ 2, 260--265.
 
 
   \bibitem{Sa}  A. S\'ark\H ozy A., (1978)  On difference sets of sequence of integers, \emph{Acta Math. Acad. Sci. Hungar.} {\bf 31} (1-2),
%125--149.
   \bibitem{Se}  A. Selberg, On the normal density of primes in small intervals, and the difference
between consecutive primes. \emph{Arch. Math. Naturvid.}, {\bf 47} (1943), 87--105.
 \bibitem{Si}  H. Siebert, (1976)   Montgomery's weighted sieve for dimension two,
\emph{Monatsh. Math.} {\bf 82},    \mbox{327--336.}
%\bibitem{Sh}
%H. N. Shapiro,   (1950) On a theorem of Selberg and generalization, \emph{ Ann. Math.} {\bf 51}, 485--497.
\bibitem{Sh.We}
 S. Shi and M. Weber,   On  Jordan  double  sums and related summatory functions,  \emph{Unif. Distribution Theory}  {\bf 18}   no.\,2, (2024),   p.\,57--76.
 \bibitem{Sh.We1}
 S. Shi and M. Weber, (2022) Value distribution of squarefree divisors of Bernoulli sums, {\it Preprint}.
 
  
\bibitem{Sp} V. G. Sprind\v zuk, (1979)
\emph{Metric theory of Diophantine approximations}, Translated from the Russian and edited by Richard A. Silverman. With a foreword by Donald J. Newman. Scripta Series in Mathematics. V. H. Winston \& Sons, Washington, D.C.; AHalsted Press Book, John Wiley \& Sons, New York-Toronto, Ont.-London. xiii+156 pp.  


\bibitem{SW}   Szewczak, Z.   and  Weber M., 
 Classical and Almost Sure  Local Limit Theorems,  
%{arXiv:2208.02700v1}, (2022), 
    \emph{Dissertationes Math}, No {\bf 589},    (2023), 97 pp.

  \bibitem{Te1} G. Tenenbaum,   \emph{Introduction \`a la th\'eorie analytique et  probabiliste des nombres}, (2008), Coll. \'Echelles Ed. Belin  Paris.
 \bibitem{Te}  G. Tenenbaum,      Crible d'Eratosth\`ene et mod\`ele de Kubilius, in:    Gy\'ory, Iwaniec,
Urbanowicz (Eds), \emph{Number Theory  in Progress, Proceedings of the Conference in honor of Andrzej Schinzel}, Zakopane, Poland 1997,
W. de Gruyter, Berlin-New-York
 (1999), 1099--1129. 
 
  \bibitem{To} L. T\'oth, A survey on Gcd-Sums functions,  \emph{Journal of Integer Sequences}, Vol. 13 (2010), Article 10.8.1.
 
\bibitem{Wal}  A. Walfisz, \emph{Weylsche Exponentialsummen in der neueren Zahlentherie.}, Mathematische Forschungsberichte, {\bf XV}, V E B Deutscher Verlag der Wissenschaften (1963), Berlin.

  \bibitem{W15}  M.  Weber,    Random   walks in analytic number theory, (2024).
   
  \bibitem{W14}  M.  Weber,     Correlation  Properties    of    Divisors 
   in the Bernoulli random   walk,
 %  random   walks - Tools and methods,
  (2023).    
 \bibitem{W11} M.  Weber, On infinite M\"obius inversion,    \emph{Publ. Math. Debrecen} Ref. no.: 9260, (2023), 1--12.

\bibitem{W13} M. Weber, On Farey sequence and quadratic Farey sums, \emph{ Research in Number Theory}, {\bf 8}, Article
number {\bf 14}, (2022)  (25 p.). 

\bibitem{W12}    M. Weber  An extension of a result of Erd\"os-Zaremba,
 \emph{Glasgow Math. J.}, {\bf 63} (1), (2021), 193--222. 

\bibitem{W10} M.  Weber, On Rozanov's theorem and strenghtened asymptotic  uniform distribution, 
%arXiv:2209.12228v1
\emph{Probab. Math. Statist.},  (2024), to appear.

 \bibitem{W9} M.  Weber,    A uniform semi-local limit theorem along    sets of multiples   for sums of i.i.d. 
   random variables, 
 %  arXiv:2209.12223v1, (2022), 
 \emph{Funct.
Approx. Comment. Math.}, (2024), to appear.  
     
\bibitem{W8} M.  Weber,  %(2021)
      Critical probabilistic characteristics of the
Cram\'er model for primes and arithmetical properties,  
  %arXiv:2105.11020v1, 
  \emph{Indian J. of Pure and Applied Math},   (2024), to appear. 
    \bibitem{W3} M.  Weber,  (2013) Instants of small amplitude of the Brownian motion   and application to    the 
 Kubilius model,  
\emph{Periodica Math. Hung.} {\bf 67}  (1), 95--113.   

 \bibitem{W2a}  M. Weber,   (2022) A Fourier analysis  of quadratic   Riemann sums and local integrals of  $ {\zeta(s)}$, \emph{submitted}.
 
 \bibitem{W2}  M.  Weber,   \emph{Dynamical Systems and Processes},  (2009)  European Mathematical Society Publishing
House, IRMA Lectures in Mathematics and Theoretical Physics {\bf 14}  
   xiii+761p.

 \bibitem{W5} M.  Weber, (2008) Sampling the integers with a complete random walk (27 p.), \emph{Unpublished}.

 \bibitem{W5a} M.  Weber, (2008) Sur la probabilit\'e $\P\{ S_n\, \hbox{est premier}\}$ (15 p.), \emph{Unpublished}.

 \bibitem{W1}  M.  Weber,    Small divisors of Bernoulli sums,   \emph{Indag. Math.} {\bf 18} No2 (2007),   281--293.  

\bibitem{W4} M.  Weber, (2007) Divisors of Bernoulli sums,    Results der Math. {\bf 51}, 141-179.

  \bibitem{W7} M.  Weber (2006) On the order of magnitude
of the divisor function,   \emph{Acta Math. Sinica} {\bf 22}, No2, 377-382. 

 %\bibitem{W5b} M.  Weber, (2006) Value distribution of squarefree divisors of Bernoulli sums  (14 p.), \emph{Unpublished}.

  \bibitem{W7a} M.  Weber, (2005) Divisors, spin sums  and the functional equation of the
 Zeta-Riemann function, \emph{Periodica Math. Hungar.}  {\bf 51} (1) 119-131.
   \bibitem{W6} M.  Weber, (2004)  An arithmetical property of
Rademacher sums,  \emph{Indag.  Math. \emph{(}N.S.\emph{)}}  {\bf 15}, 133--150.
 
 
   \end{thebibliography}  




 %%%%%%%%%%%%%%%%%%%%%%%%%%%%%%%%%%%%%%%%%%%%%%%%%%%%%%%%%%%%%%%%%%%%%%%%%%%%%%
 
 \end{document}
 
 
 
 
   \section{Concluding remark.}
  
  
  
  
  
  
We end up the paper with a    remark on distribution of lattice points visible from the origin  
We refer to section 3.8 in Apostol \cite{A}.
We say  that two lattice points $P$ and $Q$ are mutually visible  if the segment which joins them contains no lattice points other than the endpoints $P$ and $Q$. We have the following characterization: 
\medskip \par 
{\it Two lattice points $(a,b)$ and $(m,n)$ are mutually visible if, and only if $(a-m, b-n)= 1$.}
  \vskip 6 pt Equivalently two lattice points $(a,b)$ and $(m,n)$ are mutually visible if and only if $(a-m,b-n)$ is visible from the origin.
\smallskip \par
Consider now a square region defined by  $|x|\le r$, $|y|\le r$, $r$ large. Let $\mathcal N(r)$ be the number of   lattice points in this square which are visible from the origin. By taking care of symmetries, this number is 
\beq
\mathcal N(r)= 8+8\sum_{2\le n\le r} \sum_{1\le m<n \atop (m,n)=1} 1= 8\sum_{1\le n\le r} \p(n),
\eeq
$ \p(n)$ being Euler's totient function.
\vskip 1 pt The following classical estimate   is easy to prove: for $x>1$,
\beq  \label{estim.phi}
\Phi(x)= \sum_{1\le n\le x} \p(n) = \frac{3}{\pi^2}\, x^2 + \mathcal O(x\log x),
\eeq
As the total number of lattice points in the square is 
$$(2[r]+1)^2 = 4r^2+ \mathcal O(r) $$
it follows that: 
\begin{equation}\label{limF}\lim_{r\to \infty} \frac{\#(\mathcal N(r))}{4r^2 }= \frac{6}{\pi^2},\end{equation}
namely, 
\medskip \par 

 \qq \qq {\it The set of lattice points visible from the origin has density $6/\pi^2$.}
\medskip \par 
 
%Estimate \eqref{estim.phi} is easy to prove. 
Estimate \eqref{limF} can be strongly improved.   There is no loss, because of obvious symmetries, to rather  consider the set $\mathcal Q(n)$ of lattice points  in the  quadrant      $1\le k < \ell \le n$,  which are visible from the origin. 
\medskip  \par
Then for any Riemann integrable function on $[0,1]$, 
\begin{equation}\label{limF.f}\lim_{n\to \infty} \frac{1}{\#(\mathcal Q(n))}\, \sum_{(k,\ell) \in \mathcal Q(n)   } f\big(\frac{\k}{\l}\big)
=\frac{6}{\pi^2}\,\int_0^1 f(x)\dd x .\end{equation} 
  By taking $f(x)\equiv 1$, one recovers   \eqref{limF}.
\vskip 2pt
The problem of estimating the error term
  is connected with the Riemann Hypothesis (RH). There is a one to one correspondence between $\mathcal Q(n)$ and the Farey series of order $n\ge 1$
    $$\F_n=\big\{ \frac{j}{m}, 1\le j\le m\le n, (j,m)=1\big\},$$
  through the map $ (j,m)\mapsto \frac{j}{m}$
from $\mathcal Q(n)$ onto  $\F_n$.     %, see  
 % Mikol\'as  \cite{M1}, 
Let also $f:\Q\cap [0,1]\to \C$  arbitrary and $0< \s\le 1$ and define\begin{equation} \label{F_n(f)}F_n(f)= \sum_{\frac{\k}{\l}\in \F_n} f\big(\frac{\k}{\l}\big)
%, \qq F_{n,\s}(f)= \sum_{\frac{\k}{\l}\in \F_n} \frac{1}{\k^\s\l^\s}f\big(\frac{\k}{\l}\big)
\qq n=1,2,\ldots.
\end{equation} It is well-known \cite{M1} (see also \cite{W2},   Chapter  11) that the Farey fractions are uniformly distributed (mod 1).
 Thus  by the Weyl criterion,  for any Riemann integrable function on $[0,1]$,
\begin{equation}\label{limF.f} \lim_{n\to \infty}\frac{1}{\Phi(n)}\,F_{n}(f)= \int_0^1 f(x)\dd x ,\qq \quad n\to \infty.\end{equation}
    The limit  actually holds true (\cite[Th.\,2]{M1}) for any function $f$ defined on rational points of $[0,1]$, whose  Riemann sums are converging. 
Whence \eqref{limF} follows.
 \smallskip \par      Now note that $\Phi(n) =\sum_{d=1}^n dM(\frac{n}{d})$ where  $M(x)= \sum_{ \l \le x}\m(\l)$,  $\m(n)$  being the M\"obius function. Further $M(n)$ is also the Farey sum
\begin{equation}\label{mobius.sum}M(n)= \sum_{\frac{\k}{\l}\in \F_n} \cos 2\pi(\frac{\k}{\l}).
\end{equation}
%by M\"obius' inversion.
 By a result of  Littlewood,  the  RH  is equivalent to
the assertion  \begin{equation}\label{RHM}M(x)=\mathcal O_\e\big(x^{\frac12 +\e}\big). \end{equation}
Thus the  simplest example of a smooth periodic function $f(x) = \cos 2\pi x$,  thus  shows that the problem of estimating the error term ($=F_n(f)$ here)  is out of reach.  See also the recent work \cite{W13} for   unconditional results on  Farey  sums and quadratic Farey sums.

 
  

  

\section{Concluding Remarks.}\label{s10}
At the period we were interested in studying the divisors properties of the Bernoulli random walk, the topic was not much attracting attention, to say the least. Now the research in this area has received a recent growing interest.
%, 
%considered, its interest was nearly called into questions around the author. 
We steadily found the subject attractive because of the combined interaction 
%between
 of Arithmetic and Probability  which can just be promising, and are convinced  that much remains to discover, and to do. The topic has large potential of development.
 %Now the topic has received a recent growing interest. This is a satisfaction as we were always convinced about its large potential of development, also 
 %convinced about the fact that much remains to discover, and to do. 
 %We refer to the survey of Fernandez and Fernandez \cite{FF}. 
\vskip 3 pt

 
 
 
 
 $$[]\hbox{\bf partie non termin\'ee, \`a finir ou \`a sortir} []

 
\begin{proof}
 Let $n>m$. We have 
$$\P\{S_n\, \hbox{ prime} \}=\sum_{k=0}^\infty \P\{S_n\in  J_k(n) \}=\sum_{k=0}^{k_n} \P\{S_n\in  J_k(n) \}+\rho_n . $$
\begin{align*}
\P\{S_n &\, \hbox{and $S_m$ are prime} \}-\P\{S_n\, \hbox{ prime} \}\, \P\{S_m\, \hbox{ prime} \}\cr &=\sum_{k,\ell=0}^\infty
\Big[\P\{S_n\in  J_k(n), S_m\in  J_\ell(m) \}-\P\{S_n\in  J_k(n)\}\P\{ S_m\in  J_\ell(m) \}\Big] . 
\end{align*}

And  
$$\P\{S_n\in  J_k(n), S_m\in  J_\ell(m) \}-\P\{S_n\in  J_k(n)\}\P\{ S_m\in  J_\ell(m) \}
$$$$=\sum_{q\in J_k(n)\atop p\in J_\ell (m)}
\Big[\P\{S_n=q, S_m=p \}-\P\{S_n=q\}\P\{ S_m=p \}\Big] $$
But 
\begin{eqnarray}\label{radem.prim(1.7)}
{\bf Cov}\big({\bf 1}_{\{S_n =q \}}{\bf 1}_{\{S_m =p \}}\big)
&:=&  
 \P
\{S_n =q, S_m =p \}-\P \{S_n =q \}\P \{S_m =p \}  
\cr &=& \P \{S_m =p\}\P\{ S_n-S_m
=q -p\}
\cr & &\qq\qq \qq\qquad-\P
\{S_n =q \}\P \{S_m =p \} 
\cr &= &  \P \{S_m =p\}\big(\P\{ S_{n-m} =q -p\}- \P \{S_n =q \}\big).
\end{eqnarray}
% \eqno(1.7)$$
Hence 
$$\P\{S_n\in  J_k(n), S_m\in  J_\ell(m) \}-\P\{S_n\in  J_k(n)\}\P\{ S_m\in  J_\ell(m) \}
$$$$=\sum_{q\in J_k(n)\atop p\in J_\ell (m)}
\P \{S_m =p\}\big(\P\{ S_{n-m} =q -p\}- \P \{S_n =q \}\big) $$

On using  Theorem \ref{radem.t2.1}, 
  \begin{eqnarray*} \Big|  \P\{S_n=k \}-{1\over  \sqrt{ 2\pi n}} 
e^{-  {k ^2\over  2 n } }\Big|&\le & { C \over   n^{  ( 1+\d )/2}}
 \cr  \Big|  \P\{S_{n-m}=k  -\ell\}-{1\over  \sqrt{ 2\pi(n-m) }} 
e^{-  {(k  -\ell)^2\over  2 (n-m) } }\Big|&\le & { C \over   (n-m) ^{  ( 1+\d )/2}}.
\end{eqnarray*}
Hence
$$\Big|\Big(\P\{ S_{n-m} =q  -p\}- \P \{S_n =q  \}\Big)-\Big({1\over  \sqrt{ 2\pi(n-m) }} 
e^{-  {(q  -p)^2\over  2 (n-m) } }- {1\over  \sqrt{ 2\pi n}} 
e^{-  {q ^2\over  2 n } }\Big)\Big|$$
$$\le     C\Big({1
\over   n ^{  ( 1+\d )/2}}+  { 1\over   (n-m) ^{  ( 1+\d )/2}}\Big) $$
  Further  
  $$q-p\ge 2^{k-1}\sqrt n-2^{\ell -1}\sqrt m\ge \sqrt n(2^{k-1} -2^{\ell -1})\ge \sqrt n 2^{k-2}\ge \sqrt{ n-m}\, 2^{k-2}.$$ Thus 
$${1\over  \sqrt{ 2\pi(n-m) }} 
e^{-  {(q  -p)^2\over  2 (n-m) } }\le {C\over  \sqrt{  n-m  }} 
e^{-  2^{2k-5} }.$$
And 
$${1\over  \sqrt{ 2\pi n}} 
e^{-  {q ^2\over  2 n } }\le {C\over  \sqrt{   n}} 
e^{- 2^{2k-3} }.$$
Therefore
$$\Big| \P\{ S_{n-m} =q  -p\}- \P \{S_n =q  \}\Big|\le {C\over  \sqrt{  n-m  }} 
e^{-  2^{2k-5} }+       { C\over   (n-m) ^{  ( 1+\d )/2}}. $$

We deduce that 
\begin{align}&\Big|\P\{S_n\in  J_k(n),  S_m\in  J_\ell(m) \}-\P\{S_n\in  J_k(n)\}\P\{ S_m\in  J_\ell(m) \}\Big|
\cr \le &\ C\sum_{q\in J_k(n)\atop p\in J_\ell (m)}
\P \{S_m =p\}\Big| {1\over  \sqrt{  n-m  }} 
e^{-  2^{2k-5} }+       { 1\over   (n-m) ^{  ( 1+\d )/2}}\Big| 
\cr \le & \ C\sum_{  p\in J_\ell (m)}
\P \{S_m =p\}\Big({2^k \sqrt n\over \log (2^k \sqrt n)}\Big)\Big| {1\over  \sqrt{  n-m  }} 
e^{-  2^{2k-5} }+       { 1\over   (n-m) ^{  ( 1+\d )/2}}\Big| .
\end{align}
And 
\begin{align}&\sum_{k>\ell +1\atop 
2^k\le \sqrt n}\Big|\P\{S_n\in  J_k(n), S_m\in  J_\ell(m) \}-\P\{S_n\in  J_k(n)\}\P\{ S_m\in  J_\ell(m) \}\Big|
\cr \le & \  C\sum_{  p\in J_\ell (m)}
\P \{S_m =p\}\Big(\sqrt{{n\over     n-m  }}\sum_{k>\ell +1\atop 
2^k\le \sqrt n}{2^k  e^{-  2^{2k-5} }\over \log (2^k \sqrt n)}   
 \cr & \qq\qq\qq\qq \qq\qq\qq+  { 1\over   (n-m) ^{  ( 1+\d )/2}} \sum_{k>\ell +1\atop 
2^k\le \sqrt n} \#\{ J_k(n)\} \Big)
\cr \le & \  C\sum_{  p\in J_\ell (m)}
\P \{S_m =p\}\Big(\sqrt{{n\over     n-m  }} { 1\over \log  n }     e^{-B  2^{2\ell } }
 +  { n\over   (n-m) ^{  ( 1+\d )/2}\log n} \Big).  
 \end{align}
Hence
\begin{align}\sum_{\ell}&\sum_{k>\ell +1\atop 
2^k\le \sqrt n}\Big|\P\{S_n\in  J_k(n), S_m\in  J_\ell(m) \}-\P\{S_n\in  J_k(n)\}\P\{ S_m\in  J_\ell(m) \}\Big|
\cr & \le \ C {\sqrt n\over   \sqrt{ n-m } \log  n }\sum_{\ell} \sum_{  p\in J_\ell (m)}
\P \{S_m =p\}     e^{-B  2^{2\ell } }
  \cr \ & \qq+  { n\over   (n-m) ^{  ( 1+\d )/2}\log n} \sum_{\ell}\sum_{  p\in J_\ell (m)}
\P \{S_m =p\}\Big)\cr \le & \ C {\sqrt n\over    \sqrt{ n-m }   \log  n \log m } 
  +  { n\over   (n-m) ^{  ( 1+\d )/2}\log n \log m} \le   {C\over       \log  n \log m } .
      \end{align}
 
 \bigskip\par
  










   Let $X  $  be  a square integrable  random variable
   %  with basic probability space $(\O, \A, \P)$,
      taking values in the  lattice $\mathcal L(v_0,1)=\big\{v_k=v_0+ k,k\in \Z\big\}$, and let $f(k)=\P\{X=v_k\}$, for each $k$.   
   \vskip 4 pt 
Assume that 
   \beq \label{bp}\t_X   =\sum_{k\in \Z}f(k)\wedge f(k+1)>0,
   \eeq  
where $a\wedge b=\min(a,b)$.  Let $  X_i$, $i\ge 1 $  be  independent, identically distributed random variables having same law than $X$, and let $S_n=\sum_{j=1}^nX_j$, for each $n$.    Let $ \bar \m= \{ \m_k, k\in \Z\}$ be a   sequence of non-negative reals such that  $0<\m_k < f(k) $, if $f(k)>0$,   $\m_k=0$ if $f(k)=0$,    let  $ \m= \sum_{k\in \Z}\m_k$, and assume that $1- \m< \t_X$.   
\vskip 1 pt 
  Let also non-negative reals $\bar \tau=  \{ \tau_k, k\in \Z\}$, which are  solutions of the (solvable) equation:   
 $ {\tau_{k-1}+\tau_k\over 2} =f(k)  - \m_k,
 $ for all $k \in \Z$.  
Let   $\t = \sum_{k\in \Z}  \tau_k = 1- \m.$ Further let $s(t) =\sum_{k\in \Z} \m_k\, e^{ 2i \pi   v_kt}$ and    $\rho$ be such  that  $1-\t<\rho<1$.   
        Let $\widetilde X$ be the   random variable associated  to $X$  and $\bar \m$, and defined by the relation
 $ \P\{ \widetilde X =v_k\}= {\tau_k}/{\t}$, for all  $k\in \Z   $.  
\vskip 2 pt  


\subsection{The Bernoulli part   of a random variable.}\label{sub2.1}
Here we describe an ingenious coupling method, which allows one to transfer results true for the Bernoulli random walk to general square integrable random walks. Let $Y$ be
  a random variable   such that  ${\mathbb P}\{Y
\in\mathcal L(v_0,D)\}=1$. We do not   assume  that the  span $D$ is maximal.  
  Introduce the following characteristic,  
\begin{eqnarray}\label{vartheta}  \t_Y =\sum_{k\in \Z}{\mathbb P}\{Y=v_k\}\wedge{\mathbb P}\{Y=v_{k+1}\} ,
\end{eqnarray}
 where $a\wedge b=\min(a,b)$.
Obviously $\t_Y$   depends on $D$.
 Note    that we always have the relation
   \begin{eqnarray}\label{vartheta1}
0\le \t_Y<1 .\end{eqnarray}
When $Y$ has finite  variance
$\sigma^2$,   we have the following important liaison inequality,
 \begin{eqnarray}  \label{tD}\s^2 \,\ge\,  \frac{ D^2   }{4}\, \t_Y,
\end{eqnarray}
from which it follows that
 \begin{eqnarray}  \label{tDn}\qq \qq \qq \qq\  \s_n^2\,\ge  \,\frac{ D^2   }{4}\ \sum_{j=1}^n \t_{X_j}, \qq \quad n\ge 1.
\end{eqnarray}    

  Let $X$ be
  a random variable   such that  ${\mathbb P}\{X
\in\mathcal L(v_0,D)\}=1$. 
We assume  that 
\begin{equation}\label{varthteta2}\t_X>0.
\end{equation}
Let $ f(k)= {\mathbb P}\{X= v_k\}$, $k\in \Z$. Let also $0<\t\le\t_X$. One can associate to $\t$ and $X$  a
sequence  \begin{proof}[Proof]
 Plainly,
\begin{eqnarray*}{\mathbb P}\{Z=v_k\}&=&{\mathbb P}\big\{ V+\e DL=v_k, \e=1\}+ {\mathbb P}\big\{ V+\e DL=v_k, \e=0\} \cr
&=&{{\mathbb P}\{ V=v_{k-1}, \e=1\}+{\mathbb P}\{
V=v_k, \e=1\}\over 2} +{\mathbb P}\{ V=v_k, \e=0\}
\cr&=& {\tau_{k-1}+ \tau_{k }\over 2} +f(k)-{\tau_{k-1}+ \tau_{k
}\over 2}
= f(k).
\end{eqnarray*}
\end{proof} Consider now independent random variables  $ X_j,j=1,\ldots,n$,    and assume that
\begin{equation}\label{basber.j}  \t_{X_j}>0, \qq \quad  j=1,\ldots, n.
\end{equation}
Let $0<\t_j\le \t_{X_i}$, $j=1,\ldots, n$.  
 One can  associate to them a
sequence of independent vectors $ (V_j,\e_j, L_j) $,   $j=1,\ldots,n$  such that
 \begin{eqnarray}\label{dec0} \big\{V_j+\e_jD   L_j,j=1,\ldots,n\big\}&\buildrel{\mathcal D}\over{ =}&\big\{X_j, j=1,\ldots,n\big\}  .
\end{eqnarray}

Further the sequences $\{(V_j,\e_j),j=1,\ldots,n\}
 $   and $\{L_j, j=1,\ldots,n\}$ are independent.
For each $j=1,\ldots,n$, the law of $(V_j,\e_j)$ is defined according to (\ref{ve}) with $\t=\t_j$.  And $\{L_j, j=1,\ldots,n\}$ is  a sequence  of
independent Bernoulli random variables. Set
 \begin{equation}\label{dec} S_n =\sum_{j=1}^n X_j, \qq  W_n =\sum_{j=1}^n V_j,\qq M_n=\sum_{j=1}^n  \e_jL_j,  \quad B_n=\sum_{j=1}^n
 \e_j .
\end{equation}


\begin{proposition}\label{lmd}We have\begin{eqnarray*} \{S_k, 1\le k\le n\}&\buildrel{\mathcal D}\over{ =}&  \{ W_k  +  DM_k, 1\le k\le n\} .
\end{eqnarray*}
And  $M_n\buildrel{\mathcal D}\over{ =}\sum_{j=1}^{B_n } L_j$.
 \end{proposition}

We also will need the following local limit theorem for Bernoulli sums.
\begin{proposition} \label{lltber}Let
$\mathcal B_n=\b_1+\ldots+\b_n$, $n=1,2,\ldots$  where
$
\b_i
$ are i.i.d.  Bernoulli r.v.'s (\,${\mathbb P}\{\b_i=0\}={\mathbb P}\{\b_i=1\}=1/2$).  There exists a numerical constant $C_0$ such that for all positive $n$
 \begin{eqnarray*}  \sup_{k}\, \Big|  {\mathbb P}\big\{\mathcal B_n=k\}
 -\sqrt{\frac{2}{\pi n}} e^{-{ (2k-n)^2\over
2 n}}\Big|\ \le \
  \  \frac{C_0}{n^{3/2}}  .
  \end{eqnarray*}
\end{proposition}
Theorem \ref{estPdlBn.u} is combined in \cite{W9} with the Bernoulli part method,    to prove a   uniform semi-local limit theorem along    sets of multiples   for sums of i.i.d.  
   random variables.   More precisely, 
    
    %%%%%%%%%
%%%%%%%%%
\vskip 30 pt
%\section{ }
%%%%%%%%%
%%%%%%%%%
\begin{remark}\rm
 First observe
% \footnote{This remark was brought to me by Istv\'an Berkes}  
that there exists an absolute constant  $A>0$  such that for any positive integer $n$,
 $\E   \exp\{  A{(B_n-{n/2})^2/ n } \}\le 2$. 
Next, we easily have   
\begin{align}\P\{ B_n\
\hbox{ prime}\}&\le \P\Big\{ B_n\
\hbox{ prime},\, |B_n-{n\over 2}|\le \sqrt{n\log \log n}\Big\} \cr & \quad +\P\Big\{ |{B_n-{n\over 2}\over \sqrt{n}}|>  \sqrt{\log 
\log n}\Big\}\cr &\le  \P\Big\{ B_n\
\hbox{ prime},\,  |B_n-{n\over 2}|\le \sqrt{n\log  \log n}\Big\}
\cr &\quad +e^{-A\log  \log n} \ \E e^{  A{(B_n-{n/2})^2\over n } }  \cr
&\le \P\Big\{ B_n\
\hbox{ prime},\, |B_n-{n\over 2}|\le \sqrt{n\log\log n}\Big\}+ {2\over (\log n)^{ A}}  .
\end{align}

Suppose now there exists a real $A'<A$, such that for any  $n$ large enough, $\P\big\{ B_n\
\hbox{ prime}\big\}\ge (\log n)^{-A'}$.  Then
$$ \P\big\{ B_n\
\hbox{ prime and}\ |B_n-({n/ 2})|\le \sqrt{n \log\log n}\big\}>0$$
provided that   $n$ is sufficiently large. From there, we draw immediately that there exists a prime number   $p$ such that
$${n\over 2} - \sqrt{n \log\log n}\le   p\le  {n\over 2} + \sqrt{n \log\log n} $$
Let  
 $2=p_1<p_2<\ldots  $ be the sequence of consecutive prime numbers (${\mathcal P}$ being their set). Recall that   $p_k\sim k\log
k$,  
$k\to \infty$. Choose $n=n_k$ so that 
${n_k\over 2} -
\sqrt{n_k \log\log n_k}\sim p_k$. 

\vskip 3 pt  We deduce that there exists a prime number between  $p_k$ and $p_k+ 2\sqrt{n_k \log\log n_k}\sim
p_k+2\sqrt{k\log k
\log\log  k} $, in other words   $g(k)=p_{k+1}-p_k$  verifies
$$g(k)= {\mathcal O}\big( \sqrt{k\log k \log\log 
k} \big)  .$$
But according to a well-known result of   Cram\'er \cite{C},  the validity of the Riemann Hypothesis would  imply that 
 $$g(k)= {\mathcal O}\big(  k^{1/2}   
 (\log k)^{3/2}\big) .$$
Besides, the Cram\'er conjecture states that  $$\limsup_{k\to \infty}{g(k)\over (\log p_k)^2}=1 .
$$
\end{remark}


We have not addressed in the present paper the study of the system 
  \begin{equation} \label{system.Delta}{\bf \D} :=\big\{ \chi\{d|B_n\},    \  n\ge 1, d\ge 1\big\},
  \end{equation}
 where we set for $n\ge 2$ and $d\ge 1$ integers,
    \begin{equation}\label{chi} \chi(d| n)=\begin{cases} 1 &\quad \hbox{if $d|n$,}\cr
0 &\quad \hbox{if $d\not|n$} .\end{cases}
\end{equation} 
An important fact is that: 
  \begin{equation}\label{mixing}
 \hbox{\it The system ${\bf \D} 
 % =\big\{ \chi\{d|B_n\},    \  n\ge 1, d\ge 1\big\}
 $   
  is mixing.} 
\end{equation}
   That is,   the correlation function 
  \begin{equation}\label{Delta.def}  {{ \Delta}}  \big( (d,n),  (\d, m)\big)   \,=\, \P\big\{ d|B_n\, ,\,   \d|B_m\big\}-\P \{ d|B_n \}\P \{    \d|B_m \},
 \end{equation} 
verifies for each     $d$ and $\d$, 
\begin{equation}\label{corr1a1}
   \lim_{m-n\to \infty }  {{ \Delta}}  \big( (d,n),  (\d, m)\big)\,=\, 0. 
\end{equation}

Therefore   the classical   convergence criteria for sequences of random variables  directly apply to the system ${\bf \D}$, which is   a considerable advantage of this model. These criteria not only allow us to  prove almost sure limit theorems, but also provide a control of the speed of convergence, which is often very sharp.
\vskip 3 pt We thoroughly  studied  the associated correlation function 
 ${ \Delta}}  \big( (d,n),  (\d, m)\big)$
 and could identify three zones of  dependence, weak independence and strong independence, 
   according to the cases   $n<m\le n+n^c$,   $n+ n^{c} \le m\le 2n$ and   $m\ge 2n$, where  $0<c<1$. Corresponding correlation estimates are also established. Often the proofs involve long specific but unavoidable technicalities. These results are refreshed and completed in the monograph \cite{W4}. Our belief is that 
   %much remains to do. This way, the 
   study of the second order theory of a random model for divisors, should be thoroughly prospected.% It   can be calqued on the work done in the Bernoulli case,  and    more developed. However it seems to have been  also overlooked,
   %, as well our papers on divisors, 
  % as it  does not appear in particular in \cite{FF}.
  
  
  

\subsection{A  way back.}\label{sub.bin.way.back}
% along Pascal's matrices.}  
Let $f(k)$ be an arithmetical function. Consider the column vector ${\boldsymbol \nu}=\big(f(1),f(2),\ldots, f(n)\big)$. Let $L_n$ be the $n$-size   lower triangular Pascal matrix:
\begin{eqnarray*}L_n={ \left[\begin{matrix}
1        \cr
1   &1       \cr
1        &2            &1    \cr
1       &3       &3 &1       \cr
 \vdots         &   \vdots                  &\vdots   &\vdots   &\ddots    
 \cr {{n}\choose{1}}& {{n}\choose{2}}&{{n}\choose{3}}&{{n}\choose{4}}  &\cdots  &{{n}\choose{n}}\end{matrix} \right]
 %\qq j=1,\ldots, n.
\end{eqnarray*}
Then $L_n{\boldsymbol \nu}=  \,{\boldsymbol w}$ where ${\boldsymbol w}=\big(\E f(B_1),2\E f(B_2),\ldots, 2^n\E f(B_n)\big)$. Call and Velleman \cite{CV} identified $L^{-1}_n$, which is $DL_nD^{-1}$ with $D$ as indicated:
\begin{eqnarray*}L^{-1}_n={ \left[\begin{matrix}
1        \cr
  &-1       \cr
1        &-2            &1    \cr
-1       &3       &-3 &1       \cr
 \vdots         &   \vdots                  &\vdots   &\vdots   &\ddots    
 \cr {{n}\choose{1}}& {{n}\choose{2}}&{{n}\choose{3}}&{{n}\choose{4}}  &\cdots  &{{n}\choose{n}}\end{matrix} \right]
 %\qq j=1,\ldots, n.
\end{eqnarray*}
\begin{eqnarray*}D={ \left[\begin{matrix}
1        \cr
  &-1       \cr
       &{}      &1    \cr
  &{}          &{}         &-1       \cr
 &{}          &{}   &{}          &{}  &\ddots    
 \cr &{}          &{}   &{}          &{}    &   &(-1)^{n+1}\end{matrix} \right]
 %\qq j=1,\ldots, n.
\end{eqnarray*}So that knowing ${\boldsymbol w}$, we have a way back to know ${\boldsymbol \nu}$! For instance,
\beq\label{} \s_{-1}(n)=\sum_{j=1}^n  (-1)^{j+1}2^j{{n}\choose{j}}\E \s_{-1}(B_j).
\eeq

  
  
  
  
  
   %%%%%%%%%%%%%%%%%%%%%%%%%%%%%%%%%%%%%%%%%%%%%%%%%%%%%%%%%%%%%%%%%%%%%%%%%%%%%%%%%%%%%%%%%%%%%%%%%%%%%%%%%%%%%%%%
 
 
 \section{Mixing properties.} \label{s7}
    The  study the asymptotic behavior of the system 
 $${\bf \Delta} =\big\{ \chi\{d|B_n\},    \    2\le d\le n,\ n\ge 1\big\}$$
  as $n$ tends to infinity, is made in \cite{W4}   and   
  \cite{W14}. 
Introduce    the   correlation
function  
 \begin{equation*}
% \label{delta.def}
  {\Delta} \big( (d,n), (\d, m)\big)= \P\big\{ d|B_n  ,    \d|B_m\big\}-\P \{ d|B_n \}\P \{    
\d|B_m \}.
\end{equation*}
 \begin{lemma}\label{lem61} For all integers $m\ge n\ge 1$, $d\le n$, $\d\le m$, 
\begin{eqnarray*}\label{6.3}\P\big\{ d|B_n\, ,\,   \d|B_m\big\} &=&
{1\over d\d} 
\sum_{j=1}^{d-1}\sum_{h=1}^{\d-1}  e^{ i\pi  (  {j\over d}n +   {h\over
\d}m) }\cos^n{ \pi  (  {j\over d} +   {h\over \d})   }   \cos^{m-n}{ \pi   {h\over
\d} }
%\cr & &  +{\P \{ \d|B_m \}\over   d}+{\P \{ d|B_n \}\over  \d}   -{1\over d\d}.
 \cr & & +  \P\{\d|B_n\}\P\{\d|B_m\}- \big(\P
\{ d|B_n
\} -  \frac{1}{ d}\big)\big(\P
\{
\d|B_m
\} -  \frac{1}{ \d}\big).
\end{eqnarray*}  
Also,
\begin{eqnarray*}\label{6.4}& & {\bf \Delta} \big( (d,n),  (\d, m)\big)  =
\cr & &\qq \quad{1\over d\d}  \sum_{j=1}^{d-1}\sum_{h=1}^{\d-1}  e^{ i\pi  (  {j\over d}n +   {h\over \d}m) }  \cos^{m-n}{ \pi   {h\over
\d} }\, \Big(\cos^n{ \pi  (  {j\over d} +  
{h\over \d})   }-\cos^{n} {
\pi j\over d}     \cos^{n} {
\pi h\over \d}\Big).
 \end{eqnarray*} 
\end{lemma}
  %\vskip 6pt  In the example below, we show    that   ${\bf
%\Delta} \big( (d,n), (\d, m)\big)$  can be  directly depending  on the probability  $\P\big\{ d\d|B_n  B_m\big\}$.
%\begin{example} Let  $m=n+1$, $d=p$, $\d=q$ where $p,q$ are two  different prime numbers.
%  Then  
%\begin{eqnarray*}{\boldsymbol \tau}_ {p,q}(n)&:= & {\bf
%\Delta} \big( (p,n), (q, n+1)\big)+ {\bf
%\Delta} \big( (q,n), (p, n+1)\big)\cr &= &
%\P\{pq|B_nB_{n+1}\} - \frac{1}{ 2}\P\{pq| B_n+1  \}
%-\P\{p |B_n\} \P\{q |B_{n+1}\}-\P\{q |B_n\} \P\{p |B_{n+1}\}. 
%\end{eqnarray*}
% Indeed,   first note that ${\boldsymbol \tau}_ {p,q}(n)$ equals to
 %    $$ \P\{p |B_n , q |B_{n+1}\} -
%\P\{p |B_n\} \P\{q |B_{n+1}\}+ \P\{q |B_n , p |B_{n+1}\}-\P\{q |B_n\} \P\{p |B_{n+1}\}.$$
%As $\P\{p |B_n , q |B_{n+1}\} ={1\over 2} \P\{pq |B_n  \}+{1\over 2}\P\{p |B_n , q |B_ n+1 \}$,
 % we have
%\begin{eqnarray*}\P\{p |B_n , q |B_{n+1}\}+ \P\{q |B_n , p |B_{n+1}\}&=&\P\{pq |B_n  \}
% +{1\over 2}\big(\P\{p |B_n , q |B_ n+1 \}\cr & &  +\ \P\{q |B_n , p |B_ n+1
%\}\big). 
%\end{eqnarray*}
%But
%$$pq|B_n (B_{n }+1)\quad \Longleftrightarrow \quad   p |B_n , q |B_ n+1 \ \  {\rm or}\ \  q |B_n , p |B_ n+1 \ \ {\rm or}\ \
%pq|B_n    
% \ \ {\rm or}\ \ pq| B_n+1 .
%   $$
%These four events being pairwise disjoint, $\P\{pq|B_n (B_{n }+1)\} $ is the sum of their probabilities.
%  By reporting, we get
%  $$\P\{p |B_n , q |B_{n+1}\}+ \P\{q |B_n , p |B_{n+1}\}=   \P\{pq|B_nB_{n+1}\} - {1\over 2}\P\{pq|
%B_n+1  \}  ,$$   
%as claimed. 
%\end{example}
\begin{proof} 
  By using again
        formula  
 $u\chi({u |B_{n }})=\sum_{j=0}^{u-1} e^{2i\pi  {j \over u}B_{n }}$,  
 next integrating, we get
\begin{eqnarray*}\P\big\{ d|B_n\, ,\,   \d|B_m\big\}    & =&{1\over d\d}\E \Big\{
\sum_{j=0}^{d-1}\sum_{h=0}^{\d-1} e^{2i\pi  (  {j\over d}B_n+   {h\over
\d}B_m )}\Big\} 
\cr &=& {1\over d\d}  \sum_{j=0}^{d-1}\sum_{h=0}^{\d-1} \E e^{2i\pi  (  {j\over d} +   {h\over \d})B_n } \E e^{2i\pi   {h\over
\d} B_{m-n} }
\cr &=& {1\over d\d}  \sum_{j=0}^{d-1}\sum_{h=0}^{\d-1}  e^{ i\pi  (  {j\over d} +   {h\over \d}) n }\cos^n{ \pi  (  {j\over d} +  
{h\over \d})   }  e^{ i\pi   {h\over
\d} (m-n) }\cos^{m-n}{ \pi   {h\over
\d}  }
\cr &=&{1\over d\d}  \sum_{j=0}^{d-1}\sum_{h=0}^{\d-1}  e^{ i\pi  (  {j\over d}n +   {h\over \d}m) }\cos^n{ \pi  (  {j\over d} +  
{h\over \d})   }   \cos^{m-n}{ \pi   {h\over
\d} }
\cr &=& {1\over d\d}  \sum_{j=1}^{d-1}\sum_{h=1}^{\d-1}  e^{ i\pi  (  {j\over d}n +   {h\over \d}m) }\cos^n{ \pi  (  {j\over d} +  
{h\over \d})   }   \cos^{m-n}{ \pi   {h\over
\d} } \cr & &\ \  +{1\over d\d}\Big(   \sum_{h=0}^{\d-1}  e^{ i\pi      {h\over \d}m  }   \cos^{m }{ \pi   {h\over
\d} }
 +   \sum_{j=1}^{d-1}   e^{ i\pi     {j\over d}n   }\cos^n{ \pi    {j\over d}   
   }  \Big)
\cr &=& {1\over d\d}  \sum_{j=1}^{d-1}\sum_{h=1}^{\d-1}  e^{ i\pi  (  {j\over d}n +   {h\over \d}m) }\cos^n{ \pi  (  {j\over d} +  
{h\over \d})   }   \cos^{m-n}{ \pi   {h\over
\d} } 
\cr & &\ \   +{\P \{ \d|B_m \}\over   d}
 +{\P \{ d|B_n \}\over  \d}   -{1\over d\d}  .
\end{eqnarray*}
 This can also be rewritten as 
 \begin{eqnarray}\label{inter}\P\big\{ d|B_n\, ,\,   \d|B_m\big\}    & =& 
 {1\over d\d}  \sum_{j=1}^{d-1}\sum_{h=1}^{\d-1}  e^{ i\pi  (  {j\over d}n +   {h\over \d}m) }\cos^n{ \pi  (  {j\over d} +  
{h\over \d})   }   \cos^{m-n}{ \pi   {h\over
\d} } 
 \cr & & +  \P\{\d|B_n\}\P\{\d|B_m\}- \big(\P
\{ d|B_n
\} -  \frac{1}{ d}\big)\big(\P
\{
\d|B_m
\} -  \frac{1}{ \d}\big).
\end{eqnarray} By    substracting  and using (\ref{Lemma41}), this factorizes well and we get a more condensed formulation of ${\bf \Delta} \big( (d,n),  (\d, m)\big)$, namely, 
\begin{align*} {\bf \Delta} \big( (d,n),&  (\d, m)\big)
\cr =&   {1\over d\d}  \sum_{j=1}^{d-1}\sum_{h=1}^{\d-1}  e^{ i\pi  (  {j\over d}n +   {h\over \d}m) }  \cos^{m-n}{ \pi   {h\over
\d} }\Big\{\cos^n{ \pi  (  {j\over d} +  
{h\over \d})   }-\cos^{n} {
\pi j\over d}     \cos^{n} {
\pi h\over \d}\Big\}. 
\end{align*}
\end{proof} 
An important consequence is: 
 \begin{theorem} \begin{equation}\label{mixing}
 \hbox{\it The system ${\bf \D} 
 % =\big\{ \chi\{d|B_n\},    \  n\ge 1, d\ge 1\big\}
 $   
  is mixing,}  
\end{equation} namely,    the correlation function verifies for each     $d$ and $\d$, 
\begin{equation}\label{corr1a1}
   \lim_{m-n\to \infty }  {{ \Delta}}  \big( (d,n),  (\d, m)\big)\,=\, 0. 
 \end{equation}
\end{theorem}
\begin{proof}   Let    $d$ and $\d$ be  fixed. Then for each $1\le h<\d$, $\cos^{m-n}{ \pi   {h\over
\d} }\to 0$   when $m-n \to \infty$, and thereby
$$\lim_{m-n\to \infty }{\bf \Delta} \big( (d,n),
(\d, m)\big)= 0.$$ 
\end{proof}
It follows from  the study of second order properties of  this system made in \cite{W4}   and   
  \cite{W14}, that the following   zones   of  weak and strong independence can be  identified, according to three domains of values:
 %\vskip 10 pt
\begin{eqnarray*}{\rm (i)}& &\hbox{{ Weakly independent case}: $n+ n^{c} \le m\le 2n$,  $0<c<1$.}
\cr {\rm (ii)} &&\hbox{{ Strongly independent  case}:  $m\ge 2n$.}
 \cr  {\rm (iii)} && \hbox{{ Dependent case}: 
$n<m\le n+n^c$, $c>0$ small.}
\end{eqnarray*}  
 \vskip 10 pt
%--- {\ph Weakly independent case}: $n+ n^{c} \le m\le 2n$, where $0<c<1$.
  %, where a uniform rate of independence   over $d,\d$ such that $\max(d,\d)\le C \sqrt{{ n/ \log n}}$ will   be proved.  

%\vskip 6 pt
  %--- {\ph Strongly independent  case}:  $m\ge 2n$.

%\vskip 6pt   --- {\ph Dependent case}: 
%$n<m\le n+n^c$, $c>0$.
%, and we  will start with the simple  bound   
 %    $${\bf\Delta} \big( (d,n), (\d, m)\big)  \le    \P\big\{ d|B_n\, ,\,   \d|B_m\big\}  \le    \P\big\{
%d\d|B_n  B_m\big\} . $$
%\begin{equation} \label{dom1}  n+ n^{c} \le m\le 2n ,\end{equation}
%where   $0<c <1$.  

   \begin{theorem} \label{t5} % Let $0<c <1$. There exist constants $n_0$ and $C$ such that for any$  n\ge n_0$ and    $      n+ n^{c} \le m\le 2n   $
{\rm (1)} Let $\a>\a'>3/2$. There exist $n_0$ and a constant $C $ depending on $\a'$ only, such that for $m \ge n+n_0$,    any $1\le d\le n$ and $\d$ such that $\d  <{\pi \over2\sqrt{2\a}   } ( {m-n\over \log (m-n)   } )^{1/2}$, we have
%$\max(d, \d)<  {\pi \over2\sqrt{2\a}   } ( {m-n\over \log (m-n)   } )^{1/2}$. Then    
$$\big|{\bf \Delta} \big( (d,n), (\d, m)\big)\big| \le  C \,  (m-n)^{-\a'} .$$ 

{\rm (2)} For any positive integer $p$, there exist    constants $C_p $ and $n_p$ depending on $p $ only such that for $n\ge n_p$ and   if $ c> {1\over  p +1}$, $m$ is such that $  n+n^c\le m\le 2n$, we have
\begin{eqnarray*}
 \big|{\bf \Delta} \big( (d,n), (\d, m)\big)\big|
 & \le & C_p\, \bigg\{ {  (c  \log n)^{  p+1 }\over n^{ c(p+1)-1}}  +{1\over d\d}  \sum_{0< 
 {\pi j\over  d}\le 2 ( {2\a\log (m-n) \over m-n} )^{1/2}\atop 0< {\pi h\over \d}\le  ( {2\a\log
(m-n) \over m-n} )^{1/2}}     e^{-   {n\over 3}  (  {j\over d} - {h\over \d})^2    }  \bigg\}
\end{eqnarray*}



{\rm (3)} %\label{sicDelta.th}
 %Let $   \d_0 >\sqrt 2$. 
 For any  positive integer $p$, there exist constants $C_p$ and $n_p$   such
that for
 $n\ge n_p$, and 
    $    m\ge  2n $,  
\begin{eqnarray*}
 %\label{former(6.26)}
 \big|{\bf \Delta} \big( (d,n), (\d, m)\big)\big| \, \le \, C_p\,   \bigg\{    
  \, \frac{(\log (m-n))^{p+1/2}}{(m-n)^{p-1/2}}   
 +{1\over d\d} \sum_{1\le j<{d\over 2}}\  \sum_{ 0< {\pi h\over \d}\le  \p_{m-n}} 
    e^{-   {n\over 3} (   {j\over d} -
 {h\over \d})^2   }      \bigg\}
. 
\end{eqnarray*}
% $$\sup_{2\le d\le n \atop \d<{\pi\over \d_0\sqrt{2\a}} \sqrt{m\%\big( (d,n), (\d, m)\big)\big| \le C_p  \, 
%n\,\Big({  \log n\over n}\Big)^{p+1 /2}\, \Big({ \log m\over
%m }\Big)^{1/2}.$$
\end{theorem}
A technical step prior to the proof consists with estimating in the    above defined domains the quantity   
 \begin{eqnarray}\label{psi(d,n),(delta,m)}
 {\bf \Psi}\big( (d,n), (\d, m)\big)&=& {1\over d\d}  \sum_{j=1}^{d-1}\sum_{h=1}^{\d-1}  e^{ i\pi  ( 
{j\over d}n +   {h\over
\d}m) }\cos^n{ \pi  (  {j\over d} +   {h\over \d})   }   \cos^{m-n}{ \pi   {h\over
\d} }  .
% \cr  {\bf \Phi}\big( (d,n), (\d, m)\big) &=& \P\{\d|B_n\}\P\{\d|B_m\}- \big(\P
%\{ d|B_n\} -  \frac{1}{ d}\big)\big(\P\{\d|B_m\} -  \frac{1}{ \d}\big).
 \end{eqnarray}
% By (\ref{inter}),   $$\P\big\{ d|B_n\, ,\,   \d|B_m\big\}
%={\bf \Psi}\big( (d,n), (\d, m)\big)+{\bf \Phi}\big( (d,n), (\d, m)\big)  . $$
 %which after reduction due to symmetries can be rewritten as 
% \begin{align*} {\bf \Psi}\big( (d,n),& (\d, m)\big) 
%\cr &= {1\over d\d}  \sum_{1\le |h|<\d/2}   \sum_{1\le |j|<d/2}  e^{ i\pi  (  {j\over d}n +   {h\over
%\d}m) }  \cos^{m-n}{
% \pi   {h\over
% \d} } \cos^n{ \pi  (  {j\over d} +  
% {h\over \d})   }+{\bf r},  
 % \end{align*}
%  where
%\begin{equation}\label{6.5} \qq  {\bf r} =
    %\begin{cases}   
   %  0    &   \ \ \hbox{\rm if $d$ is odd,} 
%\cr   \displaystyle{{1\over d\d}   \sum_{1\le |h|<\d/2}} e^{ i\pi  (  {n\over 2}  +   {hm\over \d} ) }  \cos^{m-n}{     {\pi h\over
%\d} } .\cos^n{    (  {\pi \over 2} +  
%{\pi h\over \d})   }  &     \ \ \hbox{\rm if $d$ is even.}
% \end{cases} 
%\end{equation}
  

This  partly  depend ($m,d,\d$ being fixed)
 on the key sums below     \begin{eqnarray} \label{Fxy}
  F( x,y)=  \sum_{1\le j\le xd\,,\, 1\le h \le y\d}     e^{-  m  (  
{j\over d} -
 {h\over \d})^2    } ,
 \end{eqnarray}
where  $0<x,y<1$ are small. 
%\vskip 2 pt
%Assume that $(d,\d)=1$. Then the numbers $j\d-hd$ are all distinct. Therefore
% \begin{eqnarray} \label{Fxy}
%  F( x,y)=  \sum_{1\le j\le xd\,,\, 1\le h \le y\d}     e^{-  m  (  
%{j\over d} -
 %{h\over \d})^2    } \le C \,\frac{d\d}{\sqrt m}.
 %\end{eqnarray}
 %  \vskip 3 pt 
%Fix also some reals $\a,\a'  $,  depending on  $c$, such that\begin{eqnarray} \label{wic.a.a'}\a >\a' >\max(3/2, 1/c). 
%\end{eqnarray} This point is technical. 
 
 
 Therefore   the classical   convergence criteria for sequences of random variables  directly apply to the system ${\bf \D}$, which is   a considerable advantage of this model. These criteria not only allow us to  prove almost sure limit theorems, but also provide a control of the speed of convergence, which is often very sharp.
%  These results from  \cite{W4} are refreshed and completed in the monograph  \cite{W14}. 
 
 
 
 
 To conclude we add a few remarks on
\vskip 3 pt
  \noi{\tf Random walks in Analytic  Number Theory.}   This topic is the object of the recent work Weber \cite{W15}. Apart from the Cr\'amer model, we give two examples, among others    presented in the above.
 \vskip 3 pt
Let $X_1,X_2,\ldots$ denote an
infinite
sequence of independent Cauchy distributed random variables
(with characteristic function $\varphi(t) =e^{-|t|}$),  and $S_n= X_1+ \ldots +X_n$,  $n\ge 1$. In Lifshits and Weber \cite{LW}, in order to understand the behavior of $\zeta({1\over 2} + it)$ when $t$ tends
to infinity,   the  almost sure asymptotic behavior
of the system
$$  \zeta_n:=\zeta({1\over 2} + iS_n), \quad\qq n=1,2, \ldots  
$$
is investigated.  Put for any positive integer $n$
$${\mathcal Z}_n = {\zeta}({1/  2}+iS_n) -\E\,  {\zeta}({1/ 2}+iS_n)
=\zeta_n-\E \zeta_n. \leqno(1.19)
$$
By Theorem 1.1 in  \cite{LW},
 there exist constants $C,C_0$ such that}
\begin{equation}\label{th.LW1}
\E\,   | {\mathcal Z}_n|^2
= \log n + C
 + o(1), \  n\to\infty,\hbox{\it and for $m>n+1$, \ }\big|\E\,{\mathcal Z}_n\overline{{\mathcal Z}_m}\big|
 \le C_0 \max \Big( {1\over n},{1\over 2^{ m-n  }}   \Big).
 \end{equation}
  The explicit value of $C$ is
$$
  C=C_E -2  + 2\int_0^1 \phi(\alpha) d\alpha +
    2 \int_1^\infty \left( \phi(\alpha) -   {1\over 2\alpha} \right)
    d\alpha,
$$
where $C_E$ is the Euler constant and
$\phi(\alpha)={\alpha e^{\alpha} -2 e^{\alpha}+\alpha+2 \over  2\alpha^2
(e^{\alpha}-1)}$.
\vskip 3 pt By Theorem 1.2 in  \cite{LW}, the following results are proved, displaying a   slow growth of the
zeta function on the critical line, when sampled by the Cauchy random walk. For any real $b>2$,
\begin{equation}\label{th.LW2}
 \lim_{n\to \infty}{\sum_{k=1}^n \zeta({1\over 2} + iS_k) - n  \over
n^{{1/2}}(\log n)^b
}\buildrel{(a.s.)}\over=0 , \qq \bigg\|\sup_{n\ge 1}{\big|\sum_{k=1}^n \zeta({1\over 2} + iS_k) - n
\big|\over n^{{1/2}}(\log
n)^b }\bigg\|_2<\infty .
\end{equation}
    

  
 
